% Declared-model duty 2(d). Author source; not manuscript text.
% CURRENT BINDING SCOPE -- full GOAL_alt.md read on 2026-10-15.
% The binding reference is now the configuration-continued occupied
% spectral projector commuting with the back-averaged Hamiltonian.
% Spherical sources are explicitly angularly projected, with the
% compatible -J row. The recorded transported theorem remains a valid
% comparator, not a theorem for that spectral reference without an
% adiabatic/reference-comparison estimate. UA1--UA18's exact unprojected
% radiality obstructions, including the pending preparation compact,
% are controls and do not exclude the declared spherical truncation.
% The earlier pending compact promotion is withdrawn. Earlier uses of
% "declared transported reference" record their stated former premise.
% The exact constant-field ground-projector/no-crossing results and
% gapped static locality retain their stated mathematical domains.
% STATUS [R]: continued-spectral locality at the explicit massive-gap
% and massless subcritical finite-parent domains below. Existing [X]
% homogeneous ground/cut checks keep their exact counterpart scopes.
% Literal-transport and unprojected-source results remain controls.
\noindent\textbf{[R] Continued-spectral local area.}
Let $\mu=ma$, $\epsilon=a/L$, $E=|h_0|$ and
$h=\{n,h_0\}/2+\sum_i\{w_i,Q_i\}/2$, with the actual adjoint fields
$n=W_0^\dagger N$, $w=W_1^\dagger v$.
Constant-field continuation preserves $P_0$ at every tilt (node fillings
fixed); the ground projector equals $P_0$ if $|w\cdot q|<nE$ at every
nonzero-energy momentum. For $j=1,2,3$ equal nonzero flow components,
the threshold is $|w|/n=\sqrt{1+\mu^2/j}+\mu/\sqrt j$.

At fixed $\mu>0$ on the infinite parent, take a norm-continuous path from
$h_0$ to $h$ with uniform global gap $(-\delta,\delta)$, $\delta>0$,
bounded $n,w$, $n\ge n_*>0$, and
$|\partial^\alpha(n,w)|\le C_\alpha L^{-|\alpha|}$ for $|\alpha|\le2$.
For a regular cut, $A_{\rm lat}=A/a^2$,
\begingroup
\setlength{\abovedisplayskip}{4pt}
\setlength{\belowdisplayskip}{4pt}
\[
 \frac{|S_A(P)-S_A(P_0)|}{A_{\rm lat}}
 \le C\epsilon^2(1+|\log\epsilon|)^4 .
\]
\endgroup
\vspace{-1pt}
The state changes at first gradient; time reversal cancels the linear
entropy term. The bound does not fix the second-order coefficient.
For $0\le\mu\le\mu_*$, even antiperiodic parents and cuts comparable to $L$,
bounded scaled Fourier first moments and $\inf n-\sup|w|\ge c>0$ give
normalized error $O(\epsilon^{1/10})$ if continued $P$ equals the ground
projector and the flat parent/cut limit holds, without a uniform gap.
Both inherit the proved flat planar/spherical coefficients in their domains.
Inhomogeneous overtilted occupations need separate control.
Constants depend on $\mu$, $g=G/a^2$ and the branch.
Uniform time/allowed-$\ell$ bounds and the physical excitation cut
remain actual-history checks.

% DERIVATION / CHECK DETAIL.
% Uniform theorem: q_i=d_i cos(ap_i), d_i=sin(ap_i)/a, so |q|<=|d|<=E.
% The two double-degenerate eigenspaces of h_0 have projectors P_+,P_-.
% The actual occupied covariance is
% P=1_{v.q+NE<0} P_+ + 1_{v.q-NE<0} P_-.
% It is P_- iff neither band has crossed, up to measure-zero zeros with
% explicitly fixed filling. A positive constant N rescales energies and
% a scalar shift that causes no crossing changes no eigenvector. This is
% an exact lattice result, not a continuum covariance premise.
%
% Axial threshold: for fixed p_x, transverse sine terms only increase E.
% Put mu=|ma|, t=sin^2(ap_x). Minimizing E^2/q_x^2 gives
% min_{0<t<1} (mu^2+t)/[t(1-t)]. The stationarity equation is
% t^2+2mu^2 t-mu^2=0, hence t*=mu(sqrt(1+mu^2)-mu), and the minimum square
% root is sqrt(1+mu^2)+mu. At mu=0 the infimum 1 occurs at the nodes.
% For a direction with d equal nonzero Cartesian components, d=1,2,3,
% the exact threshold is sqrt(1+mu^2/d)+mu/sqrt(d). At fixed
% S=sum_{i=1}^d t_i, concavity of sqrt(t(1-t)) maximizes
% |n.q|=sum_i sqrt(t_i(1-t_i))/sqrt(d) at equal t_i=S/d;
% inactive components only increase E. The minimization then reduces to
% (mu^2/d+t)/[t(1-t)], proving the formula. Thus [110] and [111]
% massive flows cross before the axial value; the massless threshold is
% 1 in every direction by the small-momentum limit. A radial cubic-lattice
% field cannot use the axial massive threshold at all angular positions.
% The unrestricted condition remains |v.q|<N E at every momentum.
% For a strict crossing there are open filled/emptied pockets, separated
% from the one-band sea by codimension-one zero-energy surfaces. Across
% each regular sheet a covariance eigenvalue jumps by one. The following
% elementary covariance-variance bound proves the needed obstruction for
% literal planar slabs without assuming a matrix-valued Widom theorem.
% The example is homogeneous and irrotational, so small gradients or the
% restriction to longitudinal flow cannot repair the counterexample.
% It is a finite-energy Hermitian fermion system; no imaginary energy or
% growth rate is alleged. Normal ordering changes no covariance or entropy.
%
% [X] Cross-lane check of the obstruction by an elementary second route:
% the actual counterpart strand_Q_flow_reference.tex QF5--QF8, read only
% after this author's independent derivation was saved, proves a slab
% number-variance bound. This author and its RED rechecked it directly.
% Let A have L planes and one-face transverse area A_perp, with infinite
% longitudinal parent first and transverse thermodynamic limit before
% L->infinity. A transverse channel crosses iff
% M_perp=sqrt(mu^2+sin^2 k_y+sin^2 k_z)<(w-w^-1)/2. Call the transverse
% fraction f_perp>0. Each such channel has two roots of
% w^2 t(1-t)=M_perp^2+t in each of four monotone quadrants, hence eight
% jumps of rank two. Sum_j tr(Delta P_j)^2=16. Fourier integration by
% parts gives C_l=(2pi i l)^-1 sum_j exp(-i l k_j)Delta P_j+O(l^-2).
% Var(N_A)=2 sum_{l>=1}min(l,L)tr(C_l^dag C_l), so its log L coefficient
% is 8/pi^2 per crossed channel. Distinct-root cross terms are bounded
% oscillatory harmonic sums. For fixed mu>0, uniform bounded variation
% supplies domination in the transverse integral, even near merging roots.
% Therefore Var/[(A_perp/a^2)log L]->8 f_perp/pi^2 and, from
% h(c)>=4log(2)c(1-c), liminf S/[(A_perp/a^2)log L]
% >=32log(2) f_perp/pi^2. This rigorous positive lower bound already
% disproves a uniform area coefficient; it is not the exact log coefficient.
% Counted by total area 2A_perp these constants are halved. No scalar
% shell model, zero-mode filling assumption at finite grid, or exchange
% of the massless and thermodynamic limits is used in this proof.
% For the massless counterexample, do not take an unproved mu->0 limit:
% restrict the transverse integration to any positive-measure subset
% delta<=M_perp<M_c with delta>0. Each of those fibers has a uniform
% positive transverse gap, so the same bounded-variation proof on that
% subset gives a strictly positive slab entropy lower bound at mu=0.
%
% [X] Homogeneous evolved-state/source check: for every N(t),v(t),
% [h(t),P_-]=0, so the initial flat sea stays P_- under exact unitary
% evolution even through a reference crossing. Let P_cross be the set
% |v.q|>N E and integrate the full BZ with d^3p/(2pi)^3. The intrinsic
% source and current of this retained sea RELATIVE TO THE INSTANTANEOUS
% GROUND STATE, not to the declared transported reference, are
% rho_rel=-2 int_{P_cross} E,
% J_rel=2 int_{P_cross} q sign(v.q),
% while e_rel=N rho_rel+v.J_rel=2 int_{P_cross}(|v.q|-N E)>=0.
% Every factor two is the native band spin multiplicity. These agree with
% counterpart QF9--QF10 by direct band occupation counting. Energies are
% real; this comparison has no exponentially growing mode. No constraint-
% satisfying inhomogeneous object or global stability follows from it.
%
% Actual ground-state locality: use a fixed finite-range hopping operator
% with a global spectral gap Delta>0 about zero and smooth N(x),v(x).
% Its Fermi projector is a contour integral of the resolvent. Resolvent
% locality gives exponential decay with a length xi controlled by the
% hopping/gap ratio; resolvent differentiation gives the gradient expansion
% when xi/L<<1 and the derivatives are bounded. Uniform local spectral
% no-crossing alone is not a proof of a global gap for a nonuniform h.
% A chosen finite set of sites has a Gaussian reduced state determined by
% the restricted projector. The leading local symbol is P_0(p)=P_-(p),
% independent of N(x),v(x) while no crossing occurs. The correlation-state
% error has first-gradient terms; entropy requires a separate argument.
% A constructive sufficient global-gap test on the periodic Cartesian
% parent is available, rather than assuming that every pointwise local
% gap survives. Write h=N h_0+sum_i v_i Q_i+R,
% R=[h_0,N]/2+sum_i[Q_i,v_i]/2, and
% B=(||[h_0,N]||+sum_i||[Q_i,v_i]||)/2 >= ||R||.
% Put N_min=inf_x N(x)>v_max=sup_x |v(x)| and
% D^2=sum_i[sin(ap_i)/a]^2. Since sum_i Q_i^2<=D^2 exactly,
% ||h psi|| >= N_min sqrt(m^2+||D psi||^2)-v_max||D psi||-B
%            >= |m|sqrt(N_min^2-v_max^2)-B for ||psi||=1.
% A strictly positive final expression therefore certifies the actual
% global gap for a finite-amplitude nonuniform class. The commutators
% are bounded by nearest-neighbor N differences divided by a and
% two-step v_i differences divided by 2a, respectively. They are O(grad).
% This test is sufficient, not necessary, and is silent at m=0. It does
% not assert a gap for every pointwise subcritical coefficient profile.
% The RED proposed the norm bound; the author independently checked the
% decomposition, pointwise Cauchy--Schwarz estimate and scalar minimum.
%
% Explicit first-gradient check in midpoint/Weyl ordering:
% [N h_0,P_1]=-(i/2){partial_i(N h_0+v_j q_j),partial_{p_i}P_0}
%             =-i (partial_i v_j) q_j partial_{p_i}P_0.
% Here {h_0,partial_p P_0}=0 exactly, so the lapse term vanishes. Since
% partial_p P_0 is off-diagonal between the +/- eigenspaces, one solution is
% P_1=-i h_0 q_j(partial_i v_j)partial_{p_i}P_0/(2N E^2).
% It is Hermitian and satisfies P_0 P_1+P_1 P_0=P_1. It is generically
% nonzero. The site-symmetrized lattice symbol differs from its midpoint
% leading symbol only at even derivatives of the smooth coefficients.
%
% Why the entropy's first local gradient vanishes, in the stated symmetric
% setting: P_0 is independent of uniform v. P_1 is linear in partial v and
% independent of the local v. Time reversal relates S[N,v] to S[N,-v]
% for a fixed spatial cut; the first entropy variation around the same P_0
% is therefore even in a perturbation that changes P_1's sign, hence zero.
% A normal lapse derivative is odd under reflection/complement exchange
% of a centered pure half-space, and its projector term already vanished.
% Tangential one-derivative terms integrate to a boundary divergence.
% This removes a first-gradient entropy density; it does NOT say every
% local reduced density matrix agrees to second order. Finite asymmetric
% regions, endpoint corrections and actual excited/mixed states need their
% own error analysis. For a smooth curved pure cut, complement symmetry
% also removes a term linear in its signed mean curvature.
%
% Size of corrections: a schematic local entropy density is
% a^-2 eta_n(mu)+c_1(mu,delta) (grad ln N)^2
% +c_2(mu,delta)(grad v/N)^2+c_3(mu,delta) second derivatives+... .
% The coefficients are finite under a uniform spectral gap/margin delta;
% their mu->0 or delta->0 limits need not be bounded. Gapless infrared
% modes can yield log(L/a) factors and nonlocal terms. A relative
% O((a/L)^2) form therefore requires fixed dimensionless gap and margin,
% or an explicit UV/IR separation whose IR contribution is bounded.
% In particular a horizon with mu=0 has lambda_-~ -a^2 p_parallel^3/2
% along its aligned branch, even though its eigenvectors remain flat.
% Its small denominators are not the ordinary gapped adiabatic theorem.
%
% Allowed ways forward, all within the declared model:
% (1) In a globally gapped no-crossing star exterior/interior, apply the
% theorem and quantify xi/L. (2) At or beyond crossing, compute the actual
% same-configuration ground projector and retain any Fermi-surface term.
% (3) For collapse, compute the unitarily evolved state's covariance from
% its declared initial state; compare its ultraviolet local area density
% and its physical cross-surface correlations with the reference separately.
% The declared subtraction in (3) uses the transported incoming sea.
% It does not replace the physical state or its reference by the
% instantaneous ground projector. Option (2) is a spectral comparator,
% not the binding normal-ordering prescription.
% No alternative is silently substituted for the required interior claim.
%
% Transported-reference theorem and counterexample [R].
% At the Hartree level, for the common one-body unitary U(t,t_i),
% P_T(t)=U P_in U^dag and C(t)=U C_in U^dag. In particular for spatially
% uniform N(t),v(t), [h(t),P_-]=0 and P_T=P_- exactly, including beyond
% the pocket threshold. With P_T as reference its mean sources vanish;
% the negative rho_rel computed above belongs only to the explicitly
% different instantaneous-ground comparison. Neither prescription changes
% the entropy of a specified covariance.
%
% The pocket condition does NOT imply P_T(t)=P_-(h(t)) for a general
% smooth history. A finite-parent counterexample needs no supercritical
% flow. Let h(t)=h_0+epsilon f(t)V, V={F,h_0}/2, where real bounded F
% varies slowly in space and smooth f has compact temporal support.
% Choose epsilon so N=1+epsilon f F>0, v=0 and the massive parent remains
% gapped throughout. For an occupied mode a and empty mode b,
% V_ba=(epsilon_b+epsilon_a)F_ba/2 can be nonzero: choose unequal
% positive/negative energy magnitudes and nonparallel spinors at momenta
% connected by a small nonzero Fourier component of F. Smooth long-wave
% F still has this property. At the end t_f, h(t_f)=h_0 but
% (P_T-P_0)_ba=-i epsilon V_ba exp[-i Delta_ba t_f]
%                  int f(s)exp[i Delta_ba s] ds+O(epsilon^2),
% Delta_ba=epsilon_b-epsilon_a>0.
% A smooth pulse with a nonzero indicated Fourier transform makes this
% nonzero. The exact finite-dimensional propagator is analytic in epsilon,
% so a nonzero first coefficient proves a genuine nearby counterexample.
% It has positive lapse and zero flow at all times.
% For an explicit vertex choose F=cos(q x), an occupied p=0 beta-negative
% spinor and the positive spinor at q along one axis. In its 2x2 Dirac
% block E_q=sqrt(m^2+d_q^2), d_q=sin(aq)/a, and
% V_ba=(E_q-m)d_q/[4sqrt(2E_q(E_q+m))] is nonzero. The other spin
% component supplies a second copy, not a required extra species.
% Making the pulse adiabatic suppresses the transition; it does not make it exactly zero
% solely because no band crosses. This refutes the unrestricted equality
% in the binding paragraph, not the transported-reference declaration.
%
% For a genuinely gapped slow history, an adiabatic theorem may bound the
% difference using the gap, time derivatives, duration and initial state;
% it is additional to the spatial resolvent expansion above. Massless
% crossing regions require an explicit UV/IR state estimate or the actual
% covariance. Even a finite-time local unitary that preserves an area
% bound need not preserve its coefficient. The ground-state locality
% theorem retains its distinct scope; TL1--TL8 supplies the separate
% actual-transport estimates with their explicit history/state premises.
%
% Actual engine reference comparison: finite transport versus continuation.
% Read the complete section 13.1, "Decision (a): the transported reference",
% of .wt-solve_a/solve/specs/SPEC_final4b_engine.md. The inspected
% file's modification time is 2026-10-15 09:28:21 UTC; its header names
% the goal of 2026-10-15 06:51 UTC. This records the actual inspected
% version without substituting a source title or later implementation.
% It defines R0=P_-(h), R1=U P_in U^dag and R2 as a fixed initial count
% of occupied instantaneous eigenvectors chosen by maximal overlap with
% the preceding selected subspace. Its ambiguity rule halves the interval
% at most three times. It chooses R2 for the engine, with R1 beside as
% a control, and asserts coincidence of R1,R2,R0 whenever no level crosses.
% The pulse already proved above is a counterexample to that assertion:
% h(t_f)=h0, the mass gap remains open, R2(t_f)=R0=P0 along its smooth
% continued negative subspace, but R1(t_f) differs from P0 by the displayed
% nonzero interband amplitude. No threshold or finite continuation-step
% refinement replaces the missing temporal adiabatic limit.
%
% More generally, a differentiable R2 constructed from instantaneous
% eigenvectors commutes with h(t). If it also obeyed the declared transport
% i dot R2=[h,R2], then dot R2=0. A changing spectral-continuation projector
% therefore is not exactly the common-unitary R1 reference. At exact
% degeneracies, a selected subspace inside the eigenspace still commutes;
% discontinuous selection updates require their explicit jumps. Uniform
% commuting histories can give exact coincidence. A quantified slow,
% gapped adiabatic limit can give approximate coincidence. Neither is an
% arbitrary no-crossing history. For any physical C,
% C-R2=(C-R1)+(R1-R2),
% so the alternative source contains the additional reference difference.
% This diagnosis changes no other lane's calculation and says nothing
% against entanglement cuts actually taken on its transported physical C.
% Such a cut is reference-independent for a fixed C; the source law is not.
% The above result remains [R], with the same independently checked pulse
% theorem. The inspected specification is a claim being tested, not an
% actual cross-lane agreement on the erroneous finite-time identification.
%
% Actual smallball3 finite ground-state comparator.
% Read the matched remote list_v3_v1.txt, out/smallball3_L48.log,
% out/smallball3_v1.json and the complete mode_smallball3,
% small_ball_block, symbol_gap, flat comparator and stencil paths of
% code/obs_box3d_v3.py under compute-host:<internal>/.
% The log closes at 2026-10-15 10:10:08 UTC and supplies all four
% positions r/a=.8660,4.5552,8.5294,12.5200. This calculation restricts
% the erf approximation to Theta(-h) of a prescribed inhomogeneous
% operator; it does not transport an incoming sea or solve its sources.
% Both local comparators use mu=0 and are the same uniform flat sea.
% On radii 3,4,5,6 the actual full4 ball blocks have 123,257,515,925
% sites. Differences from the local comparator range from -.0590100
% to .0390650 nats, and normalized Hilbert--Schmidt covariance distances
% range from .00182787 to .00474377. These are finite differences, not
% a measured first- or second-gradient law. For the central radius-4
% block S=169.686884329, the w/2 and w/4 changes are -9.99e-11 and
% -2.56e-10; maximum entry changes are below 3e-14. The diagonal
% projector-defect test is about 6e-14. A diagonal test on selected
% columns is not the norm of the full-parent projector defect.
% The inspected specification section 8.5 itself names the actual
% global spectral gap as the needed width input. Its sampled uniform
% symbol minimum and a different L=16 dense spectrum do not prove
% the stated L=48 lower bound; the following argument supplies a
% conservative genuine bound for that exact profile instead.
%
% A constructive global-gap certificate for the actual smallball3 control.
% Independently applying the above commutator theorem to its inspected
% profile gives a genuine gap bound, without promoting the sampled symbol
% or another box's spectrum to a theorem. Set a=1, m=1/5, r_s=4,
% R_m=8, R_1=14, R_2=22, N=1 and
% v_r=-(3/(5sqrt(2)))[(7/4)(r/8)-(3/4)(r/8)^3] for r<=8,
% v_r=-(6/5)r^(-1/2) for r>=8, multiplied by the stated quintic taper.
% Before taper the maximum speed is 7sqrt(14)/60<.437. For a radial
% vector field, |partial_i v_i|<=max(|v_r'|,|v_r/r|). In the cubic core
% this is at most L_v=21/(160sqrt(2)); between 8 and 14 it is smaller.
% In the taper region |chi'|<=15/64, hence
% |(chi v_r)'|<=(3/5)14^(-3/2)+(15/64)(6/5)14^(-1/2)<L_v,
% and |chi v_r/r|<L_v. The same bound holds for the periodic continuation
% because the field vanishes before every face. Thus every two-step
% coefficient difference satisfies delta_i=sup_x|v_i(x+2e_i)-v_i(x)|
% <=2 L_v, and B<=sum_i delta_i/4<=3L_v/2. On the actual antiperiodic
% Cartesian parent, sum_i Q_i^2<=D^2 still holds exactly. Therefore
% dist(0,spec h)>= (1/5)sqrt(1-343/1800)-63/(320sqrt(2))>.0407.
% This bound is independent of L whenever the specified taper fits in
% the parent; here it applies to the executed L=48 profile. It is a
% bound for this prescribed instantaneous operator, not an achieved star
% or a transported-reference identification. In restored units the bound
% is a^(-1) times the displayed number, at ma=.2 and the listed r/a.
% No g is selected by a prescribed profile; its possible Gauss sources
% and their g dependence have not been solved.
% The conservative bound is below the sampled symbol gap .181496, so
% it does not by itself certify the producer's quoted erf error at
% w=.0382097. The actual w,w/2,w/4 local block convergence is separate
% finite numerical evidence; it is not an operator-norm error theorem.
% For the exact matrix function f_w(h)=[1-erf(h/w)]/2, however,
% ||f_w(h)-P_-(h)||<=erfc(gamma/w)/2 with the proved gamma above.
% At the reported w/4 this exact-filter error is below 1e-9. Before
% polynomial/evaluation roundoff, a ball with n_A modes consequently
% obeys |S(f_w(h)_A)-S(P_-(h)_A)|<=n_A h_2(erfc(gamma/w)/2).
% This uses eigenvalue perturbation plus the binary-entropy continuity
% bound, not an inference from small differences between two filters.
%
% An independent exact coefficient/tail route separates the remaining
% polynomial question from the filter error. Write sigma=w/r with
% ||h||<=r, z=1/(2sigma^2), and expand f_w(rx)=sum_k a_k T_k(x).
% a_0=1/2, a_(2n)=0 for n>=1, and
% a_(2n+1)=(-1)^(n+1) exp(-z)[I_n(z)+I_(n+1)(z)]
%                         /[(2n+1)sqrt(pi)sigma].
% Derive this by integrating the cosine coefficient by parts and using
% sin(theta)sin((2n+1)theta)=[cos(2n theta)-cos((2n+2)theta)]/2;
% exp[-cos^2(theta)/sigma^2]=exp(-z)exp[-z cos(2theta)].
% For a degree K truncation, n0 is the first n with 2n+1>K. Since
% p_n=exp(-z)I_n(z) is the symmetric Skellam probability and its moment
% generating function is exp[z(cosh t-1)], Chernoff gives the rigorous
% exact-coefficient uniform tail bound
% sum_(k>K)|a_k| <= 2 exp[-n0 asinh(n0/z)+sqrt(z^2+n0^2)-z]
%                         /[(2n0+1)sqrt(pi)sigma].
% This is directly evaluable with the recorded K=2010 at w/4; it
% supplies a possible short scalar validation without a new 3D run.
% The producer obtains coefficients by finite cosine quadrature and
% applies a floating-point matrix recurrence. Their quadrature and
% arithmetic errors are additional; no unperformed bound or run for
% them is asserted. All new analytic filter/gap statements are [R].
%
% SOURCES.
% Binding GOAL_alt.md 2(d), model paragraph; THINK sections 8,9,12
% read in full for the assigned scope. The latter's unqualified second-
% gradient reduced-state statement is corrected by the explicit P_1 above.
% Ledger predecessor lane_z/B12-006 and original
% full_numerics/open/p5_static/lane_z_reviewa/g35_cut_entropy_on_planar_flow_unit_tilt.py
% read in full. Its two-step flow has the relevant uniform single-field
% symbol; its graded observable does not establish the present theorem.
% Only the distinction between no-crossing and crossed bands carries over.
% Independent mechanics-author message agrees on BW's no-crossing domain
% and distinguishes normal energy from the tilted time generator; this is
% same-lane coordination, not the required external agreement.
% New counterpart <internal>/solve_b
% strand_Q_flow_reference.tex and its scope read fully after independent
% substantive derivations were saved. QF1--4 agree by direct diagonalization
% and constrained minimization; QF5--8 by the separate jump/number-variance
% calculation above; QF9--10 by exact covariance evolution and band counting.
% The same-lane RED independently read the actual Q source and its scope
% and checked all three comparisons. This agreement concerns prescribed
% homogeneous controls, not a solved constraint-satisfying object.
% strand_E_declared_equations.tex and scope also read fully: its explicit
% state/reference distinction agrees, while its constraint sign and full
% finite-spacing stress closure belong to the mechanics task.
% Full current author source/SOURCES and independent RED source were read;
% all mandatory repairs, including the massless subset proof, incorporated.
% The hydrodynamic GOAL was reread fully: these results concern matter
% projectors on smooth prescribed or Hartree fields, not a microscopic
% geometric law or a proof of ell independence of a solved interior.
% The current full GOAL with the transported incoming-sea prescription
% was reread. The pulse/Duhamel argument was derived directly for that
% prescription; spectral results and Q comparisons above remain explicitly
% ground-state comparisons. No instantaneous reference is substituted.
%
% ACTUAL CPU VACUUM-FLOW CONTROL [N, NAMED SUBTRACTION ONLY].
% SOURCES read in place at compute-host:<internal>/:
% cputest_v5b.sh, out_cpu/b_vacflow.log,
% out_cpu/vacflow_box_L16_m0.25_b.json, complete
% code/obs_vacflow_box_v1.py and its executed obs_box3d_v5.py dependencies
% (native CPU stencil, prescribed field, common propagator and flat sea).
% This is a site-centred antiperiodic L/a=16 parent, ma=.25, nominal
% r_s/a=2, R_m/a=1, N=1, ramp duration 4a and dt=.5a. Its quintic
% taper starts at R1/a=-2 and ends at R2/a=6; it changes the entire
% physical profile. The evaluated observable is exactly
% Tr[e_x(P_T-P_0)], with e_x={h_0,Pi_x}/2, not Tr[e_x(P_T-P_T)].
% The backward propagation of Pi_x and h_0 Pi_x implements that first
% expression directly. Raw densities therefore need not vanish and
% do not equal the binding normal-ordered source of the transported sea.
%
% The recorded proposed gradient-fit window is [3a,-2a], so every
% n_fit is zero and every ray/coalesced coefficient is null. No local
% gradient coefficient is measured by this box result. Its separately
% specified zone integral evaluates to .01900638440019845 on the 160^3
% twisted grid and .01963106258306132 on the L=16 grid. Agreement of
% the first decimal with the quoted target checks that integral's
% implementation; it does not establish a response coefficient from
% these empty fits or a history-independent local reduced state.
% All quantities are in lattice units; g=G/a^2 is unspecified by this
% imposed field, and neither a self-consistent g-dependent source nor
% independence of the averaging scale ell follows.
% FINAL CPU COMPARISON: complete current author and RED sources/SOURCES
% were read. The exact flat-subtracted observable, zero declared source
% for C=P_T, empty fit and limited zone-integral interpretation agree.
% No new response coefficient or cross-lane [X] is inferred.

% ACTUAL TRANSPORT: A GAPLESS SURFACE ESTIMATE [R, independent derivation].
% These statements use P_T=U P_0 U^dag throughout. They neither refill
% instantaneous bands nor change the declared source/field allocation.
% Work first on the infinite native parent, set a=1 and mu=ma, and fix
% mu>=0 (including mu=0 itself). Put tau=t/a after restoring units.
% On 0<=s<=tau assume real scalar N_s(x)>0 and v_s(x), with a uniform
% bounded finite-range one-body generator
% h_s={N_s,h0}/2+sum_i {v_(i,s),Q_i}/2, ||h_s||<=H_*.
% Assume each coefficient is spatially Lipschitz with constant at most
% epsilon in lattice coordinates, uniformly in s. Time continuity or
% measurability sufficient for the propagator is enough; no temporal
% adiabatic theorem or instantaneous spectral gap is used. In particular
% |v|/N may pass through or exceed every pocket threshold. H_* and the
% coefficient bounds are fixed in the limit. An incoming covariance
% different from P_0 needs its own initial error term, not an automatic
% application of this theorem.
%
% TL1. Point-node bounds and the exact commutator. The same dyadic
% Fourier proof used in declared_area_law.tex SA1 gives
% ||(P_0)_(xy)||<=C_mu(1+|x-y|)^(-3),
% ||(A P_0)_(xy)||<=C_mu(1+|x-y|)^(-4), A=h0,Q_1,Q_2,Q_3,
% away from the harmless finite-range part of the second expression.
% At mu=0, h0 P0=(h0-E)/2 and Q_i P0=Q_i(1-h0/E)/2 have conical
% degree-one rather than degree-zero singularities at the eight nodes;
% four integrations by parts on each dyadic annulus give the stated
% bound. At fixed mu>0 the symbols are smooth. This does not use a
% false gap at the massless nodes or an instantaneous occupied-set cutoff.
% For a scalar multiplier f and any of these A, [A,P0]=0 gives exactly
% [ {f,A}/2,P0 ]_(xy)
%   =(f_x-f_y)(A P0)_(xy)+{[A,f],P0}_(xy)/2.
% The first term has one factor epsilon |x-y|; [A,f] has fixed finite
% range and entries O(epsilon). Consequently
% S_s=[h_s,P0] obeys ||(S_s)_(xy)||<=C_mu epsilon(1+|x-y|)^(-3),
% with row and column l2 norms <=C_mu epsilon. No Schur l1 summability
% of the massless kernel is presumed.
%
% TL2. A deliberately conservative finite-history bound. Finite-range
% Dyson propagation and unitarity give, for q=0,3,
% sup_x sum_y ||U(s,u)_(xy)||(1+|x-y|)^q
%       <=C_q(1+H_*|s-u|)^(q+3/2),
% and the analogous column bound. To prove it, use the usual exponential
% off-cone Dyson bound outside a ball of radius C(1+H_*|s-u|), and
% Cauchy--Schwarz plus the unit row/column l2 norm inside that ball.
% Polynomial moments of the exponentially small exterior obey the same
% estimate; no translation invariance of the actual generator is needed.
% For w(z)=(1+|z|)^3 use
% w(x-y)<=C[w(x-z)+w(z-z')+w(z'-y)] in U S U^dag.
% The middle term is bounded by C epsilon times the two q=0 moments.
% In either exterior term, first contract S with one propagator row by
% Cauchy--Schwarz, using its l2 bound, then use the other q=3 moment.
% Thus ||(U S U^dag)_(xy)||
% <=C_mu epsilon(1+H_*|s-u|)^(9/2)(1+|x-y|)^(-3).
% Duhamel now proves
% ||(P_T(tau)-P0)_(xy)||<=C_mu kappa(tau)(1+|x-y|)^(-3),
% kappa(tau)=epsilon tau(1+H_*tau)^(9/2).                 (TL2)
% The power 9/2 is a sufficient propagation estimate, not a measured
% physical growth law or an optimal admissible collapse time.
%
% TL3. Turn this into a surface, not a volume, entropy estimate. Let A
% be a digital ball or a regular family of finite cuts with crossing
% count #{x in A:x+z notin A}<=C A_lat |z|, where A_lat is its boundary
% area in lattice units. The same estimate holds for planar surface
% densities by finite transverse regularization followed by its limit.
% Write T_P=1_A P 1_(A^c). For a pure covariance P,
% S_A(P)=Tr F(|T_P|), F(s)=h_2((1-sqrt(1-4s^2))/2), 0<=s<=1/2.
% The kernel decay and crossing count give
% ||T_P||_(3/2)^(3/2)<=C_mu A_lat(1+kappa)^(3/2),
% ||T_(P_T)-T_(P0)||_2<=C_mu sqrt(A_lat) kappa.
% Here the block-entry Schatten bound uses p=3/2>4/3; the relevant
% radial sum is sum_r r^(3-3p), which converges. Since
% |F'(s)|<=C s^(3/4), singular-value perturbation and Cauchy--Schwarz
% yield the genuine entropy estimate
% |S_A(P_T)-S_A(P0)|<=C_mu A_lat kappa(1+kappa)^(3/4).    (TL3)
% This proof controls the actual covariance and its full cut entropy,
% including the massless point nodes; a volume continuity bound is not
% being mistaken for a local-area theorem.
%
% In particular, if epsilon->0 and kappa(tau)->0, the transported sea
% has the same leading native surface coefficient eta_n(mu) as P0 on
% any plane for which the surface density exists, and the same proved
% sphere average on growing digital balls. For a local patch one may
% choose 1<<R/a<<L/a with epsilon=O(a/L), using the actual local normal
% and physical area. At fixed tau the normalized error is O(a/L).
% A sufficient growing-time regime is tau=o((L/a)^(2/11)) for fixed
% H_*. This is not the macroscopic collapse regime tau=O(L/a), and
% the estimate proves no failure in that longer regime. The binding
% hydrodynamic hierarchy a<<ell<<L supplies spatial scale separation;
% it does not supply this time bound, an initial-vacuum error bound, or
% a stronger long-time band-transport theorem. Its own self-consistent
% source history must satisfy those premises before TL3 can be used.
% There is no g hidden in eta: at g=G/a^2, the result is conditional
% on the actual N_g,v_g, through H_*(g), epsilon(g), tau(g), and mu.
% No ell-dependent source is changed or discarded by this estimate;
% ell independence of those histories remains a distinct requirement.

% TL4. A nonzero first-gradient entropy rate for actual transport.
% This exact planar calculation also prevents transferring the centered
% instantaneous-ground second-gradient statement to arbitrary P_T.
% Let N=1, v_y=v_z=0 and let a bounded smooth v_x(n) equal
% k(n-1/2) on the four sites n=-1,0,1,2 around the [100] cut n>=1.
% It may taper on a much larger length with |v|<1 everywhere. Start
% from P0 and use the declared one-body transport. At a fixed transverse
% momentum set B=mu beta+alpha_y sin k_y+alpha_z sin k_z, B^2=M^2,
% h0(k_x)=alpha_x sin k_x+B and E_M(k)=sqrt(M^2+sin^2 k).
% The exact native half-chain modular theorem already used in
% declared_bw_flow.tex gives K_A=b(M){X,h0}_A/2,
% X_n=n-1/2, b(M)=4 K_par(-M^2), where K_par uses its parameter.
% This is the full four-component fiber, without extra valley factors.
%
% Internal terms of h_v do not change the cut entropy at the initial
% state because K_A commutes with (P0)_A. Only the two Q-crossing
% edges contribute: (h_v)_(1,-1)=-ik/8 and (h_v)_(2,0)=ik/8.
% Therefore dS_A/dtau|0=2 Im Tr_A[K_A(h_v)_AB(P0)_BA]
% =k r(M), with
% r(M)=-(b(M)/2) int_(-pi)^pi E_M(k)cos(2k) dk/(2pi)
%     =(b(M)/2) int sin^2(k)cos^2(k)/E_M(k) dk/(2pi)>0. (TL4)
% The second expression follows by one integration by parts. Explicitly,
% the five nonzero modular entries in this trace combine to
% tr[B P_(-2)]-(i/2)tr[alpha_x P_(-1)]
%                  +(i/2)tr[alpha_x P_(-3)]
% =-2 int E_M(k)cos(2k) dk/(2pi), fixing the sign and full-spin count.
% r(0)=2/3 and r(M)~log(4M)/(4M^2) at large M. For an arbitrary v
% the exact k in this rate is
% k_d=[v_2+v_0-v_1-v_(-1)]/2
%    =a partial_x v_x+(13/24)a^3 partial_x^3 v_x+... at the cut.
% A uniform v contributes zero, as also follows from [Q,P0]=0.
%
% Integrating transverse momenta gives a strictly positive coefficient
% R(mu)=int_BZperp r(sqrt(mu^2+sin^2 k_y+sin^2 k_z)) d^2k/(2pi)^2,
% d(S/A_phys)/dt|0=(k_d/a^3) R(mu).
% For mu=0 the zero transverse channels have measure zero, and r(M)
% has the finite limit above. The half-chain first-variation formula
% holds for M>0; its transverse integration extends to mu=0 using
% the uniform boundary entanglement-rate bound for finite-range
% four-component chains. The latter is independent of the gap and
% chain length; internal Hamiltonians are ancillas and the crossing
% terms act on a fixed finite local Fock space. It justifies dominated
% convergence of the entropy difference quotients, rather than using
% a nonexistent finite entropy of the isolated M=0 half-chain.
% Thus a smooth bounded subcritical prescribed flow can have a genuine
% O(partial v) transported-entropy response at its initial time.
% No decimal for R(mu), thermal interpretation or solved spherical
% collapse is claimed. This field family is an admissible native matter
% Hamiltonian; it is not asserted to solve the declared Gauss sources.
% Its role is a precise state-dependent counterexample to an inference
% from smooth coefficients alone, and a coefficient to test on any
% actual reached field that has the corresponding local normal jet.
%
% SOURCES FOR TL1--TL4. The full current binding goal was reread.
% TL1--TL3 were derived independently before new RED transport intake,
% using the previously proved native Fourier and Schatten surface
% estimates in declared_area_law.tex SA1--SA3. The exact h_f commutator,
% propagation moment exponents and Duhamel bound are supplied here.
% TL4 was independently derived from the actual exact half-chain
% modular generator; declared_bw_flow.tex's source/domain proof and
% its primary Eisler theorem locator were read in full for that step.
% Van Acoleyen, Marien and Verstraete, arXiv:1304.5931v3, primary full
% paper pp.1--3, equations (1)--(3) and the proof through (15), was read:
% its size-independent boundary interaction rate bound supplies only
% the dominated-convergence bound, not the coefficient in TL4.
% No new counterpart result, numerical check or cross-lane X is used.

% MACROSCOPIC-TIME CONTINUATION [R, independent weighted derivation].
% TL5--TL8 were saved before reading the RED's new transport source.
% They strengthen the time range of TL1--TL3 under stronger, explicit
% space-and-time regularity. The massless proof does not insert a
% positive physical mass or replace the transported covariance.
%
% TL5. Precise regularity class. Let epsilon=a/L in (0,1), X=epsilon x
% in native coordinates, s=epsilon tau, and 0<=s<=T with T fixed.
% Write N=N(X,s), v_i=v_i(X,s), N>=N_*>0. Assume uniform bounds on
% their amplitudes and on their Fourier-measure norms
% int (1+|Q|)^8 d|fhat_s|(Q),
% for f=N,v_i and a_ij=(partial_(X_i) v_j)/N, and for partial_s a_ij.
% Constants are allowed as point masses at Q=0. Smooth periodic profiles
% or sufficiently smooth decaying profiles are examples. The displayed
% eight moments are sufficient rather than sharp. This is a specific
% quantitative hydrodynamic history class, not a consequence of the
% word "smooth" alone. In particular it controls the time derivative
% of the spatial flow gradient on the scale L. The exact Hamiltonian
% still uses endpoint symmetrization, and the initial state is P0.
% Fix 0<=mu<=mu_* and define the auxiliary positive momentum weight
% rho_e(k)=sqrt(mu^2+sum_i sin^2(k_i)+epsilon^2).
% This weight is a norm device, not a fermion mass or an altered sea.
%
% TL6. The weighted first correction closes the massless estimate.
% Put E^2=mu^2+sum sin^2 k and, away from its zero nodes,
% S_ij(k)=q_j(k) partial_(k_i) P0(k),
% B_ij(k)=-i h0(k) S_ij(k)/(2 E(k)^2).
% S_ij is bounded; B_ij has order E^(-1), and
% [h0,B_ij]=-i S_ij. Let chi(E/epsilon) vanish for E<=epsilon
% and equal one for E>=2epsilon, and define B_ij,e=chi B_ij.
% For j=0,1,2 its derivatives obey
% ||partial_k^j B_ij,e||<=C_j rho_e^(-1-j).
% Use the LEFT-ORDERED auxiliary correction
% K_e(tau)=epsilon sum_ij a_ij(epsilon xhat,epsilon tau) B_ij,e,
% where xhat is the native position operator. K_e need not be
% self-adjoint or a projector. Only the exact transported P_T is used
% in entropy; Q_e=P0+K_e is an approximate solution of its linear
% covariance equation and nothing is refilled.
%
% All following norms are operator norms on the actual infinite lattice:
% ||K_e rho_e||<=C epsilon,
% ||(i partial_tau Q_e-[h,Q_e])rho_e||<=C epsilon^2.       (TL6)
% Here is a direct Fourier-transfer proof, including the node domain.
% A scalar coefficient transfers q=epsilon Q between momenta k and k+q.
% For A=h0 or Q_j the exact commutator vertex is
% (fhat(Q)/2)[(A(k+q)+A(k))P0(k)
%                    -P0(k+q)(A(k+q)+A(k))].
% Its first term in q is -q_i fhat(Q){A,partial_i P0}/2.
% For A=h0 it vanishes; for Q_j it is -q_i fhat(Q) S_ij.
% The remainder times rho_e(k) is bounded by
% C (|q|^2+epsilon |q|)|fhat(Q)|. To check this without
% differentiating at a node, split E(k)>2C|q|, where Taylor's remainder
% is C|q|^2/E(k), from E(k)<=2C|q|, where the exact vertex and its
% bounded first term are O(|q|). Multiplication by E(k)+epsilon gives
% the bound in both regions. Thus
% [h,P0]=i epsilon sum_ij (partial_i v_j) S_ij+R,
% ||R rho_e||<=C epsilon^2.
%
% In [h,K_e], the undifferentiated coefficient term is
% epsilon N a_ij [h0,B_ij,e]
% =-i epsilon (partial_i v_j) chi S_ij and cancels this source.
% The uncancelled (1-chi) term has E<=2epsilon and hence weighted
% norm O(epsilon^2). Every coefficient commutator contributes one
% transfer q. For example [N,B_e]h0 rho_e is O(|q|) times a
% polynomial in 1+|q|/epsilon, since B_e has order rho_e^(-1).
% The external epsilon and Fourier moments make it O(epsilon^2).
% Likewise [h0,a]B_e rho_e and [Q_i,a]B_e rho_e have size O(epsilon).
% The endpoint-ordering terms {N,h0}/2-Nh0 and {v,Q}/2-vQ
% have norm and rho_e-conjugated norm O(epsilon); their commutators
% with K_e therefore also obey the weighted O(epsilon^2) bound.
% Finally partial_tau K_e=epsilon^2 sum(partial_s a_ij) B_ij,e.
% These estimates prove (TL6) by integrating the stated finite Fourier
% moments. No uniform operator-norm convergence of P_T to P0 at the
% massless nodes is assumed or required.
%
% TL7. Propagation of the weight and a local mean-square estimate.
% For the same exact vertex,
% ||[rho_e,h]rho_e^(-1)||<=C epsilon:
% |rho_e(k+q)-rho_e(k)|<=C|q| and
% ||A(k+q)+A(k)||/rho_e(k)<=C(1+|q|/epsilon).
% The Fourier moments control the resulting epsilon(|Q|+|Q|^2).
% Gronwall, also for backward propagation, consequently gives
% ||rho_e U(t,u)rho_e^(-1)||,
% ||rho_e^(-1) U(t,u)rho_e||<=exp(C epsilon |t-u|).
% Duhamel applied to Q_e, with P_T(0)=P0, now yields for tau<=T/epsilon
% ||(P_T(tau)-P0)rho_e||<=C_T epsilon.                  (TL7)
% D=P_T-P0 is Hermitian. For a site/spin basis vector e_(x,sigma),
% ||D e_(x,sigma)||<=C_T epsilon ||rho_e^(-1)e_(x,sigma)||.
% In three dimensions the latter norm squared is
% int_BZ dk/(2pi)^3 [mu^2+sum sin^2 k+epsilon^2]^(-1)<=C,
% uniformly for 0<=mu<=mu_* and epsilon>0. The eight point nodes
% are integrable in this estimate. It follows that
% sup_x tr_spin(D^2)_(xx)<=C_T epsilon^2.               (TL7a)
% This is the needed infrared estimate. A uniform spectral gap would
% give an unnecessarily stronger norm statement and would exclude
% precisely the massless case of interest.
%
% TL8. A genuine local-area limit through macroscopic times. For a
% finite region with n_A native sites and regular boundary area A_lat,
% let T_D=1_A D 1_(A^c). Equation (TL7a) implies
% ||T_D||_2^2<=C_T epsilon^2 n_A,
% ||T_D||_(3/2)^(3/2)<=C_T epsilon^(3/2) n_A,
% the latter by its rank <=4n_A and Schatten interpolation.
% Combine this with the flat surface bound and the same F' estimate
% as TL3 to obtain
% |S_A(P_T)-S_A(P0)|
% <=C_T epsilon sqrt[n_A(A_lat+epsilon^(3/2)n_A)].       (TL8)
% For balls or regular scaled shapes of native radius R this gives
% |Delta S|/A_lat<=C_T epsilon sqrt R
%                         sqrt(1+epsilon^(3/2)R).
% The flat kernel/Schatten constant here is uniform for 0<=mu<=mu_*:
% on each valley annulus of radius b,
% |partial_k^j P0|<=C_j b^(-j) uniformly in that mass interval.
% The same dyadic Fourier sum is uniformly O((1+|z|)^(-3)).
% This is a uniform transport-error bound; the separate sphere
% asymptotic is still taken at each fixed mu. A uniform small-mu
% angular expansion is not inferred from it.
% Thus for 1<<R<=C/epsilon the normalized error is at most O(sqrt epsilon),
% uniformly through tau<=T/epsilon, including mu=0 and bounded
% overtilted flows. On a local patch take 1<<R<<1/epsilon; on a
% macroscopic digital sphere R=O(1/epsilon), combine with the proved
% actual flat-sphere asymptotic. The leading coefficient is the same
% eta_n(mu)/a^2 with the actual physical normal, or its proved sphere
% orientation integral. No shell degeneracies, planar decomposition
% of a curved cut, or independently assumed local vacuum are used.
%
% The O(sqrt epsilon) bound is a controlled sufficient error, not a
% sharp gradient expansion. TL4 still exhibits a genuine linear
% initial-gradient correction; a general second-gradient remainder
% for transported entropy is therefore not established. TL8 concerns
% finite regions in the indicated joint hydrodynamic limit, not an
% infinite planar entropy density at fixed nonzero epsilon, arbitrary
% rough growing boundaries, or an arbitrary joint finite-parent limit.
% The native parent is infinite; parent/initial-state errors need
% separate estimates before importing finite numerical controls.
% For rational planar faces, their flat finite-cut limit follows
% without treating a plane as a sphere: use a lattice parallelepiped
% in an integer unimodular basis. Each fixed finite covariance
% polynomial, after inclusion--exclusion for its outside vertices,
% has intersection counts equal to products of integer side lengths
% minus the path's coordinate ranges.
% Their expansion gives the face terms with O(R) edge error.
% The finite-range polynomial approximation and surface Schatten
% bound of the area note then remove the polynomial cutoff. Each
% face term is exactly its primitive-plane chain coefficient with
% the physical area factor. This supplies the local [100], [110]
% and [111] face interpretation in the joint limit, without a new
% tilted numerical coefficient or a shell-degeneracy reduction.
%
% Binding scope. TL5 is compatible with a smooth positive-clock
% history and the declared common-unitary reference, including a
% trapping region. The Gauss law is not used to manufacture such a
% history: any actual unprojected solved history in this regularity
% class inherits TL8. The goal's spatial hierarchy alone does not
% prove its uniform time-dependent Fourier bounds, positive N_* or
% initial covariance P0, and does not exclude much longer histories.
% At g=G/a^2 and fixed ma, these constants and epsilon=a/L are those
% of N_g,v_g at the chosen ell. The result supplies no cancellation
% of ell-dependence in different source histories; that still has to
% be checked by the model's actual source evolution. It introduces
% no projected radial closure and no new geometric law.
% If TL5's constants are uniform as ell varies in the declared range,
% the leading coefficient is consequently uniform and independent of
% ell even when the finite history/error differs. This conclusion
% does not make all geometric observables or finite entropy corrections
% independent of ell. N_* bounds the declared flat-slice clock N,
% not the Schwarzschild observer's surface redshift alpha; N=1 flow
% histories can cross a trapping surface while meeting this premise.
% SOURCES: TL5--TL8 are an independent Fourier-transfer/Duhamel
% continuation, saved before reading the RED's macroscopic-time
% argument or its unfinished massless cutoff route. The previously
% read native sphere proof supplies only the flat Schatten bound.
%
% ACTUAL ORIGINAL-PAIR TRANSPORT COMPARISON.
% The complete current author and RED sources/SOURCES were read, including
% RED TR1--TR8, the exact modular/source checks, and its amended periodic
% extension premise. Author TL1--TL4 was saved before RED transport
% intake; TL5--TL8 was also saved before reading RED TR1--TR5 or its
% macroscopic cutoff continuation. RED TR6--TR7 was saved before its
% reading of TL5--TL8. These are the named independently saved routes:
% the author's direct weighted infinite-parent Fourier/Duhamel proof,
% and RED's endpoint-rank cutoff/Schatten proof with buffered AP-parent
% transfer. They agree on the leading coefficient at matching domains;
% their O(sqrt epsilon) and O(epsilon^(1/10)) sufficient error bounds
% are different estimates, not conflicting physical gradient laws.
% Author independently checked RED's exact lapse endpoint identity,
% both low-mode endpoint supports, delta>=sqrt(epsilon) Fourier tails,
% Schatten exponents, dyadic AP image estimate and causal-buffer
% transfer. The extension must retain Fourier seminorms also for
% (partial_i v_j)/N and its scaled-time derivative, as now stated.
% RED independently checked the exact left-ordered correction,
% endpoint-allocation residual, both weight conjugations, integrable
% rho_e^-2 site norm and final surface entropy inequality in TL5--TL8.
% The separate five-entry half-chain contractions agree on TL4's
% positive initial-gradient coefficient and exact full4 normalization.
% The new statements remain [R]. The actual counterpart Q comparison
% supports only the unchanged homogeneous/slab [X] statements. No
% cross-lane general-history theorem, numerical acceptance or solved
% self-consistent collapse is manufactured from this same-pair agreement.

% ACTUAL UNPROJECTED HISTORY APPLICABILITY [R; independent continuation].
% UA1--UA4 below were derived and saved before reading the RED's new
% applicability route or the new Q transported-entropy-budget source.
% They leave the published transport theorem intact. They test whether
% its premises follow from the binding native sources and spherical
% continuum Gauss law, rather than inserting prescribed fields into it.
%
% UA1. An exact source-assignment obstruction. Specify an ordinary
% admissible averaging rule, rather than leaving "averaging" undefined:
% first give each native cell its constant density e_x/a^3, then convolve
% with the normalized radial Gaussian of coordinate variance ell^2.
% Let e_x=Tr(e_x^{op} D), e_x^{op}={Pi_x,h0}/2, for trace-class D.
% The cell sequence is absolutely summable, because diagonal pinching
% and Holder give sum_x |e_x| <= ||h0 D||_1 <= E_max ||D||_1.
% Thus the cell interpolation f_a is in L1(R^3). If K_ell*f_a were
% exactly rotation invariant, then, for every rotation O,
%   exp(-ell^2|k|^2/2)[fhat_a(k)-fhat_a(O k)]=0.
% The Gaussian multiplier is nowhere zero; Fourier uniqueness makes
% f_a itself radial. A grid-cell-constant radial L1 function is zero.
% To prove this last assertion distributionally, apply each generator
% x_i partial_j-x_j partial_i. On the interior of a cell face with a
% jump, its coefficient is the jump times a tangential coordinate.
% This is nonzero on a positive-area subset of that face unless the
% jump vanishes. Different open faces cannot cancel as distributions.
% All adjacent cell constants therefore agree; L1 then forces zero.
% Consequently ANY nonzero trace-class native energy source under this
% explicit averaging rule fails exact spherical symmetry. This extends
% the compact-preparation argument in the faithful note to reached
% trace-class D(t)=U Gamma U^dag, with no analyticity of its long tails
% assumed. For point-cell allocation the corresponding conclusion is
% that a radial finite signed measure supported on aZ^3 is supported
% only at the origin; a nonzero off-origin energy cell is enough to
% obstruct radiality. The box-allocation conclusion is stronger.
%
% For spherical N(r), v(r) rhat, the scalar continuum Einstein flux is
% radial identically. It cannot equal the nonradial source just proved.
% This obstruction is present at an initial slice and is independent of
% the disputed momentum-row sign, of the reference's later transport,
% and of how positive one tries to make N. It therefore prevents an
% exact unprojected spherical history for these data before TL5's time
% or Fourier estimates become relevant. It is not a no-go for every
% possible source allocation: a noninjective low-pass filter requires
% its own actual source proof. An angular projection would change this
% conclusion by changing the binding source and is not inserted here.
%
% UA1a. Actual independent RED strengthening, checked after both saves.
% The point-readout exception of an energy measure at the central lattice
% site is also excluded for ANY nontrivial physical initial covariance
% 0<=C<=I with D=C-P0 trace class. Put S_x=(-1)^(x1+x2+x3) and R=S beta.
% Native hopping reverses S while beta anticommutes with that hopping,
% hence {S,h0}=2mR, [R,h0]=0 and R=R^dag=R^-1. Consequently
%   sum_x S_x e_x=m Tr(RD), sum_x e_x=E_D=Tr(h0 D).
% Define D_+=P_+ C P_+>=0, H_-=P_-(I-C)P_->=0 and
% N_ex=Tr(D_++H_-). Since R is band diagonal,
% |Tr(RD)|<=N_ex, whereas
%   E_D=Tr[E(D_++H_-)]>|m| N_ex
% for every C!=P0. The multiplication operator E-|m| has no l2 kernel:
% its zeros are only the isolated sine nodes. A nonzero positive trace-
% class excitation therefore gives the strict inequality. If N_ex=0,
% positivity of C and I-C forces the off-diagonal band block to vanish,
% so C=P0. These arguments retain arbitrary particle-hole coherences.
% Exact radial point averaging with an injective radial kernel would
% force e_x=E_D delta_(x,0), so E_D=m Tr(RD), a contradiction. A center
% not on a lattice site leaves no invariant atom and forces E_D=0.
% The cell-box version already forces every e_x=0 and is contradicted
% by the same strict energy inequality. Thus neither specified map has
% a nontrivial finite-excitation flat-incoming spherical initial datum.
% This strict energy inequality is NOT asserted relative to an arbitrary
% later transported P_T; UA1's cell-source injectivity remains valid
% at later times when that source is nonzero. The independent proof is
% RED AC1--AC2, actually read and verified by the author. It adds no
% momentum-sign assumption and no universal all-filter exclusion.
%
% UA2. Nonempty physical domain and the hydrodynamic distinction.
% Take the actual mixed Gaussian preparation
%   Gamma=lambda P_+ F P_+, C_in=P_0+Gamma, 0<=F<=1,
% with real nonzero finite-support F and 0<lambda<=1. The faithful
% note's native calculation gives, with dimensionless Brillouin average,
%   e_x=lambda[(e0+b0)F_x+sum_i c_i(F_{x+e_i}+F_{x-e_i})],
% e0=<E>, b0=m^2<1/E>, c_i=<sin^2 k_i/E>/(2a^2)>0,
% and e0=b0+2 sum c_i. Hence e_x>=0, J_x=0, and
% sum_x e_x=2lambda e0 sum_x F_x>0. No spectral refill or negative
% particle covariance is used. UA1 applies for every m>=0 and ell>0.
% For F_x=f(a x/L) with a smooth radial compact taper, the already
% proved source residual is O(a/L) relative to its core density,
% uniformly under normalized radial convolution. Thus UA1 does not
% say that a hydrodynamic approximate spherical description is bad;
% it says that exact unprojected equality and its existence interval
% have not been supplied by that approximation. Retaining a quantified
% angular residual is different from solving all the declared rows.
% In lattice units epsilon=a/L, mu=ma, g=G/a^2 and ell/a are independent
% controls. The weak-compactness choice g=gamma epsilon^2/(2lambda <aE>)
% keeps g rho_tilde (L/a)^2=O(gamma), while a<<ell<<L is attainable.
% UA1 persists at every finite epsilon; it is not a large-curvature or
% too-small-ell counterexample. No elliptic equation with angular
% components has been silently added to the spherical-only model.
%
% UA3. What source caps actually give, and what they do not give.
% If a total source measure has variation M_abs, Gaussian convolution
% only implies the coarse Fourier estimate
%   int(1+L|k|)^j |rhohat_ell(k)| d^3k
%      <= C_j M_abs ell^-3 [1+(L/ell)^j].
% Its constants grow in the allowed ell/L -> 0 limit. Under an existing
% bounded common history, ||dot D||_1 <= 2||h|| ||D||_1 supplies a
% time bound, but rescaling t=L s introduces L||h||, generally of
% order epsilon^-1 in lattice units. Neither estimate proves failure
% of a better state-specific bound; both show explicitly why the native
% occupancy caps, normalized smoothing, and a<<ell<<L alone do not
% furnish TL5's uniform eighth Fourier moments and scaled-time bound.
% A constructive route must control the actual source's long-wave
% Fourier distribution and current evolution, a positive-mass branch
% of the field map, and its N lower bound uniformly over the chosen
% ell range. The faithful note's local projected Picard theorem is
% not such an unprojected macroscopic-time construction.
%
% UA4. A leading hydrodynamic angular obstruction for a permitted
% particle preparation, beyond the exact finite-a cell discrepancy.
% This is an exactly computed FREE-HISTORY comparison and a necessary
% weak-field-limit test, not a claimed nonzero-G solved history.
% Set a=1, mu=m>0, F_epsilon(x)=f(epsilon x), with smooth compact radial
% 0<=f<=1, and the same Gamma=lambda P_+ F_epsilon P_+.
% In the free history h=h0, tau=s/epsilon, the macroscopic energy
% density has the weak Fourier limit
%   rho_0(X,s)=2lambda int_BZ E(k) f(X-s w(k)) d^3k/(2pi)^3,
%   E=sqrt(mu^2+sum sin^2 k_i), w_i=partial_i E=sin k_i cos k_i/E.
% Proof: at transferred momentum q=epsilon Q, the exact integrand is
%   lambda (E_++E_-)/2 tr[P_+(k+)P_+(k-)]
%        exp[-i(E_+-E_-)s/epsilon] Fhat_epsilon(epsilon Q),
% k+=k+q/2, k-=k-q/2. After the usual epsilon^3 Riemann normalization,
% the weight tends to 2E, the phase to exp[-is Q.w], and the smooth
% sampled envelope to fhat(Q). Bounded convergence applies for fixed
% Q,s,mu>0. The same normalized Gaussian, including the cell box, leaves
% this limit unchanged when ell/L->0: its multiplier at epsilon Q tends
% to one. The same holds for a fixed normalized approximate-identity
% family, not for arbitrary ell-dependent filters without that property.
% Cubic symmetry makes the second velocity moment isotropic, but the
% fourth moment has A=int E w_1^4 and B=int E w_1^2 w_2^2. Its source
% term is
%   (lambda s^4/12)[3B Delta^2+(A-3B)sum_i partial_i^4]f.
% For generic radial f the last term has a nonzero l=4 angular part.
% It is not proportional to a power of epsilon. A fully analytic
% witness is mu>=4:
%   A-3B=-3/(128 mu^3)+R_mu,
%   |R_mu|<=81/(256 mu^5)<3/(128 mu^3).
% Here <(sin k cos k)^4>=3/128 and
% <(sin k_1 cos k_1)^2(sin k_2 cos k_2)^2>=1/64;
% |E^-3-mu^-3| <= (3/2)mu^-5 sum sin^2 k_i proves the remainder bound.
% For equal small |Q| in the [100] and [111] directions the Fourier
% densities differ at order
%   lambda fhat(Q) s^4 |Q|^4 (A-3B)/18,
% which is nonzero for sufficiently small nonzero s|Q| and fhat(Q)!=0.
% This remains true as epsilon->0 for that approximate-identity averaging
% with ell/L->0. It shows why a spatially smooth radial envelope filled
% over the native positive band is not by itself a radial fluid state.
%
% This comparison also gives a precise conditional obstruction to a
% proposed regular weak-field spherical branch for these preparations.
% If on 0<=tau<=s/epsilon its native fields obey
%   int ||h_g(tau)-h0|| d tau ->0,
% Duhamel gives ||U_g Gamma U_g^dag-U_0 Gamma U_0^dag|| ->0, since
% ||Gamma||<=lambda. Each cell-energy difference is then bounded by
% ||e_x^{op}||_1 times this norm, uniformly in x. Its macroscopic
% source limit must equal the nonradial expression above. Such a
% branch cannot simultaneously have an exactly radial limiting source.
% The displayed small integrated-field condition is explicit; a mere
% assertion that g is small does not prove it. It is satisfied, for
% example, by uniform fields ||N-1||+||v||=o(epsilon) on this interval,
% at fixed mu. This is a necessary weak-field branch test, not a proof
% that the binding equations produce those fields.
%
% An allowed way forward is a different actual particle preparation
% with controlled valley momenta and rotational phase-space moments,
% followed through its coupled evolution. For instance the legitimate
% covariance lambda P_+ chi_kappa F chi_kappa P_+, 0<=chi_kappa<=1,
% restricts valley momenta; it does not alter the Hamiltonian/reference.
% To use it one must prove its native averaged source errors, positive
% mass/clock interval, angular residual, and Fourier/time bounds through
% the desired time. Small lattice momentum and continuum rotational
% moment matching are additional state controls, not consequences of
% the positive preparation or of the spatial averaging hierarchy.
% No such full coupled proof or automatic ell cancellation is asserted.
% SOURCES: native preparation and exact finite-support radial obstruction
% are in declared_faithful_vacuum.tex; UA1's all-trace-class Gaussian
% extension, UA3 bounds and UA4 ballistic Fourier/moment test are derived
% here independently. The current binding GOAL_alt.md was fully
% reread. No new RED, Q or B continuation was read for these arguments.
%
% ACTUAL APPLICABILITY COMPARISON / SOURCES.
% UA1--UA4 and RED AC1--AC5 were substantively saved before either new
% route was read. The actual complete new blocks and SOURCES were then
% compared. Author checked AC1's checkerboard identity, strict CAR band
% energy bound and zero-excitation coherence argument; UA1a records that
% accepted strengthening with its flat-incoming domain. AC2's interval
% proof of the cell-box obstruction agrees with UA1's face-distribution
% proof. The radial clock integral, smoothing estimates, finite-time
% source-map domain and all-filter limitations in AC3--AC5 agree.
% RED AC6 independently verified UA4's full4 Fourier vertex, normalized
% macroscopic limit, fourth-moment remainder and directional coefficient,
% as well as the integrated weak-field hypothesis. No g>0 free-history
% substitution or low-momentum inference is made. The original proof
% routes and their different scopes remain explicit.
% After both independent saves, the actual new
% <internal>/solve_b
% strand_Q_transported_entropy_budget.tex and its scope were read fully.
% Its finite-matrix spectral-shift/rank-trace entropy argument is sound
% at QE2--QE9's stated scope, but QE1 uses the HA ADJOINT fields and
% its QE12 comparison assumes an existing common history. It proves
% neither the unprojected point-coupled existence interval nor TL5's
% uniform field/history bounds. Its actual scope is [S], independent
% review pending; it is not used to promote the present result to [X].
% FINAL PAIRED READ: both complete current sources and all SOURCES were
% read, including RED AC6 and its exact final compact/domain additions.
% The UA4 Fourier vertex retains lambda. RED's corrected zero-source
% statement is also accepted: v=0 makes K=0, so the Gauss constraints
% alone do not fix N; N=1 is a choice, not a consequence of subtraction.
% This does not change the initial radial-source obstruction. The proved
% result is for the specified maps and domains, with no all-filter or
% solved-history conclusion. No mandatory scientific repair remains.

% NONINJECTIVE INITIAL-SOURCE CONSTRUCTION [R; independent continuation].
% UA5--UA6 were derived and saved before any new RED construction was
% read. The printed theorem/obstruction is unchanged. The point readout
% and averaging kernel below are fixed explicitly; no angular projection
% or change to the native Hamiltonian, current or transported reference
% is made. An initial constraint solution is not a coupled interval.
%
% UA5a. A positive radial averaging kernel with an actual Fourier cutoff.
% Work in lattice units a=1, mu=ma>0. Choose once a nonnegative smooth
% w supported inside (1,2), with integral one, and define
%   b(r)=int_1^2 w(t)[sin(tr)/(tr)]^2 dt, phi(x)=b(|x|)^6.
% The value at zero is by continuity. Then phi is smooth, radial,
% strictly positive, and at most one. The Fourier representation of
% sin(t|x|)/(t|x|) as the uniform probability measure on the radius-t
% sphere proves supp phihat subset {|k|<=24}. This use of spheres is
% only a Fourier identity for a smoothing kernel, never an entropy cut.
% Integration by parts in the smooth t integral gives
% b(r)=b_infty/r^2+O_M(r^-M), b_infty=(1/2)int w(t)/t^2 dt>0,
% including differentiated remainders. Thus phi~b_infty^6 r^-12,
% phi is integrable and has finite moments through order eight.
% Its logarithmic gradient is bounded: |grad log phi|<=C_phi.
% Put K=phi/int phi and K_ell(y)=ell^-3 K(y/ell). This is a fixed,
% positive, normalized radial approximate identity; Khat_ell has support
% |k|<=24/ell and is therefore noninjective. It is an explicit choice
% of averaging, not a deduction that all allowed kernels have this form.
%
% UA5b. Deconvolve the actual native cell-energy stencil, retaining CAR.
% Let E=sqrt(mu^2+sum sin^2 k_i), e0=<E>, b0=mu^2<1/E>,
% c=<sin^2 k_1/E>/2, alpha=e0+b0, and T=sum_i(S_i+S_i^dag).
% Cubic symmetry gives alpha+6c=2e0, alpha-6c=2b0>0.
% The actual source of Gamma=P_+ F P_+ is
%   e=(alpha I+cT)F=:A F, J=0,
% by the native identity in UA2, extended to nonnegative l1 F by
% trace-class convergence. Choose a target positive cell sequence
% e_L(x)=A_L phi(x/L), L>0, and beta=6c/alpha<1.
% Define F_L=A^-1 e_L by its convergent l1 Neumann series. This
% inverse is not assumed to preserve positivity for arbitrary inputs.
% Positivity for the displayed target follows from an explicit bound.
% In the weighted sup norm ||f||_(e_L)=sup_x |f_x|/e_L(x),
% ||T||<=6 exp(C_phi/L). If beta exp(C_phi/L)<1, then
% ||A^-1||_(e_L)<=1/[alpha(1-beta exp(C_phi/L))].
% Compare F_L to F_0=e_L/(2e0). Since
% A F_0-e_L=(c/(2e0))(T-6I)e_L,
%   |F_L-F_0| <= theta_L F_0,
%   theta_L=beta[exp(C_phi/L)-1]/[1-beta exp(C_phi/L)].
% At each fixed mu>0 choose L large enough that theta_L<1, and choose
% 0<A_L<=2e0/(1+theta_L). Then 0<F_L<=1 and F_L is l1. Hence
%   Gamma_L=P_+ F_L P_+, C_in=P0+Gamma_L
% is a genuine mixed Gaussian covariance with 0<=Gamma_L<=P_+,
% finite expected addition number N_f=2 sum F_L and exactly the
% desired native energy cells e_L. Time reversal gives J_x=0 exactly.
% No fractional expectation is reinterpreted as a pure fixed-N sector.
% This particular global inverse does not cover mu=0: there beta=1.
% The separate low-band construction in UA5e below closes that case;
% no massless limit of this global Neumann estimate is taken.
%
% UA5c. Exact unprojected radial source for one family of ell values.
% Use the POINT readout rho_ell(y)=sum_x e_L(x)K_ell(y-x), with the
% same fixed F_L for every ell; no cell box is inserted. If L>24/pi
% and ell>24/pi, Poisson summation gives
% ehat_L(k)=A_L L^3 sum_(n in Z^3) phihat(L(k+2pi n)).
% Khat(ell k) removes all noncentral replicas. Consequently, exactly,
%   rho_ell(y)=A_L (K_(ell/L)*phi)(y/L),
% which is positive and radial. The current remains exactly zero.
% This is a true native finite-excitation datum, not a continuum source
% substituted for it. It evades UA1a because the averaging has Fourier
% zeros and the preparation has algebraic infinite tails; the finite-
% support analyticity argument is inapplicable. The source map is fixed
% before checking the result. It is not the Gaussian or cell-box map.
% Total energy is exactly E_tot=A_L L^3 int phi, independent of ell,
% and N_f=E_tot/e0. The interior profiles still depend on ell/L; only
% their controlled limit and total charge are asserted independent.
%
% For every g=G/a^2>0 this supplies an exact initial spherical solution
% of BOTH printed momentum-sign choices:
% M_ell(r)=4pi int_0^r s^2 rho_ell(s)ds,
% N=1, v=-sqrt(2g M_ell/r) rhat.
% The center is regular, M_ell>=0, and the scalar flux equals rho_ell;
% a radial irrotational v at N=1 has zero momentum flux, as required.
% The actual matter coupling is still the point-evaluated declared h.
% To keep a fixed geometric scale and amplitude in the large-L family,
% set A_L=gamma/(g L^2) with fixed positive gamma. The CAR amplitude
% condition then holds for sufficiently large L, and
% v(LX)^2=2gamma m_delta(X)/X,
% m_delta(X)=4pi int_0^X s^2(K_delta*phi)(s)ds, delta=ell/L.
% N_f=gamma L(int phi)/(g e0). At fixed mu,g,gamma and delta<=delta0,
% the initial fields and their scaled spatial derivatives are uniform:
% convolution bounds control all source derivatives; the strictly
% positive central density makes the square root smooth at X=0;
% at infinity v is a fixed r^-1/2 radial-vector tail plus an integrable
% smooth remainder. Differentiating sufficiently often gives integrable
% derivatives uniformly in delta, and the homogeneous tail has Fourier
% degree -5/2, locally integrable in three dimensions. Splitting this
% low-momentum tail from the high-momentum derivative bound proves the
% required finite spatial Fourier moments of any fixed order, including
% those of (partial_X v)/N, at this INITIAL slice. These bounds do not
% supply its scaled-time derivative or a coupled existence interval.
% Taking a<<ell<<L and, for large gamma, ell<<L/sqrt(gamma), retains
% the stated initial geometric hierarchy. A_L, e0, L and ell here are
% lattice-unit quantities; restore L=L_phys/a, ell=ell_phys/a, mu=ma.
%
% UA5d. Test the actual native tangency, rather than declare it.
% This construction does not automatically define an invariant radial
% source domain. At its initial N=1, Gamma_L is time-reversal invariant.
% The source current q_(x,i) and flow Hamiltonian h_v are time-reversal
% odd; -i[q_(x,i),h_v] is odd because time reversal is antiunitary.
% Therefore its expectation in Gamma_L vanishes. The exact initial
% current derivative under the FULL h=h0+h_v is the free derivative,
% independent of the value of g or of the solved radial v at that slice.
% This statement uses Ddot=-i[h,D], the actual common-reference law.
%
% Its momentum formula is
% dot J_i(q)=-i Fhat_L(q) B_i(q),
% B_i(q)=int_BZ [q_i(k+q/2)+q_i(k-q/2)]/2
%             (E_+-E_-) tr[P_+(k+)P_+(k-)] d^3k/(2pi)^3.
% Here q_i(k)=sin(2k_i)/2; the transferred vector is the argument q.
% Since Fhat_L=ehat_L/A(q) with A(q)=alpha+2c sum_i cos q_i>0,
% the averaged derivative is multiplied by the scalar Khat(ell q).
% An exactly radial vector field must have Fourier vector parallel to q.
% The exact analytic B_i near q=0 has cubic-symmetry expansion
%   B_i=b_1 q_i+b_3 q_i^3+b_12 q_i sum_(j!=i)q_j^2+O(|q|^5).
% For sufficiently large FIXED mu its nonradial coefficient is nonzero:
%   b_3-b_12=-1/(6mu)+O(mu^-3).
% To derive this, use E_+-E_-=mu^-1 sum_j u_j sin q_j+O(mu^-3),
% tr[P_+^+P_+^-]=2+O(mu^-2), u_j=sin k_j cos k_j, and
% [q_i(k+)+q_i(k-)]/2=u_i cos q_i. Since <u_i u_j>=delta_ij/8,
% B_i=sin(2q_i)/(8mu)+O(mu^-3), with uniform derivatives through
% cubic order. The expansion follows analytically from the uniformly
% convergent large-mu expansion of E and P_+; no fitted threshold is
% claimed. Thus b_3-b_12 cannot vanish for all sufficiently large mu.
% More exactly, with w_i=partial_i E and
% g_ij=partial_i(h0/E).partial_j(h0/E) in Dirac-vector notation,
% b_3-b_12=int u_1[(E_111-3E_122)/12
%                   -w_1(g_11-g_22)/2+w_2 g_12-w_1].
% This integral defines the coefficient at every fixed mu>0; the
% preceding large-mu statement is the proved nonvanishing domain.
%
% Both phihat(Lq) and Khat(ell q) are nonzero in a neighborhood of zero.
% Choose q with two unequal nonzero components there. The nonzero cubic
% coefficient makes q cross dot J(q) nonzero for sufficiently small q.
% Thus for all sufficiently large fixed ma this EXACT initial state
% fails the tangent condition for a radial vector source under the full
% native point-coupled Hamiltonian, for every g>0. No differentiable
% exactly spherical history through this datum can satisfy all the
% unprojected Gauss rows, despite its exact positive radial initial
% charges. This is a result for this specified construction and mass
% domain, not an all-filter or all-state impossibility theorem.
% It also explains why a local Picard theorem for an independently
% projected source map cannot supply the missing native tangency.

% UA5e. Complete the low-band positivity and the tangency sign at mu=0.
% This is an additional derived step, not an unproved mass-limit exchange.
% It works uniformly for 0<=mu<=mu_* at fixed finite mu_*.
% On continuum functions whose Fourier support is |k|<=24/L, choose a
% real Schwartz reproducing kernel psi_L(x)=L^-3 psi(x/L), with
% psihat=1 on |k|<=24. It is only a proof device, not the averaging rule.
% Use the Banach norm ||f||_(12,L)=sup_x (1+|x|/L)^12 |f(x)|.
% For the unit coordinate shifts T, the exact identity f=psi_L*f gives
%   (6-T)f=[6psi_L-sum_i(psi_L(.+e_i)+psi_L(.-e_i))]*f.
% A central-difference Taylor estimate on this Schwartz kernel gives
% the weighted L1 norm of the bracket <=C_psi/L^2. Submultiplicativity
% of the polynomial weight therefore proves
% ||6-T||_(12,L)<=C_psi/L^2 on this low-band subspace.
% Since A=2e0[I-(c/(2e0))(6-T)] and c/(2e0)<=1/12 for ALL mu>=0,
% its low-band inverse is a convergent Neumann series once
% r_L=C_psi/(12L^2)<1. It preserves Fourier support, reality and l1
% summability of its sampled output. There are constants c_phi,C'_phi>0
% with c_phi(1+|X|)^-12<=phi(X)<=C'_phi(1+|X|)^-12. Thus for the
% same e_L=A_L phi(. /L), the constructed F_L obeys
%   |F_L-e_L/(2e0)| <= theta'_L e_L/(2e0),
%   theta'_L=(C'_phi/c_phi) r_L/(1-r_L).
% Choose L so theta'_L<1 and A_L<=2e0/(1+theta'_L). This proves
% 0<F_L<=1 and A F_L=e_L even at mu=0, without ever inverting the
% zero at the distant checkerboard momentum. Sampling gives the same
% exact native preparation and Poisson source of UA5c. For mu>0 this
% inverse agrees with UA5b on the target by uniqueness of that bounded
% inverse. At mu=0 it is a specified low-band solution, which is all
% this preparation requires. Initial CAR, mass, clock and uniform spatial
% bounds therefore hold for every fixed mu>=0, including zero itself.
%
% An exact endpoint simplification also proves that the tangency
% obstruction is present for EVERY finite mu>=0. Write u_i=sin k_i cos k_i,
% t_i=sin^2 k_i and define finite Brillouin integrals
%   A_mu=<u_1^2/E>, B_parallel=<cos(2k_1)t_1/E>,
%   B_perp=<cos(2k_1)t_2/E>.
% Shifting the E_+ and E_- denominator terms separately in UA5d gives
% the exact trigonometric polynomial
% B_i(q)=A_mu(1+cos(2q_i))sin q_i
%       -(sin(2q_i)/2)[B_parallel(1-cos q_i)
%                          +B_perp sum_(j!=i)(1-cos q_j)].
% In this identity use
% (E_+-E_-)tr(P_+^+P_+^-)
% =[(E_+^2-d_+.d_-)/E_+]+[(d_+.d_--E_-^2)/E_-],
% d=(mu,sin k_1,sin k_2,sin k_3); the spin trace is the native full4.
% The shifts reduce every denominator to E(k); odd cross integrals vanish.
% Consequently the exact nonradial cubic coefficient is
%   Delta_mu=b_3-b_12=-7A_mu/3-B_parallel/2+B_perp/2.
% This agrees with the independently saved RED coefficient reported
% after the two original UA5/AC constructions existed. The following
% integration-by-parts sign proof is an additional author derivation.
% Integrating partial_1(u_1 t_1/E) and partial_1(u_1 t_2/E) gives
%   B_parallel+2A_mu=<t_1 u_1^2/E^3>,
%   B_perp=<t_2 u_1^2/E^3>.
% Hence
%   Delta_mu=-4A_mu/3+(1/2)<(t_2-t_1)u_1^2/E^3>
%            <=-5A_mu/6<0,
% because t_2-t_1<=E^2 and A_mu>0. At mu=0 the integrations are valid
% after excising small balls around the eight nodes: u_1 t_j/E=O(r^2)
% and its surface boundary term vanishes as O(r^4). All remaining
% integrands are locally integrable. This proves nonvanishing directly
% at mass zero, with no derivative of the point-node projector assumed.
% Therefore the actual averaged initial current derivative of UA5c is
% nonradial at EVERY finite ma>=0, for every g>0 and every named ell.
% Positive radial initial charges and uniform initial spatial bounds
% do not give a differentiable exactly spherical unprojected history
% through this family. Neither the same conclusion for every possible
% preparation nor failure of a controlled approximate history follows.
%
% The initial angular error also has a controlled hydrodynamic size.
% On the entire retained Fourier support |q|<=24/L,
%   B_i(q)/A(q)=(A_mu/e0)q_i+O_mu(|q|^3),
% uniformly on any compact mass interval. The scalar denominator is
% bounded below there, including mu=0; its quadratic term is radial.
% Because phihat is continuous with compact support and |Khat|<=1,
% Fourier inversion and any fixed number of scaled X derivatives give
%   ||dot bar J+(A_mu/e0) grad_x rho_ell||_(C_X^k)
%       <=C_(k,mu_*) A_L/L^3,
% uniformly over UA5c's alias-free scales with ell/L<=delta0 and
% 0<=mu<=mu_* at sufficiently large L.
% The leading radial derivative has scale A_L/L. Thus the absolute
% angular residual is at most O(L^-2) times that scale, with no claim
% of a pointwise relative bound where the leading derivative vanishes.
% The exact nonzero curl remains: smallness is not radial tangency.
% No time-integrated angular bound or approximate coupled history is
% obtained from this initial estimate. Physical units give the relative
% factor (a/L_phys)^2, with A_L=gamma/[g(L_phys/a)^2] in this family.

% UA6. What the zero-flow lapse degeneracy actually says.
% For D=0, the displayed matter/mean-source equations give C=P_T and
% rho=J=0. With v=0 the displayed K_ij vanishes for every smooth
% positive N(t,x), so both Gauss rows vanish too. Smooth asymptotic
% normalization N->1 does not remove this algebraic degeneracy.
% One may even set N=1 initially and thereafter choose a smooth lapse
% pulse: common transport keeps D=0 exactly for every g and ell.
% However this is only a family of solutions of the DISPLAYED
% matter-plus-constraint relations, not a proved physical family of
% the full binding H2. The binding paragraph also requires geometry
% to be a fixed functional of its sources with no independent spherical
% degrees of freedom, and requires consistent slice deformation.
% Arbitrary choices of N at the same zero source do not meet a specified
% single-valued functional prescription. That prescription or its gauge-
% preserving evolution has not been supplied by the two Gauss rows.
% A translation-covariant vacuum functional and N(infinity)=1 would
% select the constant N=1. Continuity of the normalized source map from
% the positive radial J=0 branch would also select N[0]=1, since the
% nonzero branch has N=1 wherever its enclosed mass is positive. This
% is a branch-continuity premise, not division by u=0 in the vacuum row.
% Alternatively, the full vacuum Einstein
% evolution equation at flat spatial metric and K=0 would require
% partial_i partial_j N=0, again selecting N=1 with that boundary
% condition; that extra evolution equation is not imported as a proved
% consequence of the present declared constraints or matter algebra.
%
% The degeneracy therefore identifies a missing uniqueness/branch input,
% not a physical violation of the local-area theorem. Formally choosing
% N_delta=1-(1-delta)chi(x/L)f(t), 0<delta<1, with smooth 0<=chi,f<=1,
% compact spatial chi and f initially zero, illustrates that the Gauss
% rows themselves supply no uniform N_* across delta->0. L can be
% enlarged to retain a<<ell below every finite geometric length. But
% this alone neither makes those pulses allowed by the full H2 nor
% proves failure of entropy universality; v=0 in fact removes TL6's
% first flow correction. The required uniform clock bound must be
% obtained for the actual chosen geometry functional and reached state,
% not from this underdetermination or from setting N=1 by convention.
% SOURCES: the current binding model and the previously paired UA/AC
% obstruction were used. UA5's kernel, positive deconvolution, exact
% Poisson source and native current tangency test were derived here
% independently before new RED intake. UA6 separates the displayed
% constraints from the full consistency and functional requirements;
% no new geometric evolution law is asserted.
%
% ACTUAL NONINJECTIVE COMPARISON / SOURCES [R].
% The full saved RED AC7--AC12 and their SOURCES were read after UA5--UA6
% existed. The original constructions were independently saved before
% intake. RED uses a fixed nonnegative Schwartz averaging kernel with
% compact Fourier support and an averaged sinc^8 target with positive
% r^-8 tails; this source map differs explicitly from UA5a's r^-12 kernel.
% Both specify their point readout, low-frequency target, CAR covariance,
% lattice units and g before solving the initial Gauss rows. I checked
% the RED weighted inverse, Poisson alias exclusion and cubic coefficient.
% Its bound |Delta+1/(6mu)|<=19/(16mu^3), hence mu>=3 nonvanishing, follows
% by |1/E-1/mu|<=3/(2mu^3) and the full-zone elementary moments.
% The all-mass IBP sign was derived on both sides before the matching
% new proof was read. RED's independent even-cutoff chi/W inverse has
% all weighted lattice moments and reproduces the low-band target exactly;
% its second-difference positivity bound provides a second route to UA5e,
% including uniformity on compact mass intervals and mass zero. The
% current derivative and its vanished flow contribution agree for the
% actual common transport and both momentum-row signs. The initial
% spatial Fourier-tail bounds and the vacuum functional/continuity scope
% also agree. These are preparation-specific initial-source/tangency
% results, not an exact or approximate coupled existence theorem, a
% statement about every averaging rule, or a new other-lane [X] result.
% FINAL NONINJECTIVE READ: complete current author and RED sources and
% all SOURCES were read, recovering the truncated displayed range.
% This includes RED AC13, the all-mass low-band/sign completion, the
% same-preparation ell range, initial Fourier-tail bounds, and the
% quantified initial angular remainder. The compact statement retains
% the transported theorem and names only this construction's initial
% radial charges and native tangency failure. The vacuum lapse result
% concerns the displayed subsystem unless a compatible source functional
% is supplied; no physical counterexample follows from underdetermination.
% No mandatory scientific repair remains at these [R] scopes.

% MORE GENERAL NATIVE INITIAL COVARIANCE [R; independent continuation].
% UA7 is derived and saved before any new RED route is read. It retains
% precisely UA5a's fixed positive noninjective radial K and UA5c's point
% readout. It constructs radial initial rho,J and radial ACTUAL dot J,
% evading the particular one-envelope tangency obstruction above. It is
% not a coupled interval or a change to the Hamiltonian/reference law.
%
% UA7a. Four explicitly fixed preparation filters. Set a=1, mu=ma>=0,
% eta=1/100, kappa=pi/4, and
% b_eta(t)=exp[1-1/(1-(t/eta)^2)] for |t|<eta, zero otherwise,
% periodically extended to the native momentum torus. Define
% b_0(k)=b_eta(k), b_c(k)=b_eta(k-c)+b_eta(k+c),
% T_1(k)=b_kappa(k_1)b_0(k_2)b_0(k_3), with T_2,T_3 its axis permutations,
% T_4(k)=prod_i b_(pi/2)(k_i).
% These real even scalar multipliers are between zero and one, are
% supported away from every sine node, and have pairwise disjoint supports.
% They prepare the state only; they do not filter the Hamiltonian or its
% actual sources, and introduce no extra fermion component/species.
%
% For any real diagonal F_j, let Gamma_j=P_+ T_j F_j T_j P_+.
% At transferred momentum q, k_+=k+q/2, k_-=k-q/2, set
% z_j(k,q)=T_j(k_+)T_j(k_-) Tr[P_+(k_+)P_+(k_-)], and define exactly
% W_j(q)=int (E_++E_-)/2 z_j(k,q) dk/(2pi)^3,
% B_ij(q)=int [Q_i(k_+)+Q_i(k_-)](E_+-E_-) z_j(k,q) dk/[2(2pi)^3].
% The native energy and free current derivative of Gamma_j have Fourier
% transforms W_j(q) Fhat_j(q) and -i B_ij(q) Fhat_j(q), respectively.
% All spin traces are full4; all integrations are the actual native zone.
% Reflection in coordinate i makes B_ij odd in q_i and even in the other
% q coordinates. W_j is even in every q coordinate. Therefore the matrix
% M(q) with first row W_j(q) and remaining rows B_ij(q)/q_i is real,
% smooth and coordinatewise even near zero, with the quotient at q_i=0
% defined by its smooth derivative. The filters remove the massless nodes
% from these integrals; this claim is valid at mu=0, not a gapped limit.
%
% UA7b. A proved invertible source/pressure matrix, at every finite mass.
% At zero transferred momentum its columns are
% M_1(0)=(e_A,p_parallel,p_perp,p_perp)^T, its two permutations,
% and M_4(0)=(e_B,p_0,p_0,p_0)^T, where
% e_A=2<E T_1^2>, p_parallel=2<Q_1^2 T_1^2/E>,
% p_perp=2<Q_2^2 T_1^2/E>, e_B=2<E T_4^2>, p_0=2<Q_1^2 T_4^2/E>.
% Put n_A=<T_1^2>>0 and n_B=<T_4^2>>0. On supp T_1,
% .49<=sum sin^2 k_i<=.5102, Q_1^2>=.2499, Q_2^2,Q_3^2<=.0001;
% on supp T_4, sum sin^2 k_i>=2.9997 and every Q_i^2<=.0001.
% Thus p_parallel>p_perp and
% (p_parallel+2p_perp)/e_A >= .2499/(mu^2+.5102)
%                              > .0003/(mu^2+2.9997) >= 3p_0/e_B.
% For the first strict inequality, bound p_parallel below by
% 2n_A .2499/sqrt(mu^2+.5102) and p_perp above by
% 2n_A .0001/sqrt(mu^2+.49). These bounds hold for every mu>=0.
% The two anisotropic column combinations have eigenvalue
% p_parallel-p_perp, while the remaining energy/isotropic-pressure
% 2 by 2 determinant is nonzero by the displayed ratio inequality.
% Hence M(0) is invertible at every finite mu>=0. Its inverse and all
% needed q derivatives are bounded uniformly on any compact mass interval.
%
% Choose a small r_mu>0 on which
% ||M(0)^(-1)[M(q)-M(0)]||<=1/2 for |q|<=r_mu;
% existence follows from the proved determinant and smoothness, uniformly
% on each compact mass interval if desired. Fix a real smooth radial
% torus cutoff chi, one on |q|<=r_mu/2 and zero for |q|>=r_mu, and put
% Ghat(q)=chi(q) M(q)^(-1)M(0).
% This matrix multiplier is even in every coordinate, with Ghat(0)=I.
% Its convolution kernels have every absolute weighted lattice moment,
% sum_z G_jl(z)=delta_jl and sum_z z_i G_jl(z)=0. They are a preparation
% construction, not a modification of the already fixed averaging rule.
%
% UA7c. Positivity is proved after exact compensation, not assumed.
% Take the SAME positive band-limited phi of UA5a, f_L(x)=phi(x/L),
% with 24/L<r_mu/2, and w=(1/2,1/2,1/2,1/2)^T. Define explicitly
% F_j(x)=sum_l w_l (G_jl*f_L)(x).
% The relative second-difference estimate for the positive r^-12 envelope
% gives
% |F_j(x)-w_j f_L(x)|
% <=(C_phi/(2L^2)) f_L(x)
%     sum_l w_l sum_z |G_jl(z)| |z|^2(1+|z|)^12.
% All sums are finite. Choose L large enough that the right coefficient
% is <=1/4 for every j. Then f_L/4<=F_j<=3f_L/4<=3/4, and F_j is real
% and summable. The exact low-band identity is
% M(q) Fhat(q)=M(0)w fhat_L(q)
% on the whole sampled target support, with zero outside that support
% in the native Brillouin zone. No inverse-positivity assumption is used.
%
% Now set
% Gamma=lambda sum_(j=1)^4 P_+ T_j F_j T_j P_+, C_in=P0+Gamma,
% with 0<lambda<=1. Each summand is positive. Since the supports are
% disjoint, sum_j T_j^2<=I, whence
% 0<=Gamma<=3lambda P_+/4<=P_+.
% Thus C_in is a physical mixed Gaussian covariance, Gamma is trace
% class, and its mean particle count is finite. Time reversal fixes
% C_in and Gamma because T_j are real/even and F_j are real. Its
% EXACT native initial source and free current derivative are
% e_x=lambda rho_* f_L(x), J_x=0,
% dot J_i(q)=-i lambda p_* q_i fhat_L(q),
% rho_*=(3e_A+e_B)/2, p_*=(p_parallel+2p_perp+p_0)/2>0.
% The q_i in the last formula is the central transferred momentum,
% not a replacement of the declared current symbol Q_i(k).
% The identity holds on the retained low band; its lattice transform
% is periodic outside the native zone as usual.
%
% UA7d. Same exact radial readout and ACTUAL first-current tangency.
% Apply UA5a's K_ell point readout at its alias-free L,ell>24/pi.
% Its Fourier cutoff eliminates all noncentral copies exactly, giving
% bar rho(y)=lambda rho_* (K_(ell/L)*phi)(y/L)>0, bar J=0,
% dot bar J(y)=-(p_*/rho_*) grad_y bar rho(y).
% Both initial charges and this vector derivative are exactly radial.
% At the initial Gauss solution N=1 and
% v=-sqrt(2g M/r) rhat, M'=4pi r^2 bar rho, the flow contribution
% to dot J is identically zero by the SAME antiunitary argument as UA5d:
% j_x and h_v are odd, Gamma is invariant, and i[h_v,j_x] is odd.
% Consequently the displayed derivative is that of the actual full
% point-coupled Hamiltonian, for every g>0 and both momentum-row signs.
% Common reference transport gives dot D=-i[h,D] at this initial slice;
% no spectral refill or free-history substitution is used.
%
% Fix lambda=gamma/[g rho_* L^2] for any fixed gamma>0 and sufficiently
% large L. Then the initial geometry and source are exactly those of
% UA5c with A_L=gamma/(gL^2); in particular the same initial spatial
% Fourier, positive-mass/clock and geometric-length bounds apply. The
% exact count is
% N_f=lambda(3n_A+n_B) sum_x f_L(x)
%    =gamma L (3n_A+n_B)(int phi)/(g rho_*),
% using Ghat(0)=I and the Poisson identity. The same preparation works
% throughout the alias-free ell range; its total charge and N_f are
% independent of ell, while its interior depends on ell/L. Restore
% mu=ma, L=L_phys/a, ell=ell_phys/a and g=G/a^2. The filters explicitly
% select native momentum regions; this is not a cold-particle or a pure
% fixed-number preparation, and it is not an extra matter species.
%
% UA7e. What this construction settles and what remains required.
% The exact first-current obstruction in UA5 is preparation-specific:
% CAR positivity and the fixed noninjective readout do permit its evasion.
% This proves only the requested initial-current prerequisite. A full
% unprojected source history also needs the native dot rho tangent,
% higher consistency conditions, an actual coupled existence interval,
% and uniform clock/Fourier/time bounds and ell control on that interval.
% None follows from the initial compensation. The flow contribution to
% dot rho is not removed by the time-reversal argument used for dot J;
% it must be computed for this covariance and the solved initial v.
% UA8--UA9 below perform that computation and rule out energy tangency
% for this same family at sufficiently large fixed mu and large L.
% The displayed geometric constraint solution is not silently promoted
% to a solution of full H2 or to a history prepared from an earlier flat
% geometry. The reference at this initial time is P0, as in UA5.
% SOURCES: this construction uses the actual native source/current
% vertices, the already paired UA5 positive target/readout, and direct
% Fourier/covariance calculations here. UA7a--UA7e was saved independently
% before new RED or counterpart intake. No code or numerical run is used.

% UA8. Energy derivative on the SAME fixed preparation and solved flow.
% This continuation is derived before reading a new RED energy-tangent
% route. It does not re-solve F_j at later times. Work at the UA7 initial
% N=1, v(x)=V(x/L), with the actual radial Gauss V for the selected ell.
% Free h0 gives dot rho=0 by time reversal; the flow contribution remains.
% Let p be the momentum transfer in Gamma, l the transfer in v, and
% q=p+l the observed transfer (all native torus transfers are understood
% periodically). Put k_+=k+p/2, k_-=k-p/2 and
% R_j(k,p)=T_j(k_+)T_j(k_-)P_+(k_+)P_+(k_-).
% Define the exact real two-transfer kernel
% L_ij(p,l)=int Tr{[H_+(k,p,l) Z_+(k,p,l)
%                         -H_-(k,p,l) Z_-(k,p,l)] R_j(k,p)} dk/(2pi)^3,
% H_+=(h0(k_++l)+h0(k_-))/2, Z_+=(Q_i(k_++l)+Q_i(k_+))/2,
% H_-=(h0(k_+)+h0(k_--l))/2, Z_-=(Q_i(k_-)+Q_i(k_--l))/2.
% The Z's are scalar multiples of the identity. Directly evaluating
% -i[h_v,Gamma] at its two endpoints gives the exact native derivative
% dot ehat(q)=-i lambda sum_(i,j) int vhat_i(l) Fhat_j(q-l)
%                                      L_ij(q-l,l) d^3l/(2pi)^3.
% Apply Khat(ell q) to obtain the actual continuum point-readout derivative;
% it is not permissible to replace the two-transfer kernel by B_ij(q).
% The needed traces are explicit: with d(k)=(mu,sin k_1,sin k_2,sin k_3),
% Tr[h0(k') R_j(k,p)]=T_j(k_+)T_j(k_-)
%                          d(k').[d(k_+)/E_+ + d(k_-)/E_-].
% There is no odd three-Dirac-matrix trace. This makes L_ij a fixed
% finite Brillouin integral of the already specified state and flow.
% As checks, L_ij(p,0) gives the uniform-flow Q_i(k_+)-Q_i(k_-)
% energy vertex, and L_ij(-l,l)=B_ij(l). Therefore the exact total-energy
% identity is dot E=-sum_x v_i(x) dot J_i(x), using the native unaveraged
% current derivative and the point-evaluated v. No adjoint-smoothed flow
% or identification of an averaged product with this sum is inserted.
%
% UA8a. Controlled radial leading derivative, but not exact tangency.
% Define d_ij=2<E T_j^2 partial_(k_i)Q_i>. Coordinate reflections make
% the first Taylor term of the kernel
% L_ij(p,l)=d_ij p_i+(d_ij+p_ij)l_i+O((|p|+|l|)^3),
% p_ij=2<T_j^2 Q_i^2/E>.
% The remainder is cubic: simultaneous inversion of p,l makes the real
% kernel odd, and its derivatives are bounded because the preparation
% filters avoid the sine nodes. This holds uniformly on compact mass
% intervals, including mu=0. Since Ghat(q)=I+O(|q|^2),
% Fhat_j(q)=w_j fhat_L(q)+O(|q|^2)fhat_L(q) on the retained band.
% Cubic symmetry of the fixed four filters and w_j=1/2 gives constants
% d_* =sum_j w_j d_ij (independent of i),
% p_* =sum_j w_j p_ij as in UA7c.
% In the scaled coordinate X=y/L the actual derivative consequently is
% dot bar rho(y)=-(lambda/L) K_delta *
%       [d_* V.grad_X phi+(d_*+p_*)phi div_X V](X)
%                       + E_L(X), delta=ell/L,
% ||E_L||_(C_X^k)<=C_(k,mu_*,gamma,delta0) lambda/L^3
% for every fixed k, uniformly on the stated alias-free ell/L<=delta0
% range and compact mass interval. The displayed leading function is
% exactly radial, since K_delta, phi and the actual V are radial.
%
% To justify this as a bound for the actual sampled flow, not a continuum
% replacement, use the uniformly weighted Fourier measure of V proved
% in UA5c. The target has |p|<=24/L. Taylor's cubic remainder, the even
% Ghat-I term, and k scaled derivatives are bounded by Fourier moments
% through k+3 (or any larger fixed order). Sampling v folds transfers
% l+2pi n into the native torus. Every noncentral copy has macroscopic
% |L(l+2pi n)|>=cL on the retained output/target range, and those same
% higher Fourier moments bound its contribution by any prescribed
% inverse power of L. The point filter has |Khat|<=1 and needs no
% derivative in this argument. Thus native aliases are controlled, not
% discarded. The scaled initial V bounds are uniform in delta as above.
%
% The estimate controls an initial angular residual at order (a/L_phys)^2
% relative to the characteristic leading derivative. It DOES NOT show
% that the exact E_L is radial or zero. Exact first-source tangency now
% reduces to the angular part of the fully specified two-transfer
% expression above; no additional freedom to refit F_j in time exists.
% UA9 below supplies a nonzero witness for this same fixed family at
% sufficiently large fixed mu and large L. The other-mass angular
% question and the general coupled-interval/history requirements remain
% open; this initial-source control does not replace the full goal.
% SOURCES: UA8 follows by direct native Fourier endpoint multiplication
% on the saved UA7 covariance. No new RED energy calculation, counterpart
% source, code or numerical evaluation was used to obtain it.

% UA9. Energy-tangent witness for the SAME UA7 weights and fixed bump.
% This derives the coefficient for UA7's actual T_4 amplitude one and
% r^-12 envelope; a different stationary-filter weight or target exponent
% cannot supply its sign. The heavy-mass calculation was worked out from
% UA8 before reading RED AC18--AC20. RED has announced a related route;
% this saved extension makes no additional blind-intake claim about it.
%
% First extract the third-order spatial term of the EXACT UA8 expression.
% The odd kernel has a Taylor expansion through cubic order with fifth-
% order remainder, while Ghat is even. The same Fourier estimates give
% dot bar rho=(lambda/L) R_(0,delta,mu)
%             +(lambda/L^3) R_(2,delta,mu)+O_(C_X^k)(lambda/L^5).
% The first term is the radial function in UA8a. All coefficients are
% determined by the fixed initial preparation and initial Gauss V_delta.
% For ell=L^theta, 0<theta<1, delta->0 and the cubic coefficient and its
% angular part converge in each fixed local C^k norm to delta=0, since
% K_delta is a fixed approximate identity and V_delta has the proved
% uniform derivatives/Fourier moments. The remainder is uniform over
% these alias-free scales. This expansion neither refits F_j nor evolves
% an arbitrarily selected later state.
%
% It suffices to exhibit a nonzero angular part of R_(2,0,mu) at some
% sufficiently large FIXED mu. Rescale the energy row of M by mu^-1
% and its three current rows by mu. The resulting symbols, their finite
% q derivatives, their inverse, and L_ij/mu have convergent expansions
% in mu^-2 at large mu. Their limiting M(0) is nonsingular by UA7b.
% Thus R_(2,0,mu)/(2mu) converges to the following explicit scalar-band
% calculation. This is a controlled heavy-mass coefficient limit, not
% an assumption that the exact finite-mass covariance is spinless.
%
% Let b=b_eta be the EXACT UA7a one-dimensional bump. For small p define
% R(p)=int cos(2z)b(z+p/2)b(z-p/2)dz / int b(z+p/2)b(z-p/2)dz,
% u(p)=[1-R_4(p)]/8, with R_4 the same expression with cos(4z).
% Write R=R_0+R_2 p^2+O(p^4), u=u_0+u_2 p^2+O(p^4),
% d_u=1/4-2u_0>0. The conditional momentum moments are
% <Q_i^2>=(1+R_4)/8=1/4-u on an axis-center coordinate pi/4,
% and u on a coordinate centered at zero or pi/2. Their cos(2k_i)
% moments are zero, +R and -R, respectively. All these identities are
% exact for this bump, not a zero-width replacement. The product overlap
% common to all four filters cancels against the compensated F_j.
%
% Put a_c=n_A/2, b_c=n_B/2=4a_c and t_c=3a_c+b_c=7a_c; the ratio four
% follows from two axis patches versus eight stationary patches, with
% T_4 amplitude one exactly as fixed in UA7a. Let
% P_c=d_u a_c+u_0 t_c,
% zeta=[(2/3)P_c-u_2(t_c-2a_c)]/d_u.
% In the scalar-band limit, the compensated effective occupation of
% axis channel i is
% s_i(p)=[P_c/g(p_i)-u(p_i)t_c]/[1/4-2u(p_i)]
%       =a_c+zeta p_i^2+O(|p|^4),
% g(p)=sin p cos p/p. The stationary occupation is t_c-sum_i s_i.
% Here each s multiplies the common fhat_L; this formula follows by
% solving the exact limiting four-row system, not by a new preparation.
%
% The limiting number/energy-flow kernel is particularly simple. With
% q=p+l, its component i, divided by 2mu, is
% sin(q_i)cos(l_i) R(p_i)
%                     [2 sum_(j!=i)s_j(p)+s_i(p)-t_c].
% To derive the sine factor, the exact scalar-flow endpoint difference is
% [Q_i(k+p/2+l)+Q_i(k+p/2)-Q_i(k-p/2)-Q_i(k-p/2-l)]/2
% =cos(2k_i) sin(p_i+l_i) cos(l_i).
% Through cubic order the bracket including R is
% D+2zeta R_0 |p|^2+C p_i^2,
% D=(2a_c-b_c)R_0=-2a_c R_0,
% C=(2a_c-b_c)R_2-zeta R_0=-2a_c R_2-zeta R_0.
% The |p|^2 term is rotationally invariant. The remaining possible
% nonradial third-order differential expression, after inverse Fourier
% transformation, is
% -D/6 sum_i partial_i^3(V_i phi)
% -D/2 sum_i partial_i(phi partial_i^2 V_i)
% +C   sum_i partial_i(V_i partial_i^2 phi),
% multiplied by 2mu; omitted terms are radial for radial V,phi.
%
% The signs of the finite-width bump coefficients are rigorous. Let
% ell_b(z)=log b(z), and normalize the measure b(z)^2 dz. For any even
% observable f, the p^2 coefficient of its conditional average is
% (1/4) Cov(f,ell_b''). On 0<=|z|<eta, ell_b'' is strictly decreasing
% with |z|, as are cos(2z) and cos(4z). The elementary covariance identity
% Cov(f,g)=(1/2)E[(f(Z)-f(Z'))(g(Z)-g(Z'))] therefore gives
% R_2>=0 and u_2<=0. All endpoint integrations converge because the
% smooth bump vanishes faster than every power there. Also R_0>0,
% P_c>=d_u a_c, hence zeta>=2a_c/3.
%
% To show the differential expression is actually nonradial, use the
% proved far-field asymptotes of the SAME functions:
% phi(X)=phi_infty r^-u+rapidly decaying remainder, u=12, phi_infty>0,
% V_i(X)=-v_infty X_i r^-t+O(r^-19/2), t=3/2, v_infty>0.
% Differentiated remainders obey the corresponding bounds. Put s=u+t.
% The degree-four cubic harmonic H_4(n)=sum_i n_i^4-3/5 in the preceding
% expression has leading coefficient
% v_infty phi_infty(s+4)
%   {-D[s(s+2)/6+t(t+2)/2]+C u(u+2)} r^(-s-2) H_4(n).
% For the actual u=12,t=3/2 this brace equals
% -(75/2)D+168C
% =75a_c R_0-336a_c R_2-168zeta R_0
% <=-37a_c R_0-336a_c R_2<0.
% No sign from a different stationary-filter weight or r^-8 target is
% transferred into this computation. The nonzero far coefficient gives
% a finite-radius angular witness for the limiting cubic term. Continuity
% of the normalized finite-mass coefficients then gives such a witness
% for all sufficiently large fixed mu, without claiming a numerical
% threshold or an all-mass energy-tangency theorem.
%
% Fix one such mu, and then take L large with ell=L^theta, 0<theta<1.
% The approximate-identity convergence preserves the finite-radius
% angular witness in R_(2,delta,mu), while the O(lambda/L^5) remainder
% is smaller than its nonzero lambda/L^3 contribution. The actual
% dot bar rho is consequently NOT radial for this fixed UA7 preparation,
% although bar rho, bar J and dot bar J are exactly radial. This is an
% energy-tangency obstruction at finite but sufficiently large L and
% fixed ma, for either momentum-row sign and any fixed g>0 with the
% stated lambda=gamma/(g rho_* L^2). It forbids a differentiable exact
% unprojected spherical interval through these data. A small angular
% remainder, other masses, other physical covariances, or an explicitly
% controlled approximate coupled history are not excluded. Those remain
% substantive requirements of the full goal, not consequences of the
% first-current construction or of this preparation-specific witness.
% SOURCES: direct UA7--UA8 heavy-mass row normalization, exact product-bump
% conditional moments, Fourier Taylor expansion and native Gauss tail.
% No new numerical result or external mainline agreement is asserted.

% ACTUAL FOUR-CHANNEL / ENERGY-TANGENCY COMPARISON [R].
% Read the complete saved RED AC14--AC20 and their SOURCES after the
% independent author saves recorded above, including its separate exact
% two-transfer derivation. Its four-channel matrix inverse and weighted
% positive-envelope route independently give a physical finite-count
% initial covariance with radial rho,J,dot J under the fixed readout.
% The exact energy kernels agree after transfer relabeling: its
% L(q,-q)=-B(q) is our L(-l,l)=B(l), since B is odd. Both derive the
% point-paired identity dot E=-sum_x v_i dot J_i, not a product of
% separately averaged fields. The pressure term and the full four-spin
% trace in UA8 agree with the independently evaluated endpoint matrices.
%
% The two explicit preparations are not identical. RED AC14 halves its
% stationary T4 amplitude and uses an r^-8 envelope; UA7 has unit T4
% amplitude and the fixed UA5 r^-12 envelope. Its first heavy-mass sign
% therefore cannot be transferred to UA7. RED then independently computed
% UA7's actual coefficients: axial fractions 1/7, stationary fraction
% 4/7, normalized a0=-2c0/7 and a_an=-2c2/7-c0 u2, with
% u2=(7-5d0)/(21d0)+5d2/(14d0). Here c0=R_0, c2=R_2,
% d0=R_4(0), d2 the p^2 coefficient of R_4. Direct substitution into
% UA9 gives D/t_c=a0 and C/t_c=a_an, since t_c=7a_c.
% Its H4 brace 168a_an-(75/2)a0<=-37c0/7 is precisely UA9's
% unnormalized bound divided by t_c. The finite-width covariance signs,
% row-rescaled finite-mass continuity, finite-radius witness, approximate
% identity and native alias estimates agree. The witness annulus lies
% in the nonzero algebraic source tail, not in a source-free exterior.
% It excludes a differentiable exact spherical interval through this
% fixed family only in the stated large-fixed-mass/large-L domain.
% No coupled history or universal obstruction is inferred. Agreement
% inside this original pair is [R], not external [X].

% ACTUAL NEW E/Q SCOPE INTAKE [S; no imported history or count].
% After the independent energy derivation, read the complete actual
% strand_E_angular_constraint_completion.tex and its _scope.tex, and
% strand_Q_native_state_count.tex and its _scope.tex in
% <internal>/solve_b, including SOURCES.
% Both explicitly retain pending independent review. E AG7--AG18 is a
% bounded initial Hessian-constraint completion with extra constant
% Dirichlet data and HA adjoint coupling, conditional on its specified
% HT/CK preparation estimates. Its scalar sigma_2 expansion and local
% elliptic contraction do not provide a spherical point-coupled tangent,
% exterior matching or common history. AG19 is a finite-boundary charge
% identity; AG20--AG21 retain rather than eliminate angular flow.
% Q NC3--NC5 constructs an actual finite-parent spin subspace with common
% initial sources. NC6 explicitly projects the density radially; NC7--11
% use HA coupling to demonstrate leakage, and NC12--14 bound only a hard
% positive free-energy spectral window. That bound is not an invariant
% flowing-charge sector or accessible object count. These direct scope
% checks agree with the actual RED intake; neither supplies the missing
% unprojected point-coupled interval here. Unread prerequisite results,
% independent-review status, and the differing source/coupling domains
% are not promoted by this comparison.

% FINAL FOUR-CHANNEL READ / SOURCES.
% Both complete current author and RED sources and all SOURCES were
% read in untruncated ranges, including AC14--AC20, the actual UA9
% coefficient comparison, both compact statements and the new E/Q
% scope intake. The fixed four-channel CAR construction passes exact
% initial-current tangency at every finite ma; the same state's exact
% two-transfer energy derivative fails radial tangency at sufficiently
% large fixed ma and large L/a with ell/a=(L/a)^theta, 0<theta<1.
% Both momentum signs and the native point/common-transport allocation
% are retained. This excludes a differentiable exact spherical interval
% through these data, not all covariances or a controlled approximate
% history. The normalized harmonic coefficients agree only after the
% actual patch weights and r^-12/r^-8 targets are distinguished. No
% mandatory scientific repair remains at the stated [R] scope. The
% conditional transported-area theorem and existing homogeneous [X]
% scope are unchanged; a full coupled history is still required by
% the goal and is not inferred from this initial preparation.

% A BROADER PHYSICAL COVARIANCE: FULL FIRST-SOURCE TANGENCY [R].
% UA10 is derived and saved before reading any new RED continuation.
% It keeps UA5a's fixed positive noninjective K and POINT readout, the
% native h0,Q, both printed momentum signs and common transported sea.
% The physical preparation now contains particles AND holes, with finite
% expected excitation count. It is not a positive particle-only D and
% is not claimed to be cold matter or a pure fixed-number state.

% UA10a. Fix an explicit covariance and prove CAR, rather than infer it.
% Work on Z^3 with a=1, any fixed mu=ma>0, g=G/a^2>0, and the SAME
% phi and K of UA5a. Put F_L(x)=phi(x/L), a positive scalar diagonal
% multiplier with 0<F_L<=1, and H=h0/E=sign(h0), P0=(I-H)/2.
% Define once at the initial time
%   C_L=P0+nu_L F_L beta, D_L=nu_L F_L beta,
%   nu_L=gamma/(4g mu L^2), gamma>0 fixed.
% This is not assumed physical merely because its energy is positive.
% Here is a sufficient exact operator inequality establishing it for
% every fixed mu>0 and sufficiently large L.
%
% Let s_*=mu/sqrt(mu^2+3)>0. The massive torus symbol H is analytic
% in a complex strip: compactness and E^2>=mu^2 on the real torus
% give a nonzero strip and a consistent analytic square root. Thus its
% kernel obeys ||H_z||<=C_mu exp(-c_mu |z|). The proved bound
% |grad log phi|<=C_phi gives
% |sqrt(F_L(y)/F_L(x))-1|
% <=(C_phi |y-x|/(2L)) exp(C_phi |y-x|/(2L)).
% The row AND column Schur tests therefore extend
% B_L=F_L^(-1/2) H F_L^(1/2)-H
% from finitely supported vectors to a bounded operator with
% ||B_L||<=b_L:=(C_phi/(2L))sum_z |z| ||H_z||
%                                  exp(C_phi |z|/(2L))=O_mu(L^-1).
% Since {H,beta}=2mu/E>=2s_* I, the exact quadratic-form identity is
% F_L^(-1/2){H,F_L beta}F_L^(-1/2)
%       ={H,beta}+B_L beta+beta B_L^dagger
%       >=2(s_*-b_L) I.
% Choose L so b_L<=s_*/2 and nu_L<=s_*/4. With
% A_L=H-2nu_L F_L beta=A_L^dagger, this gives
%   I-A_L^2=2nu_L{H,F_L beta}-4nu_L^2F_L^2
%          >=nu_L s_* F_L>=0.
% Hence ||A_L||<=1 and C_L=(I-A_L)/2 obeys 0<=C_L<=I exactly.
% In fact C_L(I-C_L)>=(nu_L s_*/4)F_L, so C_L and I-C_L have
% that same positive quadratic-form lower bound. Inverse powers of F_L
% in the proof are only form devices; the actual covariance is bounded.
% No finite positive lower spectral margin on the whole infinite parent
% is alleged, since F_L decays. The massless ansatz is actually excluded,
% not merely unproved: any nonzero trace-class nu F_L beta at mu=0 has
% zero normal energy, contrary to UA1a's strict positive free-excitation
% energy for every nontrivial physical C-P0. This strengthening was
% checked after reading the independently saved RED AC21 below.

% D_L is trace class, Hermitian and time-reversal invariant. Its net
% added charge Tr D_L=0. Its actual positive-particle and hole counts
% relative to the flat incoming sea are equal and finite:
%   N_+=Tr(P_+ C_L)=2nu_L mu<1/E> sum_x F_L(x),
%   N_-=Tr[P_-(I-C_L)]=N_+,
%   N_ex=N_++N_-=4nu_L mu<1/E> L^3 int phi
%                    =gamma L<1/E>(int phi)/g.
% CAR proved above makes these nonnegative counts; their trace formulas
% use the actual flat-band diagonal kernel, not an imposed band filling.
% Off-diagonal particle-hole coherences of D_L are retained. Expected
% positive particles and expected holes are distinguished from net charge.

% UA10b. Exact radial charges and BOTH actual first-source derivatives.
% Direct native matrix multiplication gives at every site
%   e_x=Tr(e_x^op D_L)=4mu nu_L F_L(x), J_(x,i)=0,
%   dot e_x=0, dot J_(x,i)=0
% under the ACTUAL initial h=h0+sum_i{v_i,Q_i}/2, for every real point
% flow v. To verify the derivatives without a continuum replacement,
% first use i[h0,e_x^op]=(i/2)[h0^2,Pi_x]. The latter is spin scalar,
% so its trace against the site-diagonal beta D_L vanishes. The flow
% commutator [h_v,e_x^op] has one two-step hop and either an onsite
% mass or one nearest-neighbor hop; it has no diagonal site block to
% pair with D_L. For the current derivative, h0's one-step/mass terms
% cannot close a two-step current bond at the diagonal. The h_v/current
% diagonal products are spin scalar, whose trace against beta is zero.
% These are exact support/Clifford identities, not leading derivatives.
% Common transport means dot D=-i[h,D]; no reference term is omitted.
%
% For UA5c's alias-free scales the SAME point readout gives exactly
%   bar rho(y)=gamma/(gL^2) (K_(ell/L)*phi)(y/L)>0,
%   bar J=dot bar rho=dot bar J=0.
% The initial Gauss fields are N=1 and v=-sqrt(2g M/r)rhat,
% M'=4pi r^2 bar rho, with the regular center and normalization already
% proved in UA5c. They coincide for the two momentum signs. This is a
% genuine FULL FIRST-SOURCE-TANGENT initial preparation, unlike UA7's
% current-only pass. No change in F_L or the covariance is made in time.
% Its reference is P0 at the specified initial slice; no earlier collapse
% preparation from a flat geometry has been constructed by this identity.
% Its total energy is gamma L(int phi)/g, independent of ell at fixed
% preparation; its interior still depends on ell/L. Restore L=L_phys/a,
% ell=ell_phys/a and mu=ma. All conditions are expressed through mu,g,
% gamma,L and ell; no microscopic geometric law is inserted.

% UA10c. Test the next necessary tangent on the SAME fixed state.
% The first-source pass is not yet a coupled interval. For any twice-
% differentiable exactly spherical solution through these positive
% data, differentiating the radial Gauss rows gives dot M=dot N=dot v=0
% initially: dot rho=dot J=J=0 and M(r)>0 for r>0, with regular center
% and N(infinity)=1. Thus dot h=0 there for either momentum sign.
% Its actual second source derivative must consequently be the native
% double commutator with that same initial h, not a refitted state.
%
% For the energy source the pure-h0 and mixed h0/h_v double commutators
% have zero expectation in D_L. For the pure term use the scalar first
% commutator above; its h0 commutator has only one alpha matrix, whose
% trace with beta is zero. In the mixed terms, the scalar/grade-two
% Clifford traces vanish in the same way. Only the scalar flow squared
% transports the diagonal beta occupation. Since
% (h_v)_(x,x+2e_i)=[v_i(x)+v_i(x+2e_i)]/(8i), exactly
%   ddot e_x=(1/32)sum_(i,sign=+/-)
%       [v_i(x)+v_i(x+sign 2e_i)]^2 [e_(x+sign 2e_i)-e_x].
% This follows directly from
% -[h_v,[h_v,F_L]]_(xx)=2sum_y |(h_v)_(xy)|^2(F_L(y)-F_L(x)).
% No continuum isotropy has been imposed on this two-step expression.

% Put e_x=rho_amp phi(x/L), rho_amp=gamma/(gL^2), X=x/L,
% and v(x)=V_delta(X), delta=ell/L. Taylor expansion of the exact
% paired plus/minus stencil, then the actual K readout, gives
%   ddot bar rho=(rho_amp/(2L^2)) K_delta*
%       sum_i partial_(X_i)[V_(delta,i)^2 partial_(X_i)phi]
%       +O_(C_X^k)(rho_amp/L^4).
% The error is uniform on the same initial spatial regularity domain;
% reciprocal copies are controlled by the already proved weighted
% Fourier bounds of the sampled V_delta. The leading expression is
% generally NOT radial, despite radial phi and radial vector V_delta.
%
% A rigorous witness follows in the limit delta->0 from the SAME tails
% used in UA9: phi=phi_infty r^-12 plus rapidly decaying remainder,
% V_i=-v_infty X_i r^-3/2+O(r^-19/2), with positive coefficients.
% For phi=F r^-u, u=12, and V_i=-V0 X_i r^-3/2, the H4(n) coefficient
% in sum_i partial_i(V_i^2 partial_i phi) is
%   u(u+5) F V0^2 r^(-u-3) H4(n), H4(n)=sum_i n_i^4-3/5.
% It is strictly nonzero. The differentiated tail remainder has higher
% inverse powers, hence a fixed finite-radius annulus detects it.
% Taking ell=L^theta, 0<theta<1, preserves that witness by the fixed
% approximate identity and V_delta->V_0. The O(rho_amp/L^4) remainder
% is then smaller than the nonzero rho_amp/L^2 angular term at large L.
% Thus for EVERY fixed mu>0, g>0 and gamma>0, sufficiently large L in
% this explicit family gives nonradial ddot bar rho. It excludes a twice-
% differentiable exact spherical coupled interval through these data.
% It does not exclude other physical covariances, the massless case of
% a different construction, or a controlled approximate history. The
% initial first-source success and the second-source failure are both
% calculated on one fixed CAR preparation, without any refitting.
% SOURCES: current GOAL_alt.md reread in full; actual native source
% vertices, common transport, the previously fixed UA5 kernel/profile,
% and the independent sign-projector, Schur-form and support/Clifford
% calculations supplied here. No new RED/counterpart source, numerical
% implementation or run was read or used for this substantive save.

% ACTUAL COHERENT-PREPARATION COMPARISON [R].
% Read the complete actual RED AC21--AC24 and SOURCES after both
% substantive independent continuations were saved. Its direct form
% W^-1 C(1-C) W^-1=mu/E+R_L-F is equivalent to UA10a's squared-sign
% estimate. Both row and column bounds use the same exponentially
% decaying massive sign kernel and positive logarithmic-gradient control.
% The sufficient constants differ but each proves exact CAR, not merely
% formal positivity of the leading local symbol. The zero-mass exclusion
% follows from the previously paired strict CAR energy theorem, as above.
%
% Equal finite positive-particle/hole counts, zero net charge, the full
% native first-source cancellation and the second double commutator
% agree. In particular the hopping amplitude (v_x+v_y)/(8i) gives
% exactly the same 1/32 stencil and one-half continuum coefficient.
% The two fixed readouts/envelopes remain distinct: RED retains its
% r^-8 target, while UA10 retains the r^-12 UA5 target and kernel.
% With H4(X)=r^4 H4(n), its coefficient
% 104 gamma m_infinity c_f r^-15 H4(X) is precisely the u=8 case
% of UA10c's general formula after the one-half prefactor and
% V0^2=2gamma m_infinity. UA10's u=12 coefficient instead equals
% 204 gamma m_infinity phi_infinity r^-15 H4(n).
% Both fixed finite-radius source-tail witnesses survive delta->0 and
% the controlled native Taylor/alias remainder. No source-free exterior
% is substituted, and no covariance is refitted between derivatives.
%
% This supplies physical full-first-source-tangent initial data and a
% second-source obstruction for that SAME preparation at every fixed
% ma>0, g>0 and sufficiently large L/a in the stated ell range. It
% supplies no actual prior collapse preparation or invariant coupled
% interval. Other states and a controlled angular approximation remain
% open; agreement within this original pair does not create external X.

% FINAL ACTUAL COHERENT-SOURCE READ / SOURCES [R].
% Both complete current author and RED sources and every SOURCES block
% were reread, with the truncated display range recovered separately.
% This includes the exact saved compact cores, UA10a--c, AC21--AC24,
% both actual comparisons and RED's final coherent-source disposition.
% The massive form-CAR bounds, finite expected particle and hole counts,
% zero net charge, initial-P0 domain and massless-ansatz exclusion agree.
% The exact first-source cancellation and 1/32 second-density stencil
% use one fixed physical covariance, the actual point Hamiltonian and
% common transported reference; both momentum-row signs are retained.
% The separate r^-12 and r^-8 targets give their own nonzero harmonic
% coefficients after the stated harmonic-convention conversion. The
% finite-annulus, fixed-positive-mass, large-L and ell=L^theta limits
% exclude a C^2 exact spherical interval through these data only.
% They supply neither an earlier incoming collapse preparation nor an
% invariant domain for other states or a controlled approximate history.
% The conditional area theorem and prior homogeneous/slab X scope are
% unchanged. No mandatory scientific repair remains at the stated R
% scope; author source is READY/HELD for the original pair's handoff.

% UA11. A BROADER COVARIANCE CLASS AND SOURCE-DOMAIN NONINVARIANCE [R].
% Independent substantive continuation, saved before new RED or KM/EL/CP
% counterpart intake. The published compact core is unchanged. Work on
% the infinite native parent, a=1, fixed mu=ma>0, g=G/a^2>0, with the
% SAME UA5 positive radial K and point readout, and common incoming P0.
% Fix gamma>0 and the SAME phi, L and ell=L^theta, 0<theta<1. Nothing
% below changes the source allocation, Gauss sign or Hamiltonian in time.

% UA11a. Exact locality class, specified before testing the constraints.
% Let D=C-P0 be any Hermitian trace-class difference with 0<=C<=I,
% subject to the following explicit native matrix conditions:
% (i) D_(xy)=0 whenever 1<=|x-y|_1<=5;
% (ii) mu Tr_spin(beta D_(xx))=rho_amp phi(x/L),
%      rho_amp=gamma/(gL^2).
% The diagonal 4 by 4 blocks are otherwise arbitrary Hermitian matrices;
% in particular neither Tr_spin D_(xx)=0 nor D_(xx) proportional to beta
% is imposed. All matrix elements at separation >=6 are unrestricted
% apart from the stated physicality and trace-class conditions. This is
% a class of actual covariances, not a restriction on the source operator
% or on its positive spectral subspace. It includes nonlocal coherences
% and unequal finite expected particle/hole counts.

% This class is nonempty and much larger than UA10's single ansatz.
% Its explicit physical subset is constructed from that proved covariance
% C_*=P0+nu F beta, F_x=phi(x/L), nu=gamma/(4g mu L^2), by
%   C=C_*+nu zeta F^(1/2) T F^(1/2),
% with ANY bounded Hermitian T satisfying ||T||<=1,
% T_(xy)=0 for 1<=|x-y|_1<=5, and Tr_spin(beta T_(xx))=0.
% Choose 0<=zeta<=s_*/8, s_*=mu/sqrt(mu^2+3), and the UA10 large-L
% CAR conditions. Since C_*,I-C_*>=(nu s_*/4)F as quadratic forms,
% C,I-C>=(nu s_*/8)F>=0. Thus CAR is exact. The added term is trace
% class by the Hilbert--Schmidt ideal product, with norm <=4nu zeta
% sum F. For example T can contain any bounded sitewise spin matrix
% orthogonal to beta, and Hermitian shift coherences of length >=6,
% with total operator norm <=1. Both types may be nonzero at once.
% Their free energy is exactly zero as perturbations: h0 has range one
% and the diagonal beta trace is zero. The TOTAL state retains the
% positive energy gamma L(int phi)/g and has finite expected excitation
% count N_ex<=||D||_1<=4nu(1+zeta)sum F. No pure fixed-number, cold-state
% or earlier-collapse preparation is asserted. This construction proves
% physical membership; the obstruction below covers the entire class
% (i)--(ii), not only this small weighted perturbation subset.

% UA11b. All those covariances have the SAME first and second source jets.
% At N=1 the native h=h0+V has maximum hopping length two, the energy
% vertex e_x has length one, and the current j_xi has length two.
% Consequently e_x,j_xi,i[h,e_x],i[h,j_xi], and [h,[h,e_x]] have no
% matrix elements beyond separations 1,2,3,4,5 respectively. The distant
% coherences in (i) cannot contribute to ANY of these traces, exactly.
% It remains to evaluate them against arbitrary diagonal spin blocks.

% The diagonal of [h0,e_x]=[h0^2,Pi_x]/2 is zero. The diagonal of
% [V,e_x] is zero because a two-step hop cannot close a one-step or
% onsite energy vertex. For the current, [h0,j_xi] has zero diagonal
% by hop parity. Both V and j_xi are spin-scalar purely imaginary
% antisymmetric matrices with real point-flow coefficients, so their
% commutator has zero diagonal as well. Thus every covariance in the
% class obeys e_x=rho_amp phi(x/L), J_x=dot e_x=dot J_x=0; no vanishing
% diagonal spin trace was used. The same alias-free readout therefore
% gives exactly UA10b's positive radial rho,J and zero first derivatives.
% Both radial Gauss signs give the same regular initial N=1 and v.

% For the second energy derivative, the pure-h0 term has zero diagonal:
% its first commutator is a scalar two-step matrix with zero diagonal,
% and h0 has only onsite/one-step terms. The mixed term
% [V,[h0,e_x]] also has zero diagonal: its two scalar spatial factors
% are respectively imaginary antisymmetric and real antisymmetric.
% For the other mixed term, use Jacobi to reduce its diagonal to
% [[h0,V],e_x]. A returning path must contain one two-step V hop and
% two one-step h0/energy hops on the SAME axis. The two spin factors
% multiply to alpha_i^2=I; the three hopping coefficients are imaginary.
% This double commutator is Hermitian, so its purely imaginary scalar
% diagonal must vanish. Onsite mass terms cannot make such a path.
% These exact spatial/Clifford statements hold for arbitrary diagonal
% D_(xx), including nonzero spin trace. In the two-V term only the
% onsite mu beta part of e_x can return; its one-step alpha part has
% odd parity. Therefore the exact second density is STILL
%   ddot e_x=(1/32)sum_(i,s=+/-1)
%       [v_i(x)+v_i(x+2s e_i)]^2 [e_(x+2s e_i)-e_x].
% Here dot h=0 follows on any hypothetical differentiable regular Gauss
% branch from the common zero first sources, exactly as in UA10c.
% No extra reference term occurs because dot D=-i[h,D]. Thus the entire
% class has the same actual second density, not just its leading symbol.

% UA11c. A precise failure of source-domain invariance, not just an ansatz.
% UA10c's fixed readout/Taylor/alias proof applies to these identical
% source jets. At every fixed mu,g,gamma>0 and sufficiently large L,
% the nonzero H4 component of ddot bar rho excludes a C^2 exact spherical
% interval through EVERY covariance satisfying (i)--(ii), for both
% momentum-row signs. Arbitrary distant coherences and arbitrary onsite
% spin structure preserving the energy cells do not evade the result.
% This is an infinite-dimensional physical class with finite expected
% count, but is not all physical covariances. In particular a proposed
% repair preserving these initial charges and first derivatives must
% alter a relevant off-diagonal matrix element at native separation <=5;
% changing only occupation/spin data invisible to these vertices cannot
% repair the second tangent. Even such a near-diagonal change would
% still need all subsequent tangents and a coupled existence proof.

% A profile-general leading obstruction is also available. If the fixed
% positive radial target phi is replaced by a smooth nonnegative finite-
% mass radial f for which the SAME exact radial readout and controlled
% Taylor/alias passage have been proved, write V_i=w(r)X_i and
% m(r)=4pi int_0^r s^2 f(s)ds, w^2=2gamma m/r^3.
% The H4(X)=sum X_i^4-3r^4/5 coefficient of
% (1/2)sum_i partial_i(V_i^2 partial_i f) is exactly
%   (1/(2r)) d/dr [w(r)^2 f'(r)/r].
% If this vanished at every r>0, then m(r)f'(r)/r^4 would be constant.
% Under the stated decaying smooth finite-mass profile bounds that
% constant is zero. Wherever m>0 one would have f'=0, incompatible
% with a nonzero nonnegative finite-mass smooth target. Thus some
% finite annulus has a nonzero harmonic. This last observation is a
% conditional extension to profiles with the named readout/error proofs;
% the fixed UA5 phi already satisfies them and needs no new premise.
% No universal all-state no-go or controlled approximate-history failure
% follows. The actual incoming-history and full invariance requirements
% remain unresolved outside this explicitly characterized class.
% SOURCES / INDEPENDENCE: binding GOAL_alt.md reread in full;
% previously fixed UA5 readout and UA10 physical form/tail proof;
% independent finite-range diagonal-block and weighted-CAR calculation
% above. New RED continuation and new KM/EL/CP sources were not read
% before this substantive save. No scientific code or run was used.

% UA11d. Actual independent comparison and a global angular witness [R].
% After the independent UA11 save, read the complete RED AC25--AC27 and
% SOURCES. Its separately derived arbitrary SITE-DIAGONAL class, exact
% first cancellations, mixed-commutator identities and 1/32 second stencil
% agree with UA11b. UA11's additional separation>=6 coherences are an
% independently proved enlargement by the actual vertex ranges, not a
% change in the RED preparation. The two weighted CAR balls have different
% sufficient constants, each valid; finite particle and hole counts, not
% net charge, are retained. RED's second proof uses an integrated harmonic
% witness, which is checked explicitly next and applies to our fixed K.

% Put delta=ell/L, distinct from the preparation amplitude zeta above.
% Write H4(X)=sum_i X_i^4-3r^4/5 and
% T_delta=(1/2)sum_i partial_i(V_delta,i^2 partial_i phi),
% V_delta,i=w_delta(r)X_i, w_delta^2=2gamma m_delta(r)/r^3.
% Integration by parts gives
% I_delta=int H4 T_delta d^3X
%   =-(96pi/175)int_0^infinity r^7 w_delta^2 phi'(r)dr
%   =(192pi gamma/175)int_0^infinity
%        phi(r)[4r^3 m_delta(r)+r^4 m_delta'(r)]dr >0.
% Indeed partial_i H4=4X_i^3-(12/5)r^2X_i, while
% <sum n_i^6>=3/7 and <sum n_i^4>=3/5. Their angular combination is
% 4(3/7)-(12/5)(3/5)=48/175, fixing the first coefficient including
% the outer 1/2. The second equality uses w_delta^2=2gamma m_delta/r^3.
% Boundary terms vanish: the center is regular and phi has r^-12 tails.
% Both m_delta>0 and m_delta'>=0 are properties of the actual positive
% readout. No monotonicity of phi or a selected far-tail sign is needed.
% Continuity in delta and the uniform initial bounds in UA5c give
% inf_(0<=delta<=delta0) I_delta>0 for every fixed finite delta0.

% Our K=phi/int phi differs from RED's Schwartz K, but it is normalized,
% radial and has a finite fourth moment, which are exactly the required
% properties for the harmonic-mean identity
%   int K_ell(y-x) H4(y/L)d^3y=H4(x/L).
% Spherical averaging of a harmonic polynomial gives its center value;
% its polynomial growth and the fourth moment justify the integral.
% Thus the actual point readout preserves this witness EXACTLY, not just
% as ell/L tends to zero. The initial fields are still point evaluated.

% The exact stencil's weighted error is bounded by
% sum_x [1+|x/L|^4] |ddot e_x-rho_amp L^-2 T_delta(x/L)|
%       <=C rho_amp/L,
% uniformly for 0<=delta<=delta0. Taylor's paired-hop remainder is
% rho_amp L^-4 times a smooth integrable weight; at infinity its terms
% decay at least r^-17, since V_delta^2=O(r^-1) and phi^(4)=O(r^-16),
% with the corresponding product-derivative bounds. At the center all
% coefficients are smooth and uniformly bounded. The weighted sum is
% O(L^3), giving the displayed bound. The high-derivative integrability
% of H4 T_delta supplies the same or smaller normalized O(L^-2) Poisson
% error. For this fixed algebraic K the uniform derivative-tail bounds
% follow by splitting convolution at |z|=r/2 and putting derivatives on
% the far factor; K_delta's r^-12-j derivative tail has coefficient
% delta^9, bounded on [0,delta0]. No Schwartz-tail assumption is imported.
% Consequently the EXACT averaged second source obeys
%   int H4(y/L) ddot bar rho(y)d^3y
%       =rho_amp L [I_delta+O(L^-2)] >0
% at large L, uniformly over alias-free ell with ell/L<=delta0.
% A radial integrable source with finite fourth moment has zero such
% moment. This proves the stated noninvariance throughout UA11's broader
% physical class, without selecting a special annulus or taking delta->0.
% The physical application still requires a<<ell<<the geometric lengths;
% uniformity on this larger mathematical interval does not remove that
% hierarchy or prove ell independence of the interior or of histories.

% RED AC25--AC27 also treats other positive radial bandlimited profiles
% with differentiated decay exponent p>7 under its fixed Schwartz readout.
% Its weighted summability, positive moment and same-source conditions
% check at that stated domain. Such profiles are not silently substituted
% into our already fixed K or preparation. The overlap proved here is
% our actual phi/K and the exact initial fields, with the explicit distant-
% coherence enlargement. No all-state obstruction, actual prior history,
% new external X, angular projection, or numerical run follows.

% FINAL ACTUAL BROADER-CLASS READ / SOURCES [R].
% Both complete current author and RED sources and every SOURCES block
% were reread in untruncated ranges, including the revised compact cores,
% UA11a--d, AC25--AC28 and the exact saved RED final disposition. The
% full native vertex ranges, arbitrary diagonal-block cancellations,
% weighted CAR ball and finite expected excitation counts agree. The
% fixed positive mass, flat initial P0, specified energy cells/readout
% and absence of off-diagonal entries through distance five are retained;
% distant coherences are permitted, with no later ansatz restriction.
% The two independent angular arguments agree at their stated scopes.
% The global moment has coefficient 192pi gamma/175 and is preserved
% exactly by the actual algebraic kernel's finite-fourth-moment identity;
% its weighted errors are uniformly controlled on the stated alias-free
% ell/L interval. Physical use still requires the hydrodynamic hierarchy.
% Both momentum signs share the initial fields and the zero first source
% derivatives, so the nonradial second derivative excludes a C2 exact
% spherical interval through this characterized class for either sign.
% No general short-coherence obstruction, prior collapse preparation,
% coupled interval, finite-correction ell independence or new external X
% is inferred. The conditional transported-area theorem keeps all its
% actual-history premises. No mandatory scientific repair remains;
% the author source is READY/HELD for the original pair's handoff.

% UA12. RELEVANT FOUR-STEP COHERENCES [R, independent starting save].
% The prior handoff is published. This continuation is derived before
% any new RED or KM/EL/CP intake; the printed core is unchanged. It tests
% a materially different physical class, with off-diagonal excesses at
% distance FOUR which the UA11 second-energy vertex can see. Zero first
% derivatives below are a consequence of this preparation, not an added
% requirement on the full goal or on other candidate radial histories.
% Fix a=1, mu=ma>0, g=G/a^2>0, gamma>0, and the SAME UA5 phi,K and
% point readout. Choose one nonzero real trigonometric polynomial T(k),
% even in each coordinate, cubic invariant and pi/2 periodic in each
% coordinate, before choosing L or testing its consequences. Its site
% kernel is real, finite range and supported on 4Z^3. For example
% T=1+kappa sum_i cos(4k_i), kappa!=0, supplies actual distance-four
% coherences; no small-kappa restriction is needed beyond the CAR bound
% below. Higher fixed polynomials are also allowed. T is spin scalar.

% UA12a. Exact physical preparation and fixed source cells. Define
% W(q)=sum_z |T_z|^2 exp(-iq.z), W0=W(0)=<T^2>>0.
% W is real, even, cubic, and positive near q=0. Choose one real even
% smooth torus cutoff chi equal to one in a ball about zero and supported
% inside a ball on which W>0, and set Ghat=chi W0/W. Then G is real,
% even, has every weighted lattice moment and sum G=1. With
% F_L=G*phi(. /L), the relative second-difference argument of UA5 gives
% |F_L-phi(. /L)|<=C_T L^-2 phi(. /L).
% For sufficiently large L, F_L is positive, bounded by 3/2, summable,
% and W*F_L=W0 phi(. /L) EXACTLY on the target's low Fourier support.
% Differentiating the same smooth interpolation gives
% |grad log F_L|<=C_T/L. These are proved properties of one initial
% envelope, not envelopes refitted in time or across ell.
%
% Put A_L=T F_L T, nu_L=gamma/(4mu g W0 L^2), and
%                  C_L=P0+nu_L beta A_L.
% This is CAR-admissible for large L. Indeed S=h0/E commutes with T,
% beta commutes with T,F_L, and the massive weighted Schur argument gives
% {S,beta F_L}>=s_* F_L, s_*=mu/sqrt(mu^2+3)>0, at sufficiently large L.
% Hence {S,beta A_L}>=s_* A_L. Also
% A_L^2<=||T||^2 ||F_L||_infinity A_L. Therefore
% C_L(1-C_L)=(nu_L/2){S,beta A_L}-nu_L^2 A_L^2>=0
% whenever nu_L ||T||^2 ||F_L||_infinity<=s_*/2. Spectral calculus then
% gives 0<=C_L<=1. This proof needs no inverse of T and permits its zeros.
% A_L is positive trace class, ||C_L-P0||_1=4nu_L W0 sum phi_L<infinity.
% Expected particle and hole numbers are equal, with
% N_ex=4nu_L mu<T^2/E> sum phi_L; net charge is zero. These are Gaussian
% expectations, not a pure-state or fixed-number assertion. In particular
% E_tot=4mu nu_L W0 sum phi_L=gamma L(int phi)/g.
%
% Every A_L matrix element has x-y in 4Z^3. Exact native spin traces give
% rho_x=4mu nu_L (A_L)_xx=rho_amp phi(x/L), rho_amp=gamma/(gL^2),
% J_x=dot J_x=0. The pure-h0 energy commutator is spin scalar, while the
% only surviving spin trace of the flow energy commutator pairs A_L
% with a two-step bond. That entry vanishes, so dot rho_x=0 as well.
% Thus the fixed readout yields exactly the same positive radial initial
% rho,J and solved point fields N=1,v=V_delta(x/L) as UA5c, for BOTH
% momentum signs. The finite expected-count and incoming-P0 qualifications
% are the same as UA10; no preceding collapse or radial interval is supplied.

% UA12b. Exact second derivative retains the newly relevant coherences.
% On any hypothetical regular C2 spherical interval the zero first
% sources force hdot=0 initially. Clifford traces in D=nu beta A leave
% only the two-flow contribution, for every scalar site matrix A:
%       ddot rho_x=-4mu nu [V,[V,A]]_xx.
% For a fixed axis put t_i(x)=[v_i(x)+v_i(x+2e_i)]/8, so
% V_(x,x+2e_i)=-i t_i(x). Direct multiplication, using A support in 4Z^3,
% gives
% ddot rho_x/(4mu nu)=2 sum_i {
%  t_i(x)^2(A_(x+2e_i,x+2e_i)-A_xx)
% +t_i(x-2e_i)^2(A_(x-2e_i,x-2e_i)-A_xx)
% +t_i(x)t_i(x+2e_i) A_(x+4e_i,x)
% +t_i(x-2e_i)t_i(x-4e_i) A_(x-4e_i,x)
% -2t_i(x)t_i(x-2e_i) A_(x+2e_i,x-2e_i) }.
% Unlike UA11's stencil, this expression contains physical distance-four
% coherences. It uses common Ddot=-i[h,D] and the actual point fields.
%
% Write b4=<T^2 cos(4k_1)>/W0 and
% c=(1+b4)/2=<T^2 cos^2(2k_1)>/W0. Cubic symmetry makes c independent
% of the chosen axis; 0<c<1 for every nonzero polynomial T. The exact
% midpoint kernel of T F_L T has symbol
% T(k+q/2)T(k-q/2) Fhat_L(q). Its even q expansion and the fixed finite
% kernel moments therefore give, uniformly on scaled compact sets,
% ddot rho_x=(rho_amp/L^2) U_delta(x/L)+O(rho_amp/L^4),
% U_delta=c sum_i partial_i(V_delta,i^2 partial_i phi)
%             +(2c-1)sum_i partial_i(V_delta,i phi partial_i V_delta,i).
% One can check the coefficient without symbol calculus: the diagonal
% terms above give (1/2)partial_i(v_i^2 phi_i), and the distance-four
% terms give (b4/2)partial_i^2(v_i^2 phi). Their sum is the displayed U.
% The order-L^-2 correction to the midpoint envelope cancels out of
% this leading second difference. At T=1, c=1/2, the formula reduces
% exactly to UA10's leading coefficient. No continuum advection rule
% was substituted for the native two-step coupling.

% UA12c. The relevant coherences cannot remove this family's full angular
% derivative. First take delta=ell/L to zero in the smooth initial fields.
% Put V_i=w(r)X_i, w^2=2gamma m(r)/r^3,
% m(r)=4pi int_0^r s^2 phi(s)ds. The coefficient of
% H4(X)=sum_i X_i^4-3r^4/5 in U_0 is exactly
%  (1/r)d/dr { [c w^2 phi'+(2c-1)w w' phi]/r }.
% UA5's fixed profile has phi=F_infty r^-12 plus a remainder decreasing
% with all derivatives faster than every additional power, F_infty>0.
% Thus m=M_infty-d r^-9 plus such a remainder, d=4pi F_infty/9>0.
% Substitution, using w w'=(w^2)'/2, gives this harmonic coefficient as
%  51gamma M_infty F_infty(10c-1) r^-19
%       +gamma d F_infty(312-1248c) r^-28
%       +a remainder decreasing faster than every displayed power.
% If c!=1/10 the first coefficient is nonzero. If c=1/10 it cancels,
% but the second is (936/5)gamma d F_infty>0. Consequently EVERY fixed
% nonzero T in this class has a finite scaled annulus with a nonzero H4
% component, even the occupation moment tuned to cancel the leading tail.
%
% Fix that annulus first. For ell=L^theta, 0<theta<1, the actual radial
% fields V_delta tend to V_0 smoothly there. The SAME positive radial
% K_delta tends to the identity and preserves angular harmonic degree.
% The finite T moments and the smooth G envelope control the Taylor
% remainder and point-sampling aliases as in UA10. Hence the exact
% averaged ddot rho is nonradial at all sufficiently large L, uniformly
% for theta in any fixed compact subinterval of (0,1). Both signs of
% the momentum row have the same initial fields and obstruction.
% The argument also covers any sequence ell(L)->infinity with ell/L->0:
% fix the nonzero annulus, then choose delta sufficiently small and L
% sufficiently large. Uniform integrable derivatives of the sampled
% profiles give each noncentral Fourier copy O_J(L^-J); multiplying
% by the fixed cutoff Khat(ell q) and integrating gives an alias bound
% O_J(L^(3-J)ell^-3) for the normalized smooth profile. J can be chosen
% arbitrarily large. Thus no power-law rate for ell is needed for this
% initial-derivative conclusion. The same estimate applies to the
% actual AC29--AC32 conversion below, including its cubic first jet.
% The
% physical hierarchy also retains ell<<every geometric length; fixed
% gamma changes those constants but not the large-L separation.
% This rules out a C2 exact spherical interval for this larger fixed
% preparation class, not arbitrary short coherences or states with
% nonzero radial first derivatives. No covariance is changed in time.
% It does not solve the full unprojected-history or approximation task.
% All quantities use fixed mu,g,gamma,T and L=L_phys/a,ell=ell_phys/a.
%
% SOURCES / INDEPENDENCE FOR UA12. The binding goal was reread in full.
% The actual UA5 profile/readout and UA10 massive form estimate were
% reread before the new construction. The weighted T preparation,
% exact five-term two-flow contraction, leading c tensor and two-tail
% obstruction were derived here before any new RED continuation or
% KM/EL/CP intake. No code, numerical run, source projection or new X.

% UA12d. A necessary-profile equation explains why tail tuning is not
% invariance. This further independent step uses no new counterpart.
% Replace phi only for this conditional extension by a fixed strictly
% positive radial bandlimited f comparable to (1+r)^(-p), p>7, with
% bounded logarithmic gradient and the corresponding differentiated
% decay bounds. Require the same
% low-band envelope/CAR and actual readout/error proofs; they hold by
% the displayed estimates for such a profile and the fixed UA5 K.
% If the leading tensor U_0 were radial everywhere, its H4 coefficient
% would vanish. Decay then sets the integration constant to zero:
%       c w^2 f'+(2c-1)w w' f=0.
% Since c>0, f>0 and w^2>0, integration using log|w| gives
%       f=A (m/r^3)^n, n=(1-2c)/(2c), A>0,
% using w^2=2gamma m/r^3 and absorbing its constant into A.
% For n<=0 the right side cannot tend to zero at finite total mass.
% For 0<n<1, integration of m'=4pi r^2 f gives
%       m^(1-n)=(4pi A/3)r^(3(1-n))+C,
% inconsistent with finite positive m(infinity). At n=1, m is a
% positive power of r and also has no finite nonzero limit. Thus only
% n>1 could work, and the regular finite-mass solution must be
%       f(r)=A[(4pi A/3)+C r^(3(n-1))]^(-n/(n-1)), C>0.
%
% Band limitation excludes this remaining family. Because f is L1
% with compact Fourier support, its restriction to a complex coordinate
% axis is entire (the Fourier integral converges on every compact
% complex set). Real analyticity and radial smoothness at zero force
% 3(n-1)=2j for a positive integer j in the displayed nonconstant
% solution; otherwise its first nonconstant power is not an even
% integer. But for 2j the denominator has finite complex zeros, and
% the positive exponent n/(n-1) gives a pole or branch singularity,
% contradicting that entire extension. The constant solution has
% infinite mass and is already excluded. Hence no profile in the
% specified positive bandlimited finite-mass class makes this tensor
% radial for any fixed nonzero T. This supplies a necessary invariance
% obstruction, not a permission to refit f as a function of time.
% Passing its nonzero finite-annulus witness to the actual native
% derivative retains the fixed-profile, large-L and ell/L->0 controls
% of UA12c. The original fixed phi/K proof above does not depend on
% this extension or on a uniform bound over all T or all profiles.

% ACTUAL AC29--AC32 COMPARISON AFTER BOTH INDEPENDENT SAVES [R].
% Read the complete new RED block/SOURCES only after UA12a--d and its
% independent distance-two construction were saved. Its A_b=I+sum_i
% b_i(U_i^2+U_i^-2), sum|b_i|<=1/4, is a DIFFERENT fixed preparation
% from UA12's pi/2-periodic T. Its exact W_b, even inverse, massive
% weighted form bound, equal finite particle/hole counts and positive
% initial energy cells all check. Its flow-energy contraction is
% dot rho_x=-mu sum_i b_i[(v_x+v_(x+2i))(F_x+F_(x+2i))
%                         -(v_x+v_(x-2i))(F_x+F_(x-2i))],
% whose two-transfer symbol is exactly
% -2i sum_i b_i sin(q_i) cos(l_i) cos(p_i) vhat_i(l)rhohat(p)/W_b(p),
% q=p+l. Thus nonzero first derivatives are actually tested, not
% excluded by a condition imposed on candidate histories.
%
% Unequal b_i give the stated leading degree-two angular tensor. For
% equal nonzero b, expansion of the exact symbol gives the three
% negative cubic derivatives with coefficients 1/6,1/2,1/2 and the
% radial c_iso term. The independently recomputed tail brace is
% alpha(alpha-2)/2+beta(beta-2)/2+(alpha+beta)(alpha+beta-2)/6.
% At its actual alpha=-8,beta=-3/2 this is 365/6, as saved. If the
% SAME construction is instead explicitly made with our already fixed
% phi/K, alpha=-12 gives 243/2, also nonzero. Its envelope/CAR and
% exact-source proof extends to this fixed target by the same weighted
% inverse estimates; its first-tangent implication uses its own actual
% source cells and solved point flow. No tail coefficient or readout is
% identified across the two distinct preparations without that conversion.
%
% RED's additional Hermitian trace-class perturbations at separation>=4
% cannot change initial rho,J or dot rho, whose exact vertex ranges are
% <=1,2,3. Its weighted CAR ball proves nonemptiness with those allowed
% perturbations. Dot J may change, which is harmless for the already
% nonradial first-energy obstruction. No later band restriction is used.
% This gives a complementary first-tangent exclusion for the displayed
% distance-two family; UA12 supplies a second-tangent exclusion when
% four-step coherences are the relevant entries. Both retain fixed
% mu>0, g,gamma>0, finite expected count, common transported initial P0,
% both momentum signs and the physical ell hierarchy. Neither is an
% all-state theorem or a construction of an unprojected coupled interval.

% FINAL ACTUAL SHORT-COHERENCE READ / SOURCES [R].
% Both complete current author and RED files and every SOURCES block
% were reread in untruncated ranges, including the exact revised cores,
% UA12a--d, AC29--AC34 and the saved RED final disposition. The binding
% goal was reread in full. The independent substantive saves preceded
% counterpart intake. No mandatory scientific repair remains.
% The massive weighted CAR forms, finite expected particle/hole counts,
% fixed diagonal source cells and exact native commutators agree. The
% two-step family has a nonradial first energy derivative, including its
% cubic term when the three coefficients agree. The four-step family
% has a nonradial second energy derivative, including the occupation
% moment that cancels its leading far-tail coefficient. The necessary
% profile equation and its entire-function exclusion have the repaired
% signed-flow domain f>0,w^2>0 and use log|w|. The distinct r^-8/r^-12
% targets retain their own coefficients and averaging kernels.
% The fixed-annulus Taylor and Poisson estimates in AC34 agree with
% UA12c: noncentral aliases are O_J(L^(3-J)ell^-3), so a fixed sufficiently
% large J controls each derivative order for every ell->infinity,
% ell/L->0 sequence within the physical hierarchy. No finite correction,
% endpoint scale or actual history is thereby made ell independent.
% The C1/C2 exclusions concern the respective fixed physical initial
% classes, under common incoming-reference transport and actual point
% coupling, for both momentum-row signs. Nonzero radial derivatives are
% allowed; arbitrary relevant coherences, prior incoming preparations and
% exact or controlled approximate coupled intervals remain open. No new
% external X, source projection, field, refit, code or run is introduced.
% The conditional transported-area theorem retains every history premise.
% The author source is READY/HELD for the original pair's handoff.

% UA13. AN ARBITRARY-COHERENCE LOCALIZATION OBSTRUCTION [R, independent save].
% This continuation follows the published short-coherence handoff and is
% saved before any new RED or optional R/KM/EL/CP intake. It keeps the
% printed core unchanged. It imposes a condition on the ACTUAL energy
% sequence, not a finite-range covariance ansatz or zero source jets.
% No scientific code or run is used.
%
% UA13a. Fix the infinite parent, finite mu=ma>=0 and g=G/a^2>0.
% Set a=1 in this proof. Let C be ANY physical covariance 0<=C<=I,
% D=C-P0 trace class and nonzero, with P0 the flat incoming sea. Its
% spin, particle-hole and arbitrarily near or distant coherences are
% unrestricted. Define e_x=Tr[{Pi_x,h0}D/2], and assume only
%            sum_x exp(sigma |x-x_c|)|e_x|<infinity
% for some sigma>0 and prospective continuum centre x_c. The source
% may have either sign at individual sites. It is the physical native
% energy sequence, not an independently chosen positive density.
% Choose ANY fixed normalized radial L1 point kernel K, including the
% actual noninjective UA5 K, and K_ell(y)=ell^-3 K(y/ell), ell finite.
% Positivity of K is permitted but is not needed for this exclusion.
%
% Suppose bar rho(y)=sum_x e_x K_ell(y-x) is exactly radial about x_c.
% The finite measure nu=sum_x e_x delta_(x-x_c) has a Fourier transform
% holomorphic in the tube |Im k|<sigma, by the displayed exponential
% moment. Khat_ell(0)=1 and continuity supply an open ball on which it
% is nonzero, even when Khat has zeros elsewhere. For each rotation O,
% radiality gives nuhat(k)=nuhat(Ok) on that ball. Both sides are real
% analytic on the connected real momentum space, so the identity theorem
% extends the equality to every k. Fourier uniqueness makes nu rotation
% invariant. A finite atomic measure on the shifted cubic lattice can
% then have support only at zero: a nonzero atom away from the centre
% would require its entire uncountable rotation orbit. Thus e_x is zero
% except possibly at the single central site, if x_c is a lattice site.
%
% The native checkerboard/CAR identity now contradicts physicality. With
% S_x=(-1)^(x1+x2+x3), R=S beta, one has {S,h0}=2mu R and [R,h0]=0.
% Put G=P_+ C P_+>=0, H=P0(I-C)P0>=0 and N_ex=Tr(G+H). Then
% E_D=Tr(h0D)=Tr[E(G+H)]>mu N_ex>=|mu Tr(RD)|
% for every nontrivial D. Strictness holds because E-mu has no nonzero
% l2 kernel on the infinite sine parent; its zeros are isolated momentum
% points. If N_ex=0, positivity of C and I-C removes the interband block
% and gives C=P0. But central support would require
% S_(x_c) E_D=sum_x S_x e_x=mu Tr(RD), a contradiction. A centre not on
% the lattice permits no atom and instead gives E_D=0, also impossible.
% Thus this class admits NO nontrivial exactly radial initial energy
% source for ANY normalized radial L1 point averaging kernel. This
% covers arbitrary relevant covariance coherences and every finite mass,
% not just the previously tested scalar or polynomial preparations.
%
% For the cell-box allocation, the same exponential moment makes the
% Fourier transform of f_D=sum_x e_x 1_(cell x) analytic in a tube.
% The same low-frequency argument makes f_D itself radial. The actual
% cell-face rotation proof in UA1 then makes this integrable step
% function zero, again contradicting E_D>0. The kernel need not be
% injective in either version. The source map is translation convolution,
% with no angular projection or additional geometric equation.
%
% UA13b. Nonemptiness, scales and the allowed escape are explicit.
% The exponential-source class is nonempty at EVERY finite mu>=0:
% C=P0+lambda P_+ F P_+, 0<lambda<=1, with any nonzero finite-support
% real 0<=F<=1, is physical and trace class. Its exact positive energy
% sequence is the nearest-neighbour W_mu*F proved in UA5, hence finite
% support even at mu=0. This example proves nonemptiness only; UA13a
% covers all physical covariances with the stated source moment.
% The exclusion holds for every finite ell/a and g and hence throughout
% the declared hierarchy. It is exact, not a uniform lower bound on
% angular error as the source length grows. Smooth Gaussian or other
% localized envelopes can have arbitrarily small hydrodynamic angular
% errors without exact equality. The previously constructed noninjective
% radial data have algebraic tails and fail every positive exponential
% source moment; they do not contradict this result. More generally a
% possible escape must lose this analyticity condition, but need not
% have a power-law tail: nonanalytic subexponential tails are also open.
% This is an initial flat-P0 theorem, not the false assertion that the
% free-energy inequality holds relative to every later transported sea.
%
% UA13c. A necessary interval structure for localized reached sources.
% Independently of the flat-P0 inequality, suppose an actual common-
% transport interval has trace-class D(t), exponential moments for BOTH
% native e_x(t) and J_xi(t), and exactly spherical averaged sources at
% each time under the same fixed radial kernel. Apply UA13a's analytic
% argument to scalar and vector Fourier transforms. The native energy
% measure must be E(t)delta_(x_c), and the native vector measure must
% vanish: it has no noncentral rotation orbit, while a vector at the
% centre invariant under all rotations is zero. These are consequences
% of exact spherical closure, not extra zero-derivative requirements.
% For the actual smooth strictly positive UA5 K and E(t)>0, the regular
% Gauss rows give N=1 and v from M(r)=E(t)4pi int_0^r s^2 K_ell(s)ds,
% for either momentum sign. Hence h=h0+V with V=sum_i{v_i,Q_i}/2.
% Since the current vertices are fixed, the EXACT point-pairing identity
% gives dE/dt=-sum_(x,i) v_i(x,t) dJ_xi/dt=0. Indeed
% sum v_i dot J_xi=i Tr([h,V]D)=i Tr([h0,V]D),
% whereas dot E=i Tr([V,h0]D). The bounded coefficients and trace-class
% D justify the sums. Thus E,N,v must be time independent on any such
% positive-mass interval. This does not require the covariance itself
% to be stationary, and does not prove that such reached data exist.
% At a flat-incoming initial slice UA13a excludes them altogether.
% At a later P_T the missing free-energy inequality is not imported;
% existence of a central-source stationary-geometry state remains a
% separate problem. No arbitrary lapse choice is used at positive mass.
%
% A sufficient propagation domain for these source moments is
% ||W_sigma D(0)W_sigma||_1<infinity, W_sigma(x)=exp(sigma|x|/2).
% For bounded finite-range h(t) on a fixed finite time interval,
% W_sigma h W_sigma^-1-h is bounded by the finite hop range and the
% coefficient bound. The weighted propagator and its inverse therefore
% have finite exponential-in-time norm bounds. Factoring the common
% D(t)=U D(0)U^dag proves ||W_sigma D(t)W_sigma||_1<infinity; bounded
% finite-range energy/current vertices then give their exponential
% moments. This sufficient domain is not necessary and is not asserted
% for every massless preparation. The source-only initial theorem above
% already has a nonempty massless domain. No uniform growing-time or
% ell-independent propagation estimate is inferred from this finite-time
% statement.
%
% SOURCES / INDEPENDENCE FOR UA13: binding GOAL_alt.md reread in
% full; actual native checkerboard/CAR proof UA1a, source allocation UA5,
% and exact common-transport point-pairing identities in this source.
% The analytic-tube continuation, arbitrary-coherence scope and necessary
% interval structure were derived and saved before any new RED route
% or optional counterpart intake. No external X is assigned. The
% nonanalytic source class, actual unprojected coupled existence and
% controlled approximate-history/ell tests remain unresolved.

% UA13d. Actual original-pair comparison and stronger reached conditions [R].
% Read the complete independent RED AC35--AC38 and SOURCES only after
% UA13a--c and that route were substantively saved. Its analytic-tube,
% low-frequency cancellation and native CAR/checkerboard proof agrees
% with UA13a for arbitrary coherences and a SOURCE-only moment condition.
% Its weighted trace-class propagation proof agrees after converting
% W_kappa=exp(kappa|x|) to our W_sigma=exp(sigma|x|/2), sigma=2kappa.
% Our cell-box and shifted-centre arguments and the massless nonempty
% source-only example add their stated domains; no weighted-covariance
% hypothesis is silently attached to the broader initial theorem.
%
% RED's additional conserved-charge route was checked directly. R=S beta
% commutes with every site multiplier, with h0, and with every Q_i because
% its hops have even parity and Q_i is spin scalar. Thus [R,h[N,v]]=0
% for the ACTUAL endpoint-ordered Hamiltonian, including nonconstant N.
% Common transport conserves Tr(RD). At a central native energy source,
% S_(x_c) E(t)=mu Tr(RD_in), so E is constant without first assuming J=0
% or using a strict free-energy inequality at the later reference.
% In particular a massless reached analytic radial point source has
% E=0, independent of the earlier history. With analytic vector sources
% as well, J=0; at positive E and the actual strictly positive UA5 K,
% the regular continuous inward branch has the static N=1,v of UA13c.
% This agrees with the independently derived exact point-pairing proof
% of E'=0. It is not an arbitrary-lapse construction at zero mass charge.
%
% RED AC37's further stationary spectral prerequisite also checks. Fix
% a time inside a proposed positive interval, call its trace-class excess
% D_*, and use the resulting STATIC bounded native h. For every local
% finite-rank source vertex b, Tr[b exp(-iht)D_* exp(iht)] extends to an
% entire function of complex t, with nth derivative bounded on the real
% axis by ||b|| ||D_*||_1(2||h||)^n. Its constant source value on an open
% interval therefore extends to all real times for this autonomous
% comparator. No earlier physical history is changed by this analytic
% continuation, and its transported reference is not replaced.
%
% On Hilbert--Schmidt operators the bounded self-adjoint commutator
% L_h:A->hA-Ah has a spectral theorem. The Cesaro averages of
% exp(-it L_h)D_* converge in that norm to D_av in ker L_h, since their
% scalar multiplier is bounded by one and tends to zero off zero.
% Their trace norms are <=||D_*||_1. Trace-norm lower semicontinuity
% under Hilbert--Schmidt convergence gives D_av trace class with the
% same bound. Pairing with every local source vertex passes to the limit,
% so D_av has the same central energy E and zero native current.
% It is self-adjoint and compact. Each nonzero finite-dimensional
% eigenspace of D_av is h-invariant; diagonalizing h there places the
% range of D_av in h's point-spectral subspace. Consequently a positive
% interval in this localized class REQUIRES a nonzero stationary
% point-spectral source with those exact native values. If this actual
% h has no point spectrum, such an interval is impossible. No absence
% of gravitational bound states has been proved or assumed here.
% D_av is a necessary source test, not a newly prepared covariance
% against P0 or an instantaneous replacement of the transported sea.
% The actual native eigenstates, their physical reference compatibility
% and the central-source identities still need a construction or proof.
%
% Both routes therefore constrain arbitrary relevant coherences under
% explicit localization premises, for every finite mu,g and allowed ell.
% They retain the nonanalytic-tail escape already realized by actual
% initial data. Finite expected count alone implies none of these moment
% conditions. The source width ell of a central smoothed charge is not
% silently declared an allowed geometric hierarchy or a new exclusion;
% actual geometric derivative scales still have to be checked. No
% external X, additional state, source projection, code or run follows.

% FINAL ACTUAL SOURCE-REGULARITY READ / SOURCES [R].
% Both complete current author and RED sources and every SOURCES block
% were read in untruncated ranges, including UA13a--d, AC35--AC39,
% both compact statements and the exact reached-source wording repair.
% The independent substantive saves preceded actual counterpart intake.
% The low-frequency analytic-continuation and native checkerboard/CAR
% proofs agree for arbitrary initial covariance coherences, finite
% expected excitation count and every finite mu>=0. The exponential
% premise is on the actual energy SOURCE; weighted covariance decay is
% only a separately stated sufficient propagation domain. No such
% moment follows from finite count alone. Point and cell-box allocation,
% noninjective kernels, shifted centres and the nonempty massless domain
% retain their exact specified scopes.
%
% For exactly spherical reached point averages, analytic energy/current
% sources are central/zero, and conserved R=S beta fixes their energy,
% which vanishes at mu=0. On the specified positive-kernel positive-mass
% regular branch, the Gauss fields are static for both momentum signs.
% The independent native point-pairing proof agrees. Bounded-h entire
% traces and the Hilbert--Schmidt Cesaro argument give only a necessary
% stationary point-spectral source; they do not assert absence of bound
% states, replace P_T by P0, or prepare a new physical covariance.
% Nonanalytic source tails, actual unprojected coupled existence and
% controlled angular approximations remain open. The conditional area
% theorem retains its incoming-state, space/time Fourier, clock, fixed-T
% and allowed-ell premises. No new external X, projection, time refit,
% scientific run or additional geometric law is introduced.

% UA14. THE SPHERICAL SOURCE FIBRE IS NOT A NATIVE INVARIANT SECTOR [R].
% Independent substantive continuation after the UA13 handoff. Saved
% before reading RED AC40--AC42 or a new counterpart route. This is a
% statement about arbitrary covariances in an explicit physical order
% domain, not another fixed occupation ansatz. The printed core is left
% unchanged until the original-pair comparison. No code or run is used.
%
% UA14a. Fix the actual infinite parent and the same point readout,
% any finite mu=ma>=0, g=G/a^2>0 and one finite allowed ell/a. Use a=1
% and tau=t/a below. At a specified slice let R be the ACTUAL transported
% incoming projector, C a physical covariance, and D=C-R trace class.
% Suppose its actual averaged sources are exactly spherical and admit
% a regular real point-coupled N,v for the chosen printed momentum sign.
% No first source derivative is assumed zero. Require a finite-site
% ORDER margin: for a projection Pi_B onto a finite set of complete
% native spin blocks, C>=eta Pi_B and I-C>=eta Pi_B, eta>0. This is an
% explicit global operator-form premise, stronger than faithfulness of
% the reduced state on B; the latter alone would not permit the following
% perturbation. All other short/long, spin and band coherences are free.
%
% Choose distinct axes i,j and a site r. Put u=r+e_j, t=r+2e_i,
% s=r+2e_i+e_j, and choose this four-site set away from the spherical
% centre. Only r,s need belong to B. Define the Hermitian finite-rank
% perturbation
%       Z_r=i alpha_j (|r><s|-|s><r|),       ||Z_r||=1.
% Its only off-diagonal separation has Manhattan length three. Hence
% Tr(e_x Z_r)=Tr(j_xk Z_r)=0 at EVERY site, since the actual fixed
% vertices e_x={Pi_x,h0}/2 and j_xk={Pi_x,Q_k}/2 have ranges one and
% two. It also has Tr Z_r=0. Thus C+epsilon Z_r has exactly the SAME
% initial native and averaged sources, and the same N,v selected by
% the declared source functional. For |epsilon|<=eta it remains CAR:
% -Pi_B<=Z_r<=Pi_B and the two displayed order inequalities apply.
% Its difference from R is trace class. The reference is not changed.
%
% The two common-transport derivatives differ exactly by
%       delta dot e_x=epsilon i Tr([h,e_x]Z_r).
% There is no spectral-reference derivative and no response of h to
% this initial perturbation, because every initial source is unchanged.
% The lapse part {N,h0}/2 has range one, so its commutator with e_x
% has range at most two and contributes zero, even for varying N and
% nonzero mass. The sole contribution is the actual point flow V.
% Define A_i(z)=[v_i(z)+v_i(z+2e_i)]/8, so
% V_(z,z+2e_i)=-i A_i(z), h0_(z,z+e_j)=-i alpha_j/2.
% Direct multiplication, retaining BOTH possible paths, gives
% [V,e_x]_(r,s)=(alpha_j/4)
%       [A_i(u)(delta_(x,r)+delta_(x,u))
%                     -A_i(r)(delta_(x,t)+delta_(x,s))].
% Therefore, with the full4 spin trace,
%       i Tr([h,e_x]Z_r)=2[A_i(u)(delta_(x,r)+delta_(x,u))
%                         -A_i(r)(delta_(x,t)+delta_(x,s))].
% This exact four-site source is nonzero whenever either A_i(r) or
% A_i(u) is nonzero. Neither its sign nor its existence depends on mu.
% In a nonzero regular inward spherical flow, suitable off-centre bonds
% exist, for example in any positive-mass exterior region with x_i>0.
% The proof uses the actual sampled v for either momentum-row sign;
% it does not replace it by an affine or continuum matter generator.
%
% UA14b. Exact count preservation and the averaging obstruction.
% Take two translated copies of the preceding bond, with pairwise
% disjoint four-site output sets away from the centre and nonzero flow
% coefficients. Enlarge B to contain their four input sites, retaining
% the same explicitly required order margin. With q_b=Tr[(I-2R)Z_b],
% choose a nonzero real vector (c_1,c_2) satisfying sum_b c_b q_b=0 and normalize
% max_b|c_b|=1. If both q_b vanish any nonzero choice works. Then
% Z=sum_b c_b Z_b has ||Z||=1, Tr Z=0 and Tr[(I-2R)Z]=0.
% The expected particle and hole counts relative to the SAME R are
% individually unchanged, because their variations are Tr[(I-R)Z]
% and -Tr(RZ), both zero. The native total free energy and every native
% source are also unchanged initially. For R=P0 the translated bonds
% have identical q_b, so the simple difference Z_1-Z_2 suffices.
% For a reached R the displayed two-bond nullspace condition
% uses that actual projector; no replacement by P0 is made.
%
% Disjoint output supports make the derivative map injective on this
% two-bond span. Thus delta dot e is a nonzero finite atomic measure
% supported away from the spherical centre. Under ANY normalized radial
% L1 convolution kernel its average cannot be radial: the multiplier
% is nonzero near zero, the finite-source transform is entire, and
% radiality would analytically extend to a rotation-invariant atomic
% measure supported only at the centre. There is no central atom here.
% The same conclusion holds with the specified cell-box allocation by
% the actual radial-step-function proof. Noninjective filters do not
% remove this finite-support witness. No lower bound uniform as ell
% grows is asserted; this is an EXACT finite-spacing statement.
%
% Consequently, if the base state's actual first energy derivative is
% radial, every nonzero sufficiently small perturbation epsilon Z fails
% that necessary tangency. More generally at most one member of the
% physical affine line C+epsilon Z can have a radial energy derivative,
% since the difference between any two derivatives is nonradial.
% No condition was imposed on the value of a radial first derivative;
% nonzero radial evolution is allowed. All members have the same Gauss
% initial data and excitation counts. The source functional therefore
% cannot make the ENTIRE physical same-source fibre an invariant
% spherical sector on this order domain. An invariant subset selected
% by additional coherence conditions is not excluded or constructed.
% The perturbation is fixed at the initial slice, not refitted in time.
% At a reached slice this tests physical alternative data with that same
% reference; it does not assert that the alternative has the same prior
% self-consistent incoming history.
%
% UA14c. This physical order domain is nonempty at EVERY finite mass.
% Unlike the beta-F ansatz, the following construction includes mu=0.
% Write P_+=I-P0 and choose a real scalar diagonal 0<F<=I/2, summable.
% Define one physical mixed Gaussian covariance
%       C_*=P0+P_+ F P_+-P0 F P0,
%       D_*=P_+ F P_+-P0 F P0={S,F}/2,
%       S=I-2P0=h0/E.
% Its particle/hole blocks are separate contractions, so CAR is exact.
% Moreover C_* and I-C_* both dominate the pinching
% P_+FP_++P0FP0=(F+SFS)/2>=F/2. Hence every finite set has the actual
% order margin required above. This argument needs neither a mass gap
% nor an exponentially decaying P0 kernel. Its equal finite expected
% particle and hole counts are 2 sum_x F_x each, and ||D_*||_1<=
% 4 sum_x F_x. It is not a fixed-number or pure-state preparation.
%
% The exact native source calculation gives
%       e_x(D_*)=2(W_mu F)_x,          J_x(D_*)=0,
% where W_mu is the same positive nearest-neighbour energy stencil
% established in UA5 for P_+ F P_+. This also follows directly from
% D_*={S,F}/2: h0 S=E, and the full spin trace of S at every separation
% is zero. The negative-band hole contribution equals the positive-band
% energy contribution; no species or valley multiplier is added.
%
% To exhibit exact spherical initial data with the FIXED noninjective
% UA5 kernel, retain its positive bandlimited phi(x/L),
% w0=What_mu(0)=2<E>, and the real even low-band inverse G_mu with
% Ghat_mu=chi/What_mu; What_mu denotes the stencil's Fourier symbol. Set
%       F=eta_L w0 G_mu*phi(. /L),
%       eta_L=gamma/(2g w0 L^2),    gamma>0 fixed.
% The already proved weighted second-difference estimate, including
% mu=0, gives eta_L phi/2<=F<=3eta_L phi/2. At sufficiently large L,
% F<=1/2, and exactly 2W_mu F=gamma phi(. /L)/(gL^2).
% Poisson alias exclusion then gives the same positive EXACT radial
% averaged energy and J=0 for every allowed alias-free ell, from one
% fixed physical covariance. The initial regular Gauss fields have
% N=1 and nonzero inward v for either momentum sign. This supplies the
% order domain and initial source fibre only: C_* itself is not claimed
% to have radial source derivatives or an earlier collapse history.
%
% Restoring units uses tau=t/a, mu=ma, g=G/a^2, L=L_phys/a and
% ell=ell_phys/a. The displayed energy cell is a times its physical
% energy, and its tau derivative is a^2 times the physical derivative.
% No additional dimensional coefficient enters the four-site identity.
% Fixed g, gamma and L determine the counts and the actual point flow;
% the hierarchy a<<ell<<geometric lengths is retained. Pure covariance
% projectors do not satisfy the two-sided order-margin premise, and
% no theorem about all such pure states is inferred. A special physical
% nonanalytic-source invariant class, actual coupled existence, uniform
% time/clock/ell estimates and charge-loss/state accounting remain open.
%
% UA14d. Arbitrarily many independent tangency constraints.
% Suppose the same covariance has the two-sided GLOBAL operator margin
% on every finite set of sites (with a set-dependent positive eta), and
% its nonzero spherical flow permits arbitrarily many disjoint witnesses
% above. The explicit C_* satisfies the order premise, and its positive
% finite-total-energy radial flow supplies the bonds. Choose any n such bonds,
% with disjoint four-site outputs away from the centre. On their span
% Z(c)=sum_(b=1)^n c_b Z_b, every native initial source and Tr Z vanish.
% Fixing the separate particle and hole counts adds only the one real
% equation sum c_b Tr[(I-2R)Z_b]=0, leaving dimension at least n-1.
% A sufficiently small ball in this finite-dimensional subspace is CAR
% by the actual order margin on its finite union of input sites.
%
% The energy-derivative map on this count-preserving subspace is injective
% even MODULO radial averaged functions. Indeed any linear combination
% whose average is radial is a finite atomic source with radial average;
% the analytic argument forces central support, while disjoint off-centre
% output supports force each nonzero-flow coefficient to vanish. The
% cell-box conclusion follows from the same step-function argument.
% Consequently the first radial energy-tangency condition imposes at
% least n-1 independent real linear conditions on this physical,
% count-preserving same-source chart, for every n. At a fixed source
% fibre h is fixed, so this is an affine-linear condition on C, not a
% linearization that dropped field-response or reference terms.
%
% In particular the tangent-admissible part of this fibre has empty
% relative interior in these physical finite-rank directions. No finite
% collection of additional linear coherence moments can make the entire
% remaining physical fibre invariant: choose n larger than their number
% plus one and impose those moments together with the count equation;
% a nonzero derivative witness remains. This last statement is expressly
% about finite LINEAR moment closures, not every nonlinear or infinite
% set of coherence conditions. It supplies a necessary consistency
% hierarchy, not an exclusion of a specially chosen invariant state set
% or a construction of a coupled interval. No zero radial derivative
% was required and no history was refitted. The result is exact at each
% finite a,ell; it gives no lower bound uniform in a/ell.
%
% SOURCES / INDEPENDENCE FOR UA14: actual native endpoint h,e,j, common
% transported subtraction, UA5's proved all-mass low-band source inverse,
% and the explicit CAR pinching, two-path commutator and finite-measure
% analyticity arguments above. UA14a--c's same-source-fibre obstruction
% was derived and saved before reading RED AC40--AC42; UA14d is the
% finite-dimensional linear-algebra extension saved after that comparison.
% The preceding
% optional R intake is unrelated conditional scattering/state evidence
% and supplies no premise to this source-invariance proof. No external
% X, source projection, time refit, new field or scientific run is used.

% ACTUAL AC40--AC42 COMPARISON / SOURCES [R].
% After the substantive UA14 save, read the complete independently saved
% AC40--AC42 and its SOURCES in redteam_declared_local_area.tex. Direct
% rechecking gives the exact support lemma: if P+D is CAR and Df=0,
% positivity of C(1-C) gives DPf=0, so supp(D) reduces P. If additionally
% trace-class D commutes with bounded h, each nonzero eigenspace of D is
% finite dimensional and h-invariant; its support is therefore point
% spectral and also reduces P. This necessary condition applies to a
% PHYSICAL stationary excess, not automatically to UA13d's signed
% Cesaro comparator or to a state with constant sources alone.
%
% For the specified finite-mass J=0 branch, N=1 and v tends to zero.
% Its finite-range point-flow perturbation is norm-limit finite rank,
% giving the stated free essential bands. Outside the closed free
% bands the free resolvent has an exponentially summable kernel. The
% outer flow can be made a contraction also after a sufficiently small
% exponential conjugation, while the inner inhomogeneity has finite
% support. The weighted Neumann solution equals the unweighted one by
% uniqueness. Thus finite off-band eigenstate fillings have exponential
% source moments, and the already proved initial flat-P0/CAR obstruction
% excludes their exact spherical initial source if nonzero and physical.
% This does not exclude edge accumulation, embedded states, arbitrary
% nonanalytic histories, or all possible invariant coherence subsets.
% UA14 instead tests the full same-source fibre without stationarity.
% The two results have separate hypotheses and neither is promoted to
% an all-history theorem, an absence-of-bound-states claim or new [X].

% ACTUAL AC43--AC45 / CONDITIONAL ORDER PROPAGATION [R].
% Read the complete actual saved AC43--AC45 and SOURCES after UA14d.
% AC43 independently recomputes the two paths and full4 spin factor,
% count-nullvector, CAR pinching and finite-measure obstruction; its
% checks agree. AC45 verifies the arbitrary finite number of independent
% derivative constraints, including the explicit finite-linear-moment
% restriction. The comparison leaves special invariant subsets open.
%
% AC44 supplies a useful separate sufficient propagation theorem. Its
% proof was independently saved by RED before author-source intake;
% the author had considered this route in reasoning but had not saved
% it, so no independent author provenance is claimed for this addition.
% If C0,I-C0>=kappa F and F_x>=c(1+|x|)^(-p)>0, 0<=p<infinity, put
% B_p(T)=int_0^T sum_d(1+|d|)^(p/2)||h_d(t)||_infty dt for the
% actual finite-range generator h=sum_d M_(h_d)T_d. On an ALREADY
% EXISTING common-transport interval assume B_p(T)<infinity. With
% W=(1+|X|)^(p/2), the finite-shift weight ratio and backward Dyson
% expansion give ||W U(T)^dagger W^-1||<=exp B_p(T). Therefore
% ||F^-1/2 U(T)^dagger Pi_B||<=c^-1/2 exp(B_p(T)) w_B,
% w_B=max_(x in B)(1+|x|)^(p/2). Cauchy--Schwarz in the form of
% U F U^dagger then gives the ACTUAL global operator bounds
% C(T),I-C(T)>=kappa c exp(-2B_p(T)) w_B^-2 Pi_B.
% The calculation was rechecked directly, including the backward
% propagation and the form-domain use of the unbounded inverse weight.
% UA14c has kappa=1/2 and a polynomial lower envelope at fixed L,g,mu,
% so the required finite-site margins persist on any such supplied
% interval without a gap or an instantaneous sea. The bound need not
% be uniform on a collapse-time or allowed-ell family. It proves neither
% existence nor boundedness of that history, and alternative data at a
% reached slice are not assigned the same earlier coupled trajectory.

% FINAL ACTUAL SOURCE-FIBRE READ / SOURCES [R].
% Both complete current local-area sources and every SOURCES block were
% reread in untruncated ranges, including UA14a--d, AC40--AC45, both
% exact compact statements and the saved RED final disposition. The
% count-nullvector permits a zero component; the global finite-site
% operator margin covers the actual disjoint witness inputs. It is not
% inferred from reduced-state faithfulness. The full4 four-path trace,
% all-mass CAR pinching and actual low-band energy-symbol normalization
% agree, with finite expected particle and hole counts relative to the
% SAME actual transported reference. Alternative reached-slice data are
% not assigned the same earlier self-consistent incoming history.
%
% These fixed physical perturbations preserve every initial source and
% both counts but give independent nonradial first-energy derivatives.
% Nonzero radial derivatives remain allowed. Arbitrarily many disjoint
% witnesses with their own global margins exclude closure by finitely
% many additional LINEAR moments; no special nonlinear or infinite
% invariant subset or pure-state history is excluded. The weighted
% propagation comparison retains 0<=p<infinity and applies only to an
% already existing bounded finite-time history from the exhibited
% algebraic mixed preparation. It supplies neither that history nor
% uniform positive-clock, time or allowed-ell bounds.
%
% The separate stationary CAR-support and off-free-band localization
% tests retain their physical initial-P0 hypotheses. Constant sources
% do not imply a stationary covariance, and the signed Cesaro source
% comparator is not silently made physical. No all-bound-state or
% all-history exclusion follows. The point/cell allocation, variable
% lapse, either momentum-row sign, every finite ma and g=G/a^2 domains
% remain explicit. The conditional area theorem and prior X statements
% are unchanged. No new external X, code/result, run, source projection,
% time refit or additional review role is introduced. No mandatory
% scientific repair remains; author source is READY/HELD for handoff.

% UA15. POSITIVE-PARTICLE SOURCE FIBRES: INDEPENDENT SPIN-KERNEL ROUTE [R,S].
% Substantive continuation saved before reading RED AC46--AC49. The
% published compact theorem is unchanged. The aim is the actual native
% common-transport source problem, not a replacement of its point readout
% or an assumption that radial time derivatives vanish. The construction
% below supplies exact physical zero-hole source-null tangents and a
% detecting criterion; the actual spherical-flow detection is a separate
% step, not asserted merely from that criterion.
%
% UA15a [R]. Put a=1, mu=ma>=0, P=P_+(h0), n_k=d(k)/E(k), and
% S(k,p)=P(k)P(p). On any compact momentum-pair set away from nodes and
% n_k=-n_p, c(k,p)=Tr S(k,p)=1+n_k.n_p has a positive lower bound.
% For a smooth matrix kernel T(k,p)=P(k)T(k,p)P(p), define
%   T_perp(k,p)=T(k,p)-[Tr T(k,p)/c(k,p)] S(k,p).
% This is the Hilbert--Schmidt orthogonal projection to trace-zero
% kernels: Tr[S(k,p)^dag Z(k,p)]=Tr Z(k,p) for every such band block,
% and ||S(k,p)||_HS^2=c(k,p). A symmetric real smooth momentum cutoff
% and T(p,k)=T(k,p)^dag make the resulting Z Hermitian, trace class,
% positive-band supported and rapidly decreasing in both site indices.
% It obeys Tr_spin Z(k,p)=0 at EVERY momentum pair, not just after an
% integral. Consequently its exact native cell energy and each current
% vanish at every site:
%   Tr(e_x Z)=0,  Tr(j_xi Z)=0,  Tr Z=0.
% Indeed the positive-band energy vertex reduces to (E_k+E_p)/2 times
% the spin trace, while the two-step current vertex is the scalar
% [Q_i(k)+Q_i(p)]/2 times that trace. No low-momentum expansion, angular
% average or vanishing-gradient assumption enters these identities.
%
% The CAR domain is a positive-band ORDER margin, not UA14's stronger
% margin on the entire one-body space. Suppose 0<=Gamma<=P and
%   Gamma, P-Gamma >= kappa P F P,
% where kappa>0 and the positive scalar F_x is bounded below by a
% polynomial tail c_F(1+|x|)^(-p_F), 3<p_F<infinity. Since Z has rapid
% decrease in both indices, T_F=F^(-1/2) Z F^(-1/2) is bounded. Thus
% +/-Z<=||T_F|| F, and PZP=Z implies +/-Z<=||T_F|| PFP.
% For |s| ||T_F||<=kappa, Gamma+sZ stays between 0 and P. Therefore
% C_s=P0+Gamma+sZ is CAR-admissible, has NO holes, and retains its
% exact finite expected particle count if Gamma is trace class. All
% initial native rho,J, hence the actual fixed readout and solved fields,
% are identical for this entire fixed-preparation interval. The already
% constructed compensated Gamma=P F P has this margin when 0<F<=theta<1;
% its algebraic tail and exact radial initial readout are retained.
% At mu=0 the cutoff avoids the finite nodes; no mass gap is assumed.
%
% UA15b [R,S]. Let an actual common initial point Hamiltonian h be
% fixed by those common sources, and let A=i[h,e_w], with e_w the
% native energy vertex of a real zero-angular-average continuum test
% through the SAME readout. Summable w makes A trace class for bounded
% h. Compress A to the positive band. If its spin-trace-orthogonal
% part (P A P)_perp is nonzero on one compact pair set as above, there
% is a smooth Hermitian Z of UA15a with Tr(A Z)!=0. For smooth kernels
% take a nonnegative symmetric cutoff times (PAP)_perp; its pairing
% with A is the integral of the positive squared Hilbert--Schmidt norm.
% For a general trace-class A, smooth band-kernel approximation on that
% compact set gives a nonzero pairing while keeping the trace-zero
% constraint exactly. Thus the actual initial averaged energy derivative
% changes by s Tr(AZ), although all initial sources/counts are identical.
% At most one member can have zero value of this angular derivative.
% This would exclude invariance of that positive-particle source fibre,
% not a special invariant subset, pure state or actual coupled history.
% A proof that the criterion holds for the solved radial field is still
% required; a nonconstant prescribed flow alone is not taken as that proof.
%
% A useful exact algebraic diagnostic is available without a run. For
% any real four-vector b and nonparallel unit n_k,n_p, set
% T=P(k)(b.Gamma)P(p). Clifford tracing gives
%   ||T_perp||_HS^2=(1-n_k.n_p)|b_perp|^2,
% where b_perp is orthogonal to span(n_k,n_p) in R4. At N=1, the
% spin-vector part of i[V,e_w], V=sum{v_i,Q_i}/2, has the exact kernel
%   b(k,p)=(i/2)int [V(k,r) what(r-p)-what(k-r)V(r,p)] d(r) dr/(2pi)^3,
% up to endpoint terms whose positive-band blocks are multiples of S.
% Here V(k,r)=sum vhat_i(k-r)[Q_i(k)+Q_i(r)]/2. For distributional
% vhat this is read through the original bounded finite-range product.
% The real-vector norm identity extends sesquilinearly to complex b.
% A Fourier diagnostic with a longitudinal flow transfer epsilon e1
% and test transfer epsilon e2 contains the intermediate difference
% d(p+epsilon e2)-d(p+epsilon e1). At p=0, mu>0, this vector is
% orthogonal to the two endpoint d vectors and nonzero. At mu=0 use
% p=t e3 away from a node instead. This verifies an available spin
% direction for such coefficients; these individual plane waves are
% NOT a localized actual spherical source or its required angular test.
% The actual field/test convolution and its nonzero angular detection
% remain the marked step before applying the criterion to a coupled slice.
%
% SOURCES / INDEPENDENCE. Full current binding GOAL reread. This block
% uses native full4 Clifford traces, the existing exact energy/current
% vertices, the paired UA5/UA14 positive-band preparation and direct
% trace-class/form estimates. It was saved before reading the original
% RED's new positive-particle source-fibre continuation. No new external
% source, code/run, changed allocation, refitted state or X is used.
% g=G/a^2 enters the actual solved h; all form domains above retain
% fixed finite mu, the actual ell and existing preparation/history scope.

% UA15c. ACTUAL AC46--AC49 COMPARISON: A COMPLETE CURRENT-FIBRE TEST [R].
% Read the original RED's complete independent AC46--AC49 and SOURCES,
% then its exact AC14--AC16 patch definitions and native response symbols.
% The following route is attributed to that independent save, not to
% UA15a's different spin-kernel calculation. It closes the actual initial
% current-detection question on its stated positive-band margin domain;
% it does not by itself evaluate UA15b's distinct energy-spin criterion.
%
% Let chi_A be the two smooth patches about +/-pi/4 e1, and chi_B the
% eight stationary +/-pi/2 patches, with the RED's amplitude 1/2 on
% each B patch. All use its fixed even product bump of radius .01.
% They are real, separately even, disjoint, avoid the sine nodes, and
% chi_A^2+chi_B^2<=1. Define the ACTUAL full4 symbols
% W_j(q)=int chi_j(k+)chi_j(k-)(E_++E_-)
%                         Tr[P(k+)P(k-)]/2,
% B_(j,i)(q)=int chi_j(k+)chi_j(k-)[Q_i(k+)+Q_i(k-)]
%                         (E_+-E_-)Tr[P(k+)P(k-)]/2,
% where k+/-=k+/-q/2 and integrals use d^3k/(2pi)^3.
% At q=0, w_j=W_j(0)=2int chi_j^2 E>0 and
% p_(j,i)=partial_i B_(j,i)(0)=2int chi_j^2 Q_i^2/E.
% Cubic symmetry gives p_(B,1)=p_(B,2). The axial patch has
% p_(A,1)/w_A>=.24/(mu^2+.6),
% p_(A,2)/w_A<=10^-4/(mu^2+.4), so p_(A,1)>p_(A,2) at every
% finite mu>=0. These are bounds on the fixed supports, not fitted
% dispersion moments or extra fermion species.
%
% The construction works both for RED's r^-8 envelope/readout and for
% our ALREADY FIXED UA5 phi,K with r^-12 tails. To check the latter
% without identifying the two kernels, set f_L(x)=phi(x/L) and
% H_L(x)=-(Delta_X phi)(x/L). Its central Fourier transform is
% L^5|q|^2 phihat(Lq), its lattice sum is exactly zero at alias-free L,
% and |H_L|<=C f_L by the proved derivative/tail estimates. Put
% U_A=W_B*H_L/(w_A+w_B), U_B=-W_A*H_L/(w_A+w_B),
% Z_L=P chi_A U_A chi_A P+P chi_B U_B chi_B P.
% The smooth W_j kernels have every absolute weighted moment, so
% |U_j|<=C_U f_L and sum U_j=0 for either envelope. Exactly
% W_A U_A+W_B U_B=0, while time reversal makes every J_x of Z_L zero.
% Its trace also vanishes: Tr Z_L=2sum_j<int chi_j^2>sum_x U_j=0.
%
% Define B_L=sum_(j=A,B) P chi_j f_L chi_j P. For ANY trace-class
% 0<=Gamma<=P satisfying
%        Gamma, P-Gamma >=eta B_L, eta>0,
% and exactly radial positive initial averaged rho with J=0, the
% covariances P0+Gamma+s Z_L are physical for |s|C_U<=eta. They have
% the same finite expected particle count, zero holes, and exactly the
% same native source arrays. All other coherences of Gamma are allowed.
% A nonempty example uses a smooth even low-transfer cutoff zeta and
% Ghat_j=w_j zeta/W_j, chosen where W_j!=0 and zeta=1 on the target
% support. With A_L=gamma/[g(w_A+w_B)L^2], set
% F_j=A_L G_j*f_L and Gamma_*=sum_j P chi_j F_j chi_j P.
% The kernels have sum one and zero first moments, hence
% F_j=A_L[f_L+O(L^-2 f_L)] and 0<F_j<=1/2 for sufficiently large L.
% Thus Gamma_*>=A_L B_L/2 and P-Gamma_*>=B_L/(2||phi||_infty).
% Its exact cells are e_x=gamma f_L/(g L^2), J_x=0. The fixed UA5
% point readout is exactly radial in its alias-free 1<<ell<<L range,
% and has the regular positive branch N=1, v=-sqrt(2gM/r)rhat for
% either momentum sign. N_f=2A_L(sum_j<int chi_j^2>)sum_x f_L is
% finite and proportional to L/g at fixed mu,gamma. This is a mixed
% Gaussian initial slice, not a supplied earlier collapse or an interval.
% The domain also includes the current-tangent four-channel preparation
% of UA7 at sufficiently large L: every channel envelope is bounded
% below by a positive multiple of f_L and can be made <=1/2. Its axial
% and stationary multipliers are e^3 chi_A and 2e^3 chi_B, respectively,
% because UA7's one-dimensional seed is e times the RED seed. Thus the
% same B_L lower margin follows; its upper margin follows from
% Gamma<=P/2. That preparation's exactly radial dot J is generally
% nonzero. Thus the tested domain does not impose a zero derivative
% as an extra condition; UA9 retains its separate energy-jet limitation.
%
% On EVERY member of this source fibre the ACTUAL initial h is the same.
% Since Z_L is time-reversal even whereas i[h_v,j_xi] is odd, the flow
% contribution to its current derivative is EXACTLY zero. Common
% transport of C and P0 supplies dot D=-i[h,D]; no spectral refill or
% extra reference response is used. The remaining native difference is
% delta dot Jhat_i(q)=-i s Hhat_L(q)
%     [B_(A,i)(q)W_B(q)-B_(B,i)(q)W_A(q)]/(w_A+w_B).
% The bracket is q_i d_i+O(|q|^3),
% d_i=w_B p_(A,i)-w_A p_(B,i), and
% d_1-d_2=w_B[p_(A,1)-p_(A,2)]>0. At q=(t,t,0), sufficiently small
% and nonzero, q_1 delta dot Jhat_2-q_2 delta dot Jhat_1!=0 for s!=0.
% The ACTUAL Khat(ell q) is nonzero near zero at every finite ell,
% so the averaged derivative has nonzero curl too. A radial vector
% must have Fourier transform parallel to q. Thus at most one member
% of this physical affine family has a radial derivative, which may
% itself be nonzero. The entire positive-particle source fibre is not
% invariant. No exceptional member, pure state or special invariant
% nonlinear class is excluded by this argument.
%
% UA15d. Finite linear supplements and the precise history limit [R].
% The RED's disjoint-annulus extension also checks. Replace H_L by any
% finite list of real radial Schwartz h_j(x/L), with disjoint small
% Fourier annuli away from zero, within the fixed nonzero low-transfer
% neighborhood of Khat and W_j. The profiles have zero lattice sum
% and are bounded by C_j f_L; their actual derivative curls have disjoint
% Fourier supports and are nonzero. Any k specified real linear moments
% cost at most k dimensions. With n>k, a nonzero admissible combination
% preserves those moments and every native source/count, yet has a
% nonradial derivative. The fixed annuli can be chosen once for all
% ell/L<=delta0 in the alias-free hierarchy; no time or ell refit occurs.
% This excludes closure of the WHOLE margin-domain fibre by finitely
% many additional linear moments. It does not exclude an infinite or
% nonlinear invariant selection. At a later slice the occupied reference
% is P_T, not P0, so this incoming-band construction is not silently
% reused there. Its common initial N=1,J=0 premise covers both momentum
% signs; a variable-lapse reached slice needs its own native calculation.
%
% SOURCES / ACTUAL COMPARISON. RED AC46--AC49 and its prerequisite
% AC14--AC16 were actually read after the independent UA15a--b save.
% Their compensation, form bound, native current kernel, time-reversal
% argument and small-transfer curl were checked directly. The extension
% from its fixed r^-8 Schwartz readout to our fixed UA5 r^-12 algebraic
% K uses only the already proved weighted convolution, positivity and
% alias identities; the kernels themselves are not equated. This is
% original-pair R agreement, not an external X or an actual invariant
% history. The full coupled existence, uniform clock/time/ell premises,
% emitted Gauss energy and accessible count remain separate requirements.

% FINAL ACTUAL POSITIVE-PARTICLE READ / SOURCES [R].
% Both complete current author and RED local-area sources and every
% SOURCES block were reread in untruncated ranges through UA15a--d and
% AC46--AC50. The exact revised compact statements, explicit B_L domain,
% and corrected e^3 chi_A / 2e^3 chi_B seed conversion were read as saved.
% The independent substantive saves preceded actual counterpart intake.
% No mandatory scientific repair remains at the stated scopes.
%
% The independent spin-kernel projection supplies exact trace-class
% positive-band source/count-null directions under its form margin;
% its actual solved-flow ENERGY detection remains explicitly conditional.
% The RED's independently saved compensated occupation route instead
% supplies a completed CURRENT test on the actual N=1 point-flow initial
% branch. The native W/B symbols, all-mass patch inequalities, exact
% source/count cancellation, time-reversal identity and nonzero curl
% agree. Both the fixed r^-8 and r^-12 envelopes retain their own
% readouts and weighted/alias bounds; their kernels are not identified.
%
% The tested domain has finite expected particles, zero holes and the
% actual initial P0 reference, with arbitrary additional coherences
% permitted by its explicit positive-band margins. It contains the
% four-channel example with nonzero radial dot J. Source-preserving
% physical changes show that the entire fibre, even after finitely
% many additional linear moments, is not an invariant spherical class.
% No vanishing radial derivative is imposed. Both momentum-row signs
% share this initial N=1,J=0 test; later variable-lapse/P_T slices are
% not substituted into the incoming-band calculation. Special invariant
% subsets, pure histories, prior preparations and actual unprojected
% coupled existence remain open. The conditional area theorem retains
% its fixed-time, clock, Fourier, g and allowed-ell premises. No external
% X, scientific run, source projection, refit or extra review is introduced.

% UA16. AN ARBITRARY-COHERENCE INTERVAL OBSTRUCTION AND ITS REGULAR
% WEAK-GRAVITY LIMIT [R; independent substantive continuation].
% The printed core is unchanged pending the original pair's comparison.
% This route was derived and saved before reading any new RED interval
% or charge/work continuation. It allows nonzero radial time derivatives
% and assumes neither stationarity nor a finite preparation ansatz.

% UA16a. Exact free continuity uses a nonradial native divergence.
% Work in lattice units a=1, finite mu=ma>=0, on the infinite parent.
% Let 0<=C0<=I, D0=C0-P0 be trace class and nonzero, and use the ACTUAL
% free common transport D(t)=exp(-it h0)D0 exp(it h0). Its reference
% is P0 because [h0,P0]=0; no instantaneous band refill is involved.
% Define the native cell means
% e_x(t)=Tr[{Pi_x,h0}D(t)/2],
% j_(x,i)(t)=Tr[{Pi_x,Q_i}D(t)/2].
% They are in l1, with ||e||_1<=Emax||D0||_1 and
% ||j_i||_1<=||Q_i||||D0||_1, Emax=sqrt(mu^2+3), ||Q_i||=1/2.
% The same estimates apply to their trace-norm time derivatives.
%
% For transfer q with every |q_i|<pi/2 their Fourier transforms obey
%     partial_t ehat(q,t)=-i sum_i tan(q_i) jhat_i(q,t).       (UA16.1)
% This identity is for arbitrary band/spin coherences. At midpoint k,
% the commutator of h0 with the symmetrized energy vertex reduces under
% the spin trace to [h0(k+q/2)^2-h0(k-q/2)^2]/2
%   =sum_i sin(2k_i)sin(q_i)/2,
% while the symmetrized Q_i vertex is sin(2k_i)cos(q_i)/2.
% Equivalently the same relation follows as an identity of bounded
% Fourier-smeared one-body vertices, paired with trace-class D. Thus no
% pointwise momentum kernel or low-momentum state assumption is needed.

% Apply the SAME fixed translation-invariant readout to e and every
% Cartesian current component: normalized radial K_ell in L1 with
% POINT allocation, or that same convolution after normalized cell-box
% allocation. Its common Fourier multiplier is nonzero in a ball about
% q=0, including for the fixed noninjective cutoff kernel of UA5.
% Multiplication preserves (UA16.1) for the actual continuum averaged
% rho,J. No angular source projection is taken. Suppose both are exactly
% spherical on a nonempty open time interval: rho is a radial scalar
% and J a radial vector. Since they are L1, on 0<|q|<q0<pi/2 write
% rhohat(q,t)=F(|q|,t), Jhat_i(q,t)=q_i H(|q|,t).
% Comparing q=r e1 and q=r(1,1,1)/sqrt(3) in (UA16.1) gives
% [r tan r-sqrt(3)r tan(r/sqrt(3))] H(r,t)=0.
% The bracket is strictly positive for 0<r<pi/2, since tan r/r is
% strictly increasing. Consequently H=0 and partial_t F=0 throughout
% that low-transfer ball. Radial derivatives were allowed; the exact
% native equation has forced this particular low-transfer staticity.

% UA16b. That staticity is impossible for a nontrivial physical
% finite-excitation covariance, including nonanalytic algebraic tails.
% For any Schwartz test psi whose Fourier support lies in the preceding
% ball, the averaged energy pairing is Tr[A_psi D(t)], with
% A_psi={w_psi,h0}/2 trace class. Indeed
% sum_x|w_psi(x)|<=||K_ell||_1 sup_y sum_x|psi(y+x)|<infinity;
% the normalized cell-box factor has the same bound. By UA16a this
% pairing is constant on the open interval. Since h0 is bounded, it
% extends to an entire function of complex t with trace-norm bound
% proportional to exp(2 Emax |Im t|). It is therefore constant for all
% real t, without assuming spatial analyticity of D0 or of its sources.
%
% The free sine Hamiltonian has no point spectrum: each level set of
% +/-sqrt(mu^2+sum_i sin^2 k_i) has Lebesgue measure zero. In its
% momentum kernel the Hilbert-Schmidt Cesaro average of D(t) tends to
% zero. The time multiplier tends to zero off equal-energy pairs,
% which have measure zero, and is bounded by one; dominated convergence
% applies to the square-integrable kernel of each band block. Pairing
% with the Hilbert-Schmidt A_psi shows that the preceding constant is
% zero. Hence the continuous rhohat(q,0) vanishes in a ball about zero.
% In particular its value at zero, E_D=Tr(h0 D0), vanishes.
%
% CAR contradicts this last conclusion. With E=|h0|,
% E_D=Tr[E P_+ C0 P_+]+Tr[E P0(I-C0)P0]>0
% for every C0!=P0 with trace-class D0. Both terms are nonnegative;
% E has no l2 kernel even at mu=0. If both vanish, positivity of C0
% and I-C0 also removes their off-diagonal band blocks, giving C0=P0.
% Thus NO nontrivial such FREE history has exactly spherical averaged
% rho and J on any open interval, for any finite mu and the stated
% fixed readout. This covers arbitrary physical coherences, particles
% and holes, and nonzero radial derivative candidates. It is not merely
% a failure of one preparation or of one finite number of tangent tests.
% The noninjective positive initial examples above make the initial-data
% domain nonempty; this theorem excludes their whole free interval,
% not the exact radiality already established at their initial slice.

% UA16c. A necessary singularity of any exact spherical weak-g branch.
% Consider actual common-transport point-field histories at g>0 on a
% common fixed interval [0,T0], T0>0, with initial reference P0,
% physical C_g(0) and D_g(0)=C_g(0)-P0 trace class. Suppose
% D_g(0)->D_*!=0 in trace norm and
% delta_g=int_0^T0 ||h_g(t)-h0|| dt ->0.                (UA16.2)
% Keep the preceding readout FIXED (or convergent in its source L1
% operator norm to such a fixed readout); do not let its cutoff collapse
% to q=0 as g decreases. Duhamel and common transport give, uniformly,
% ||D_g(t)-exp(-it h0)D_*exp(it h0)||_1
% <=||D_g(0)-D_*||_1+2||D_*||_1 delta_g.
% The source maps above are bounded from trace class to continuum L1.
% Exactly radial scalar/vector functions form closed subspaces there.
% If every actual g-history were exactly spherical, its limit would
% therefore be a forbidden free history from UA16a--b. The limiting
% covariance P0+D_* is physical by norm closure. Hence such a regular
% branch CANNOT exist. This is a necessary condition on actual coupled
% histories, not a statement that all finite-g histories have this limit.

% The condition has a concrete Gauss-law sufficient domain, for either
% printed momentum sign. In dimensionless variables let actual spherical
% fields obey N(infinity)=1, M=r(v/N)^2/(2g)>0 for r>0, and assume
% uniformly in t,g that
% B_M=sup_(r,t,g) M(r,t)/r<infinity,
% B_J=sup_(t,g) int_0^infinity r^(3/2)|J(r,t)|/sqrt(M(r,t)) dr<infinity.
% The absolute radial momentum row gives
% |log N(r,t)|<=2sqrt(2)pi sqrt(g) B_J,
% ||v||_infinity<=exp(2sqrt(2)pi sqrt(g)B_J)sqrt(2g B_M).
% Its sign is immaterial in these absolute inequalities; it is not
% repaired. For the native point operator, writing c=2sqrt(2)pi B_J,
% ||h_g-h0||<=Emax[exp(c sqrt(g))-1]
%                 +(3/2)exp(c sqrt(g))sqrt(2g B_M).
% Thus fixed T0 and these ACTUAL positive-mass/clock/source controls
% imply (UA16.2). Their existence is not assumed as a consequence of
% CAR, finite count or radial projection. Center and exterior integrals
% must have the stated finite bounds along the actual reached history.

% Restoring units, t/a is the dimensionless time, M is a times physical
% enclosed energy, r is physical radius/a, mu=ma and g=G/a^2. This is
% an exact finite-spacing theorem. It allows a fixed hydrodynamic scale
% hierarchy if the actual family supplies it, but no uniform allowed-ell
% or large-time premise is manufactured. Branches with noncompact
% initial covariance limits, diverging count/size, collapsing readout
% transfer support, shrinking time interval or nonsmall integrated
% fields are NOT excluded; nor are special histories at a fixed nonzero
% g. In particular a continuum/object scaling that changes those data
% cannot be replaced by this fixed-data limit. Actual finite-g nonlinear
% invariance, its positive clock/mass and physical emitted-energy/count
% requirements remain open. No time refit, HA coupling, source projection,
% strong operator-sector replacement or additional quantum species is used.
%
% SOURCES / INDEPENDENCE: actual h0, Q_i and site-symmetrized vertices;
% bounded-generator time analyticity, direct free momentum-kernel Cesaro
% averaging and CAR positivity, all proved above. The actual radial Gauss
% row enters only through its displayed absolute-value consequence.
% UA16 was saved before any new RED free-interval/weak-g derivation or
% charge/work source was read. Existing source-regularity, finite-family
% and source-fibre results are retained, not treated as this proof.
% No new external X, scientific run or complete coupled history is claimed.

% UA16d. Finite expected excitation count suffices; trace-class excess
% is not necessary for the FREE interval obstruction. This strengthening
% was saved before a new RED derivation of it was read. Keep 0<=C0<=I
% but assume only
% N_ex=Tr[P_+ C0 P_+]+Tr[P0(I-C0)P0]<infinity,
% with C0!=P0. The difference D0 can fail to be trace class. Nevertheless
% D0^2<=P_+ C0 P_++P0(I-C0)P0 implies ||D0||_HS^2<=N_ex.
% Its diagonal band blocks are trace class; its off-diagonal blocks
% need only be Hilbert--Schmidt. This explicitly includes coherent
% finite-excitation preparations not covered by a trace-norm premise.

% Native e_x,j_(x,i) are now l2 sequences, not asserted l1. Their
% averaged sources are well-defined distributions by the actual
% Schwartz readout tests, whose vertices are trace class and hence
% Hilbert--Schmidt. In a ball inside the reciprocal cell their Fourier
% transforms are L2 functions multiplied by the bounded common readout
% factor. The exact continuity identity holds there in L2. Exact radial
% scalar/vector distributions therefore have rhohat=F(r,t),
% Jhat_i=q_i H(r,t) almost everywhere on this ball. Subtract the angular
% mean of A(r,n)=sum_i r n_i tan(r n_i) in (UA16.1). Its angular
% variance is strictly positive at each 0<r<q0, by the axis/diagonal
% inequality already proved. Thus H=0 almost everywhere and F is
% time independent there. This avoids evaluating an L2 Fourier function
% at a selected direction. Hilbert--Schmidt time analyticity and the
% SAME band-kernel Cesaro proof then show that this low-transfer energy
% distribution vanishes, just as in UA16b. No spatial moment is assumed.

% To identify its total energy without an unjustified sum of signed
% native cells, use a smooth exhaustion. Fix real Schwartz phi with
% compact Fourier support about zero and phi(0)=1, and put
% w_R(x)=int phi(y/R)K_ell(y-x)d^3y,
% A_R={w_R,h0}/2, with the same cell-box factor if declared. Each A_R
% is trace class, w_R is uniformly bounded and tends pointwise to 1.
% At sufficiently large R its no-alias Fourier transform is
% R^3 phihat(Rq) Khat(ell q), with the optional bounded cell multiplier.
% The already paired exact native interband formula QV1--QV5 in
% declared_exchange_kernel.tex gives
% ||P0 A_R P_+||_HS^2
%   =int |what_R(q)|^2 B_mu(q)d^3q/(2pi)^3 <= C_mu,K,phi/R,
% using B_mu(q)<=C_mu|q|^4 at EVERY finite mu, including zero.
% The massless estimate retains the eight point-node proof; a gap is
% not inserted. The source vertices are still those of the actual
% readout, not an angular projection or a replacement global charge.

% The trace-class diagonal blocks of D0 therefore give, by bounded
% strong convergence of A_R, the positive canonical free energy
% E_free=Tr[E P_+ C0 P_+]+Tr[E P0(I-C0)P0].
% The off-diagonal contribution is bounded by C R^(-1/2)||D0||_HS
% and tends to zero. Hence Tr(A_R D0)->E_free>0. But the established
% low-transfer energy vanishing makes Tr(A_R D0)=0 for every sufficiently
% large R, a contradiction. This proves the free interval theorem for
% ALL physical finite expected particle-plus-hole counts, arbitrary
% coherences and non-Gaussian states through their covariance. It does
% not replace actual excitation count by finite net charge or presume
% absolute summability of the local energy in this larger domain.

% The regular coupled-branch conclusion likewise extends if
% D_g(0)->D_*!=0 in Hilbert--Schmidt norm, the initial excitation counts
% are uniformly finite, and (UA16.2) plus the fixed readout hold. Fatou
% on increasing finite-mode number tests makes the limiting covariance
% finite-count and physical. Duhamel gives the displayed bound of UA16c
% with Hilbert--Schmidt norms. Every fixed readout test has a Hilbert--
% Schmidt vertex, so the actual sources converge as distributions;
% radiality is closed in that topology. The forbidden free limit again
% follows. For varying readouts one must instead supply convergence
% of these actual test vertices in Hilbert--Schmidt norm. This extension
% supplies no compactness of an arbitrary family and no bound on a
% field/history that has not actually been constructed. Singular large-
% count/size/time/readout limits and fixed nonzero-g histories remain
% open. The trace-class proof above is retained as a separate simpler
% route, not silently assigned this broader domain.
%
% SOURCES / ADDITIONAL INDEPENDENCE: direct CAR square inequality,
% Hilbert--Schmidt Fourier/propagator estimates, and the already fully
% paired native interband variance bound QV1--QV5. This extension was
% derived after the actual EC11/QR11 comparison but before reading any
% new RED finite-count extension of UA16; those charge/work and normality
% proofs are distinct from the smooth-exhaustion argument here.

% UA16e. ACTUAL AC51--AC54 COMPARISON / SOURCES [R].
% Read all of the actual saved AC51--AC54 after UA16a--d. RED checked
% UA16a--c after its distinct EC11 charge/work save; its AC52 finite-count
% extension was derived after that intake but before reading UA16d.
% The author's UA16d was likewise saved before AC52 intake. The direct
% CAR square, distributional continuity and smooth-energy exhaustion
% proofs agree, including the full4 interband normalization, R^-1
% variance and the massless point-node bound. Neither proof identifies
% a non-trace-class excess with an unspecified Tr(h0 D).

% AC53 strengthens the author's strong-HS sufficient limit correctly.
% A uniform N_ex bound gives a weakly sequentially compact HS ball.
% For any weak-HS limit D_*, finite-rank matrix tests make P0+D_*
% physical; increasing finite particle/hole projections and Fatou
% give finite excitation count. Uniform Duhamel control against the
% free propagation of the SAME D_g(0) costs at most
% 2 sqrt(sup N_ex) int||h_g-h0||. Every actual tested source therefore
% converges distributionally to that of D_*. Exact spherical symmetry
% is closed under this limit, so UA16d forces D_*=0. All weak cluster
% points being zero, D_g(0) tends weakly to zero in HS. This does not
% force its global count or canonical free energy to tend to zero;
% the family may lose local compactness. In particular any fixed finite
% excitation-mode subspace with a positive uniform occupation rules
% out this regular branch. The printed conclusion uses this checked
% broader weak limit, not an assumed trace-norm compactness theorem.

% Also checked AC54's actual finite-g necessary test. Let B range over
% native vertices obtained from Schwartz energy tests minus their
% rotations, and covariantly rotated current-test differences. With
% ||B||_HS<=1, the bounded variational diagnostic
% d_*(D0,T)=[2||D0||_HS]^-1 sup_(t in [0,T],B)
%                  |Tr[D0 exp(it h0)B exp(-it h0)]|
% lies strictly between zero and 1/2 inclusive for any fixed physical
% finite-count C0!=P0 and T>0. Otherwise all tested free sources would
% be spherical on an open interval, contradicting UA16d. Any genuine
% spherical history from those SAME data must satisfy
% int_0^T||h-h0||dt>=d_*. This follows directly by subtracting its zero
% angular pairing and using the HS Duhamel bound. The actual B_M,B_J
% estimate above gives AC54's stated sufficient exclusion inequality.
% This is not a uniform threshold over changing data or an existence
% result above it; no prescribed profile is substituted into its bounds.

% One useful actual-charge consequence of the weak-HS result is local.
% On a regular family with B_M bounded, ||u_g||_infinity=O(sqrt(g)).
% At any FIXED finite radius R the actual radial mass row reads
% M_g(R,t)=int_(|y|<R)rho_g(y,t)d^3y
%                    +s int_(|y|<R)u_g(|y|,t)J_g(y,t).yhat d^3y,
% with s=+1 compatible and s=-1 literal. The first term is a fixed
% trace-class energy vertex paired with D_g(t), hence tends to zero.
% For the second, its readout weights obey
% |w_(g,i)(x)|<=||u_g||_infinity b_R(x),
% b_R(x)=int_(|y|<R)|K_ell(y-x)|d^3y in l1 and l2.
% The cell-box allocation has the same bound with the effective kernel.
% Its current vertex thus has HS norm O(sqrt(g)); bounded ||D_g||_HS
% makes the second term vanish. Therefore M_g(R,t)->0 at every fixed
% R,t. Uniformity on the common compact time interval follows from
% norm-continuity of the free test orbit and the uniform Duhamel bound.
% This excludes a uniformly positive enclosed mass at a fixed radius
% in that regular bounded-count limit, not a growing-radius object or
% a finite-g interval. No total boundary-charge vanishing is inferred.

% Complete general finite-g nonlinear invariance, reached positive
% clock/mass and allowed-ell bounds, detector/Gauss transfer and the
% accessible count remain actual tasks. The finite-count interval proof
% is a materially broader necessary obstruction, not a replacement
% source map or an assertion that weak-g compactness must hold for every
% object scaling. No additional numerical run or external X is used.

% UA17. FINITELY LOCAL QUANTUM PREPARATIONS AND ACTUAL REACHED SOURCES
% [R; independent substantive continuation, published core unchanged].
% This is a preparation-domain implication of the already paired source-
% analyticity obstruction, not another fitted covariance family. It is
% saved before any new original-RED route on this local-preparation
% implication is read. It retains arbitrary correlations and all native
% coherences created by the preparation. A special nonlocal, nonanalytic
% source class is not excluded, and no stationary-state premise is added.

% UA17a. Work on the infinite native parent, in lattice units, at every
% finite mu>=0 and g>0. Let omega_0 be the incoming pure sea P0. Fix a
% FINITE set A of sites and a trace-preserving local quantum channel,
% with Heisenberg action Lambda(O)=sum_alpha L_alpha^dag O L_alpha,
% L_alpha in CAR(A), sum_alpha L_alpha^dag L_alpha=I. One may take
% parity-homogeneous Kraus operators, as appropriate to a physical
% parity-preserving channel. No Gaussianity or quadratic preparation is
% imposed. Let omega_in=omega_0 composed with Lambda, and let C_in be
% its physical one-body covariance. For x,y outside A the even bilinear
% a_y^dag a_x commutes with every L_alpha, so its expectation is unchanged.
% Hence, with D0=C_in-P0 and Pi_A the full4 site projection,
% (1-Pi_A)D0(1-Pi_A)=0,
% D0=Pi_A D0+D0 Pi_A-Pi_A D0 Pi_A.                    (UA17.1)
% In particular D0 is trace class of rank at most 2 rank(Pi_A), with
% ||D0||<=1. Positivity of C_in makes its particle and hole diagonal
% excitation blocks positive, so the actual expected excitation count
% is finite and at most ||D0||_1<=2 rank(Pi_A). The initial positive
% free excitation energy is bounded by Emax times that count and is
% strictly positive if C_in!=P0, also at mu=0 since the nodes are null
% sets. This is a physical state preparation, not an arbitrary Hermitian
% perturbation. A nonselective channel is essential: conditioning on
% a local outcome need not leave distant observables unchanged.

% UA17b. Supply any actual bounded finite-range POINT history h(t) on
% 0<=t<=T<infinity, and transport BOTH the physical state and its incoming
% reference by the same CAR automorphism. Thus D_t=U D0 U^dag, even
% when omega_in is non-Gaussian or U has no incoming-Fock implementer.
% For each lambda>0 define J_lambda(t) as the larger of the row and
% column sums of ||h_xy(t)|| exp(lambda|x-y|). It is finite for bounded
% native coefficients because the hopping range is at most two.
% The Dyson series and the triangle inequality for path lengths give
% the corresponding weighted row/column bound for U by
% exp(int_0^T J_lambda(t)dt). Therefore for the fixed finite-radius
% support B_x of an actual site energy/current vertex b_x,
% ||Pi_(B_x) U Pi_A||
% <=exp[-lambda dist(A,B_x)+int_0^T J_lambda(t)dt].     (UA17.2)
% Insert UA17.1 in Tr(D_t b_x), using b_x=Pi_(B_x)b_x Pi_(B_x).
% Each of its three terms is at most
% ||D0||_1||b_x||||Pi_(B_x)U Pi_A|| in absolute value. Consequently
% |e_x(t)|+sum_i|j_xi(t)|<=C_lambda,T exp[-lambda dist(x,A)]
% for EVERY lambda>0, uniformly on this supplied finite interval.
% The fixed support radius is absorbed in C_lambda,T. Taking a larger
% lambda proves every exponential source moment is finite. No decay
% of the massless P0 kernel, gap, or weighted trace norm of w D0 was
% assumed; D0 can have long off-diagonal sea tails. This is precisely
% why the local-channel domain is distinct from the weighted-S1
% sufficient source domain checked in the exchange note.

% UA17c. Fix the actual readout BEFORE testing the state: a real normalized
% radial L1 convolution kernel K_ell, and either point allocation or
% the native piecewise-constant cell-box allocation. Its Fourier
% transform is nonzero near zero, even when it is noninjective elsewhere.
% The exponential source bounds in UA17b put the ACTUAL reached arrays
% in the analytic domain of UA13 for every finite supplied time. Thus
% exact spherical averaged sources have the following necessary forms.
% With point allocation about a lattice-site center, e_x=E delta_(x,0)
% and j_xi=0. With cell-box allocation, e_x=0 and j_xi=0 everywhere.
% For the latter statement one deconvolves K only near zero, continues
% the analytic transform, and obtains a radial scalar and covariant
% radial vector that are constant on each open Cartesian cell. A scalar
% jump across a face violates rotational distributional derivatives
% unless the two neighboring constants coincide; integrability then
% makes the scalar zero. A constant vector on an open cell can transform
% as a radial vector only if it is zero. No angular projection is used.

% Consequently a regular-center positive-mass spherical Gauss branch
% with this CELL-BOX readout cannot pass through ANY such finite reached
% time: rho=J=0 gives M'=0 and M(0)=0 for either printed sign. This
% conclusion does not say that the clock is fixed in the zero-source
% v=0 equations. It excludes positive mass there, not lapse freedom.
% It is a finite-g necessary obstruction for the stated local-channel
% preparation/common-history domain, with no requirement that radial
% derivatives be zero imposed on the candidate beforehand.

% For POINT allocation define S_x=(-1)^(x_1+x_2+x_3) and R=S beta.
% The exact native identities {S,h0}/2=mu R and [R,h[N,v]]=0 give
% E=sum_x S_x e_x(t)=mu Tr(R D_t)=mu Tr(R D0)           (UA17.3)
% whenever the reached sources have that central form. The trace and
% sum are legitimate because D_t is S1 and the sources are exponentially
% summable. Therefore massless local-channel preparations, or massive
% ones with Tr(RD0)=0, likewise cannot supply a positive-mass radial
% reached branch at any finite time. For mu>0 with nonzero conserved
% R charge this argument leaves the already stated static central-source
% possibility open; it does not assume that its necessary point-spectrum
% comparator is a physical stationary covariance. With the specified
% strictly positive radial kernel and a regular positive branch, J=0
% yields N=1 and E constant, as in UA13. These are consequences of
% exact radiality, not extra first-derivative restrictions on the goal.

% At the INCOMING slice itself, C_in!=P0 cannot have exactly spherical
% averaged sources for either allocation at any mu. Its compact source
% and finite count are in UA13's physical initial domain, where strict
% free energy exceeds |mu Tr(RD0)|. That argument is not transferred
% to a later reference P_T; UA17.3 is the separate reached-state result.
% Initial configurations obtained by postselection, nonlocal infinite-
% extent channels, or infinite preparation-time limits are not covered
% merely because they are normal finite-count states. Those can have
% the nonanalytic tails left open in UA13 and remain part of the genuine
% coupled-history question. Nor does a sequence of approximate spherical
% sources have to satisfy this exact finite-spacing theorem termwise.

% All claims are at fixed a,mu=ma,g=G/a^2, actual kernel/ell and finite T;
% no uniform macroscopic-time or allowed-ell estimate follows from the
% exponential constants in UA17.2. They depend on the actual bounded
% history, not on a presumed Gauss existence theorem. Finite parents
% and their limits require their actual boundary control. The proof
% neither supplies an unprojected seed nor changes the field/source
% map, reference, Hamiltonian or momentum sign. SOURCES / INDEPENDENCE:
% exact local-channel CAR commutation, finite-range weighted Dyson
% propagation, and the already paired UA13 analytic-source and native
% checkerboard identities; this new preparation-domain implication is
% saved before new original-RED intake. No code, run or external X.

% ACTUAL AC56--AC58 LOCAL-PREPARATION COMPARISON / SOURCES [R].
% Read the complete independent RED continuation after UA17 was saved.
% Its even-observable channel argument and two-block covariance identity
% agree. Its separate finite-range OPERATOR-NORM Dyson bound is valid:
% for B_T=int_0^T||h(s)||ds and distance d, every product of fewer
% than d/2 range-two matrices has zero tested block. Therefore
% ||Pi_C U Pi_A||<=sum_(2n>=d) B_T^n/n!
%                  <=exp[B_T exp(2s)-s d], s>0.
% The two-block form D0=Pi_A D0+(1-Pi_A)D0 Pi_A also gives
% ||Pi_C D_t Pi_C||<=2||Pi_C U Pi_A|| directly, using ||D0||<=1.
% Applying the actual finite trace norms of the source vertices proves
% the same all-exponential reached-source bounds without assumptions
% on the massless sea's spatial decay. This is a distinct proof from
% UA17b's weighted row/column propagation.
%
% RED supplies a further exact hard-number statement. Every Kraus
% operator in the d_A=4|A| mode algebra is a polynomial of degree at
% most 2d_A in its Majoranas. Each factor is linear in incoming
% particle/hole creation and annihilation operators. Acting on the sea
% therefore puts every Kraus vector in excitation sectors 0,...,2d_A,
% and their mixture retains that support. Common quadratic transport
% preserves it in the transported excitation fibre. This strengthens
% UA17a's expected-count bound and uses no Gaussian closure or fixed
% incoming-Fock implementer for the later evolution.
%
% The actual initial strict-energy and point/cell analyticity proofs
% agree. In the reached state their conclusions are only the stated
% central-source/J=0 restriction, the cell-box positive-mass exclusion
% and the massless/zero-checkerboard-charge point exclusion; the flat-P0
% energy inequality is not reused at P_T. A local operation on an
% already reached sea retains the locality proof but not that initial
% inequality. Arbitrary nonlocal finite-count states, postselection,
% infinite preparation-time/region limits and controlled approximate
% spherical histories keep their separate unresolved domains. Both
% substantive saves preceded actual comparison. No new X or run.

% UA18. FINITE-DURATION LOCALIZED NATIVE PREPARATIONS CLOSE THE
% ZERO-CHECKERBOARD-CHARGE POINT CASE AT EVERY MASS [R; independent].
% This substantive continuation is saved before reading any new RED
% reached-point or localized-drive result. It is a domain statement
% for physical initial preparations, not an asserted preparation by
% the self-consistent binding equations or a new interaction in them.
% The compact core is unchanged at this save.

% UA18a. Fix a=1, every finite mu>=0 and g>0, the infinite parent and
% its incoming free sea P0. Choose a finite preparation duration tau
% and real bounded native pulse coefficients with support in one fixed
% finite site region A during -tau<=s<=0. Define
% h_prep(s)=h0+V(s),
% V(s)={delta N(s),h0}/2+sum_i{delta v_i(s),Q_i}/2.
% The pulses may be completely nonradial and vary in time; an actual
% positive preparation clock can be imposed without changing the proof.
% Their finite-range endpoints lie in a fixed enlargement A_* of A,
% so V=Pi_(A_*) V Pi_(A_*), and int||V(s)||_1 ds<infinity.
% Let U_prep be its one-body propagator and define the PHYSICAL
% incoming covariance C0=U_prep(0,-tau)P0 U_prep(0,-tau)^dag.
% Keep P0 as the reference of the subsequently posed initial-value
% problem. This selects a physical state, not a reference refill
% during its later common evolution. Since [h0,P0]=0,
% D0=C0-P0=-i int_(-tau)^0 U_prep(0,s)[V(s),P0]
%                                      U_prep(0,s)^dag ds.
% Hence D0 is trace class,
% ||D0||_1<=2 int||V(s)||_1 ds, and its actual expected particle
% plus hole count is finite, bounded by ||D0||_1. Tr D0=0, so the
% two expected counts are equal. No small pulse, adiabatic limit,
% spectral gap or zero-momentum approximation is used. The prepared
% state is a pure Gaussian CAR state; D0 generally has infinite rank
% and is not the finite-region-channel covariance of UA17.
% Nontrivial preparations exist at every finite mu: one nonzero
% finite-site lapse pulse has [e_x,P0]!=0 by the paired strict native
% interband-variance theorem. A sufficiently short nonzero pulse
% therefore gives C0-P0=-i tau[V,P0]+O(tau^2) nonzero in trace norm.
% A positive pulse keeps the preparation clock positive.

% The exact native symmetry R=S beta, S_x=(-1)^(x_1+x_2+x_3), commutes
% with h0, every diagonal pulse coefficient and every Q_i. Thus
% [R,h_prep(s)]=0 and [R,U_prep]=0. The trace-class Duhamel formula
% and cyclicity for its finite-rank commutator give
%       Tr(RD0)=-i int Tr[R[V(s),P0]] ds=0.          (UA18.1)
% This is the actual conserved checkerboard charge, not net CAR
% particle number. Common later native evolution also commutes with R,
% so Tr(RD_t)=0 at every reached time. The mass mu need not vanish.

% UA18b. The ACTUAL reached sources have all exponential moments even
% though this preparation is quasilocal and its covariance has sea tails.
% Supply any bounded finite-range point-field history on 0<=t<=T,
% with finite fixed T and integrable coefficient bounds. The physical
% covariance and the reference then both use its U(t,0), so
% D_t=U(t,0)D0 U(t,0)^dag. Put
% G(t,s)=U(t,0)U_prep(0,s), -tau<=s<=0.
% Each native source vertex b_x=e_x or q_(x,i) is supported in a fixed
% finite enlargement B_x of x, with uniformly bounded trace norm.
% Inserting UA18a and using one localized endpoint of V gives
% |Tr(b_x D_t)|
% <=2||b_x||_1 int_(-tau)^0 ||V(s)||
%                           ||Pi_(B_x)G(t,s)Pi_(A_*)|| ds. (UA18.2)
% For every lambda>0 the finite-range weighted Dyson estimate gives
% ||Pi_(B_x)G(t,s)Pi_(A_*)||
% <=C_(A_*,B,lambda) exp[-lambda dist(x,A_*)+
%          int_(-tau)^0 J_lambda^prep+
%          int_0^T J_lambda^actual],
% where J_lambda is the actual bounded row/column hopping sum with
% weight exp(lambda|x-y|). All factors are finite at the fixed time.
% Taking lambda greater than any desired source weight proves
% sum_x exp(lambda'|x|)(|e_x(t)|+sum_i|j_xi(t)|)<infinity
% for every finite lambda', uniformly on 0<=t<=T.
% No exponential decay of P0, weighted trace-class assumption on D0
% or instantaneous sea is used: the second leg in UA18.2 is bounded
% by ||P0||=1. In particular the proof includes the massless nodes.

% A finite sequence of additional local trace-preserving CAR channels
% is allowed if each Kraus operator K_l in its finite region B obeys
% [K_l,Q_(R,B)]=0, Q_(R,B)=dGamma(Pi_B R Pi_B). This is neutrality
% under the onsite R action, additional to ordinary trace preservation.
% Each such channel changes the covariance only in rows or
% columns meeting its finite region, by the even-observable argument
% of UA17, and adds a finite-rank difference. Its Kraus neutrality
% preserves the expectation of the centered R charge. Inserting those
% differences between the finite propagator segments proves the SAME
% source localization and finite-count conclusion, by a finite sum
% of UA17 and UA18.2 terms. This extension can be non-Gaussian.
% It assumes a fixed finite number of finite-region operations;
% no uniform infinite protocol or postselection limit is inferred.

% UA18c. Use the SAME fixed real normalized radial L1 point readout,
% or its normalized cell-box allocation, as in UA17. It can be
% noninjective. At ANY finite reached time, suppose the actual averaged
% energy and current are exactly radial scalar and radial vector.
% Source analyticity forces the point arrays to be e_x=E delta_(x,0),
% j_x=0 (center at a native site); the cell-box arrays vanish outright.
% For point allocation the exact checkerboard identity now gives
% E=mu Tr(RD_t)=0 at EVERY finite mu by UA18.1. Thus both allocations
% give rho=J=0. A regular-center Gauss branch has M=0 and cannot be
% positive-mass there, for either printed momentum sign. This excludes
% a positive-mass exactly spherical reached interval for this entire
% preparation class, without requiring radial derivatives to vanish.
% It does not fix the clock in the zero-source v=0 subsystem.

% The preparation/reference distinction is essential. If the pulse is
% itself part of the DECLARED common history starting with physical
% C=P0, its transported reference undergoes that same pulse and D
% remains identically zero; its incoming-reference pairs are not a
% gravitating payload. UA18 instead tests the explicitly selected
% physical covariance above as initial data relative to the flat P0
% of a subsequent initial-value problem. It does not assert an earlier
% self-consistent collapse or silently reset a reference within one
% physical history. The admissible mathematical initial-state domain
% is nonempty; this result restricts which preparation protocols can
% realize an exact unprojected positive-mass spherical history.

% Arbitrary nonlocal finite-count initial states, R-charge-changing
% operations, infinite duration/region limits and nonanalytic source
% tails remain outside this theorem. They retain the full coupled
% existence and uniform clock/time/ell question. Approximate radial
% histories in a controlled hydrodynamic limit are also not excluded.
% Constants depend on mu=ma, g=G/a^2 through any supplied actual fields,
% the fixed pulse/operations, fixed ell and finite duration. The
% exponential Dyson constants supply no uniform macroscopic time or
% allowed-ell bound. No angular projection, alternative coupling,
% field equation, numerical run or external X is introduced.
% SOURCES / INDEPENDENCE: exact native checkerboard commutation,
% trace-class Duhamel, finite-range weighted propagation, local
% trace-preserving CAR commutation and the already paired UA13/UA17
% analytic-source implication. All displayed pulse/count/source steps
% were derived before a new original-RED continuation was read.

% ACTUAL AC60--AC61 PULSE / NEUTRAL-CHANNEL COMPARISON [R].
% Read the complete saved RED proof and SOURCES after UA18. Its
% independent check was saved before actual UA18 intake and explicitly
% acknowledges the author's prior route announcement. The trace-class
% Duhamel and finite-count bounds agree. The operator-norm range-two
% Dyson series is a second proof of UA18.2's localized-leg estimate:
% every product of fewer than dist(A_x,B)/2 hopping factors has zero
% tested block, yielding exp[B_tot exp(2alpha)-alpha dist(A_x,B)].
% The fixed finite-time estimate needs no massless sea-tail bound.
% The nontrivial short lapse-pulse example uses the actual nonzero
% interband site-energy block already proved for all finite masses.

% RED gives a precise weaker neutral-channel condition, also accepted:
% sum_l K_l^dag Q_(R,B)K_l=Q_(R,B). The charge increment is the
% expectation of this LOCAL difference, hence zero even for a
% non-Gaussian current state; no subtraction of infinite charges is
% used. The sufficient Kraus commutation premise is now explicit above.
% Finite-region covariance increments are finite rank independently
% of the current state. A finite sum of their propagated contributions
% and the pulse Duhamel terms proves the stated source moments and
% finite expected counts. A nonzero-duration pulse does not imply
% hard finite excitation-sector support.

% Both proofs distinguish preparation of initial data relative to
% P0 from the binding common reference transport, which would leave
% an empty payload empty through the same pulse. The all-mass reached
% zero-source conclusion requires the displayed conserved R neutrality
% and exact spherical point/cell sources; generic trace preservation
% alone is insufficient. General nonlocal initial data, nonneutral
% channels, infinite protocols and approximate hydrodynamic histories
% remain separate. No new X, code, run or Gauss existence is inferred.

% SC1--SC3. CONTINUED-SPECTRAL LOCALITY: QUANTITATIVE REGULAR-CHART
% CONSTRUCTION [R, independently saved; actual paired check in SC6].
% This is a new binding-reference calculation, after the full revised
% GOAL and owner-qualified reviews were read. It uses the already
% paired pure-projector entropy inequality and flat kernel estimates
% as inputs; no new RED spectral-locality source was read first.
% The earlier literal-transport theorem is a comparator. The revised
% compact states the checked SC/SL continued-reference domains.

% SC1. State and operator domain first. Set a=1, epsilon=1/L and
% mu=ma. Consider an actual configuration path s in [0,1] of native
% self-adjoint operators
% h_s={n_s,h0}/2+sum_i{w_si,Q_i}/2,
% where n_s,w_si are the ACTUAL back-averaged site coefficients.
% They are not point samples substituted for W^dag(N,v). Assume their
% spatial Fourier measures have uniformly bounded first moment in
% scaled momentum: int |Q| (|nshat|(dQ)+sum_i|wsihat|(dQ))<=M1.
% Constant components cost zero. On a finite periodic parent this is
% the corresponding exact discrete Fourier sum, with its actual
% native representatives; no unproved alias cancellation is used.
% Suppose h_s is norm continuous and has an occupied spectral
% projection P_s continued from P_0, with P_s norm continuous.
% Require actual Riesz contours Gamma_s selecting these occupations
% and a uniform finite bound
% J_epsilon=sup_s (1/(2pi))int_(Gamma_s)
%                         ||(z-h_s)^(-1)||^2 |dz|.
% Their occupied sets need not be the lowest energies; occupations
% remain fixed through zero crossings. These contours and their
% bound are premises on the actual branch, not consequences of a
% pointwise symbol threshold. On an infinite parent this statement
% requires the corresponding isolated spectral sets. At a finite
% parent it applies to separated occupied/unoccupied clusters.

% The exact native Fourier commutator has a uniform first-moment
% bound at every finite mu, including massless nodes:
%             ||[h_s,P_0]|| <= C epsilon M1.            (SC1)
% For a lapse transfer q with endpoints k_+,k_-, write h_+=h0(k_+),
% E_+=|h_+| and likewise -. Its multiplier is exactly
% (E_+-E_-)[I-h_+h_-/(E_+E_-)]/4.
% Its norm is <=|E_+-E_-|/2<=C|q| off the nodes. At a zero endpoint
% use the original undivided expression and its specified filling;
% h at the other endpoint is O(|q|), giving the same bound. A flow
% transfer has multiplier
% [Q_i(k_+)+Q_i(k_-)] [P_0(k_-)-P_0(k_+)]/2.
% If both endpoint energies exceed a fixed multiple of |q|, the
% derivative bound ||grad P_0||<=C/E and |Q_i|<=E give O(|q|);
% otherwise both energies and the averaged Q_i are O(|q|).
% Thus it is O(|q|) also at the eight massless point nodes.
% Integrating these exact multipliers against the actual coefficient
% Fourier measures proves SC1. No mass-gap denominator is used there.

% SC2. The actual spectral branch remains close, when its contour
% bound is controlled. The exact resolvent commutator gives
% [P_s,P_0]=(1/(2pi i))int_(Gamma_s)
%                  (z-h_s)^(-1)[h_s,P_0](z-h_s)^(-1)dz.
% Hence ||[P_s,P_0]||<=kappa_epsilon,
% kappa_epsilon=C epsilon M1 J_epsilon. If kappa_epsilon<1/2,
% continuity from P_s=P_0 selects the SMALL projection-angle branch:
%              ||P_s-P_0||<=sqrt(2) kappa_epsilon=:d_epsilon.
% Indeed the singular principal angles have sin(theta)cos(theta)
% bounded by kappa. A path from zero angle cannot cross pi/4, at
% which that product is 1/2; on the selected branch sin(theta)<=
% sqrt(2)sin(theta)cos(theta). The same norm-projection argument
% applies on an infinite parent. A commutator estimate without this
% continuous occupation/path premise would also permit orthogonal
% commuting projectors and would NOT prove closeness.

% For a cut A with n_A native sites, pure-projector entropy stability
% then gives the explicit estimate
% |S_A(P_s)-S_A(P_0)|
%   <=C d_epsilon sqrt[n_A(A_lat+n_A d_epsilon^(3/2))].   (SC2)
% Here A_lat is its physical area/a^2 and the flat cross-block
% Schatten-3/2 sum must obey the already proved C A_lat bound for
% the stated regular cut/parent limit. Proof: X=chi_A P_s chi_Ac,
% Y=chi_A P_0 chi_Ac have norm <=1/2,
% ||X-Y||_2<=2sqrt(n_A)d_epsilon and
% sum singular(X)^(3/2)<=C[A_lat+n_A d_epsilon^(3/2)].
% Entropy is the sum of F(s)=h2((1-sqrt(1-4s^2))/2) over these
% singular values. The inequality |F'(s)|<=C s^(3/4), followed by
% the singular-value l2 bound and Cauchy--Schwarz, proves SC2.
% Thus if J_epsilon,M1 remain bounded and R/a<=C/epsilon,
% the normalized TRANSPORT-FREE SPECTRAL comparison error is
% O(sqrt(epsilon)); the flat planar/sphere limits determine the
% same leading local coefficient at each fixed mu. No rate for the
% separate flat-sphere convergence is silently asserted. Finite-parent
% sequences must also satisfy the actual flat-covariance/cut limit;
% a parent independent of R is not an infinite-area limit.

% SC3. What this supplies and what it does not. A fixed positive
% mass, sufficiently subcritical globally gapped branch, connected
% without crossings to the free sea, supplies a conventional contour
% with bounded J_epsilon. The earlier actual global-gap certificate
% is one sufficient test; a local uniform symbol gap alone is not.
% The first-gradient covariance and centered pure-cut entropy parity
% calculation above still distinguishes an O(gradient) state change
% from the vanishing first local entropy coefficient on that chart.
% SC2 is a controlled leading-coefficient estimate, not a sharp
% second-gradient expansion with a computed correction coefficient.
%
% Constant n,w have P_cont=P_0 exactly at every tilt when occupations
% are followed in the fixed momentum eigenspaces. The GROUND
% projector coincides under |w.q|<n E (with zero filling fixed).
% Along equal-component directions d=1,2,3 its exact massive
% threshold is |w|/n=sqrt(1+mu^2/d)+|mu|/sqrt(d).
% In particular an axial threshold cannot certify a radial field
% uniformly in direction. Actual back-averaged coefficients, not
% unaveraged geometric values, belong in these native tests.
%
% At mu=0 on the infinite parent, the initial sea is already gapless
% and these ordinary Riesz-isolation premises do not hold. The
% gap-independent commutator SC1 survives, but it does NOT supply
% J_epsilon or SC2 for an arbitrary massless continued projector.
% A low-energy/spectral-continuation estimate is still needed there.
% More generally a collapsing occupation separation can make
% epsilon J_epsilon large, even with smooth positive clock and slow
% spatial fields. This is an unresolved uniform estimate, not a
% counterexample to the binding model. A supplied Hartree chart from
% RR1 gives local continuation; it does not by itself give the
% uniform parent/time/ell spectral constants used in this theorem.
% Physical C need not equal this reference P_s; its actual excitations
% and cut must be compared separately. No literal propagated sea,
% angular-source obstruction, refilling or new dynamics is substituted.
% SOURCES / INDEPENDENCE: exact native endpoint commutators, Riesz
% calculus, continuous projection angles and the already paired
% area/TL Schatten entropy estimate; revised GOAL and owner reviews.
% This substantive save precedes new RED spectral-locality intake.
% No code/run, external X or actual collapse result is claimed.

% SC4--SC5. GAPLESS SUBCRITICAL FINITE-PARENT COMPLEMENT
% [R, independently saved before actual new RED SL-source intake].
% This is a different sufficient domain from SC2. It does not assume
% a uniform spectral gap, but DOES require that the actual continued
% reference still equals the ground projector: no refilling is used
% to impose that equality. A spatial symbol threshold alone does not
% establish the required occupation history.

% SC4. Take even antiperiodic cubic parents of side L_p comparable
% to epsilon^(-1), cuts of size R comparable to L_p, uniformly away
% from wrapping, and the actual scaled coefficient Fourier bound SC1.
% Assume n(x)>=n_*>0 and sup_x |w(x)|<=n_*-c, with c>0 fixed.
% This sufficient condition is stronger than local subcriticality.
% All mu in a fixed finite interval [0,mu_*] are allowed. Write
% E=|h0|, P0=(1-h0/E)/2 and
% h_d=P0 h P0+(1-P0)h(1-P0). Exact SC1 gives
%                  ||h-h_d||<=C epsilon.                  (SC4a)
% The actual finite-range commutators also give
% ||h psi|| >= c ||E psi||-C epsilon||psi||,
% because h=n h0+sum_i w_i Q_i+O(epsilon),
% ||n h0 psi||>=n_*||E psi|| and
% ||sum_i w_i Q_i psi||<=sup|w| (sum_i||Q_i psi||^2)^(1/2)
%                         <=sup|w| ||E psi||.
% The same coercivity estimate holds for h_d by SC4a.

% The block FORMS have the stronger signs
% (1-P0)h(1-P0)>= (1-P0)(cE-Cepsilon)(1-P0),
% P0 h P0<=P0(-cE+Cepsilon)P0.                            (SC4b)
% Here are the needed actual-symbol bounds, rather than an invoked
% semiclassical positivity assertion. On the positive block the lapse
% form is {n,E}/2. Its difference from E^(1/2)nE^(1/2) has Fourier
% multiplier (sqrt(E_+)-sqrt(E_-))^2/2, bounded by C|q|.
% On the negative block it has the opposite sign. Put t_i=Q_i/E.
% The difference of {w_i,Q_i}/2 and
% E^(1/2){w_i,t_i}E^(1/2)/2 has multiplier
% (sqrt(E_+)-sqrt(E_-))
% (Q_i(k_+)/sqrt(E_+)-Q_i(k_-)/sqrt(E_-))/2.
% Both factors are uniformly one-half Holder in k, also at mu=0;
% the second follows from Q_i=O(E), grad Q_i=O(1) and
% grad(Q_i/sqrt(E))=O(E^(-1/2)), splitting |q| above/below E.
% Thus this multiplier is O(|q|). Integrate against the actual
% Fourier coefficients to get the O(epsilon) form remainders.
% The norm of sum_i w_i t_i is <=sup|w| by the same row
% Cauchy--Schwarz estimate and sum_i t_i^2<=1. This proves SC4b.

% Let N_p be the number of parent sites. Elementary valley counting
% on this exact AP grid gives
% rank 1_(E<=r) <= C N_p(r+epsilon)^3.
% The coercivity estimate and the min--max principle therefore imply
% rank 1_(|h|<=delta)+rank 1_(|h_d|<=delta)
%                         <=C N_p(delta+epsilon)^3.        (SC4c)
% Indeed every vector in either low-energy subspace obeys
% ||E psi||<=(delta+Cepsilon)||psi||/c; apply min--max to E^2.
% SC4b likewise shows that the negative projector of h_d differs
% from P0 by rank at most C N_p epsilon^3: the wrong-sign subspaces
% have E-form at most Cepsilon, so min--max applies to E itself.
% Arbitrary fixed filling of actual zero eigenvalues only adds to
% these explicitly bounded low-energy ranks.

% Choose a smooth scalar f_delta equal to 1 on (-infinity,-delta]
% and 0 on [delta,infinity), between 0 and 1. Its scaled smooth
% derivative has integrable Fourier transform, so the Duhamel
% functional-calculus formula gives
% ||f_delta(h)-f_delta(h_d)||<=C||h-h_d||/delta.
% Decomposing each negative projector into f_delta plus its supported
% low-energy remainder proves
% P_-(h)-P0=B_delta+T_delta,
% ||B_delta||<=Cepsilon/delta, ||T_delta||<=3,
% rank T_delta<=C N_p(delta+epsilon)^3.                    (SC4d)
% This remains true with a specified partial zero-mode filling.
% It uses the actual negative spectral projector, and hence applies
% to the continued reference ONLY on the no-crossing domain stated
% above. It is not a definition of the reference inside pockets.

% SC5. Entropy consequence without a mass gap. Since N_p/n_A stays
% bounded in this stated parent/cut scaling, SC4d gives
% ||chi_A(P_-(h)-P0)chi_Ac||_2<=C sqrt(n_A)d,
% d=epsilon/delta+(delta+epsilon)^(3/2).
% Choose delta=epsilon^(2/5), so d=O(epsilon^(3/5)).
% For the pure-projector entropy function of cross-cut singular values,
% |F'(s)|<=C s^(5/6), 0<=s<=1/2. The earlier flat Schatten-3/2
% area bound also bounds its Schatten-5/3 sum. The same singular-value
% interpolation/Cauchy--Schwarz proof as SC2 consequently gives
% |S_A(P_-(h))-S_A(P0)|
%       <=C d sqrt[n_A(A_lat+n_A d^(5/3))].
% With n_A=O(epsilon^(-3)), A_lat comparable to epsilon^(-2),
% this is the normalized bound
%            |S_A(P_-(h))-S_A(P0)|/A_lat
%                         <=C epsilon^(1/10).             (SC5)
% Every constant retains the uniform actual coefficient, subcritical,
% parent/cut, mu and source-branch hypotheses; the flat finite-parent
% leading-coefficient limit is still a required separately established
% limit. This is a finite-parent low-energy-count proof, not an
% infinite-parent operator-norm comparison or a sharp gradient series.
% It proves a massless sufficient case where the minimum gap may tend
% to zero. It does not cover a continued occupied projector differing
% from the ground state, a supercritical trapping interior, an arbitrary
% occupation degeneracy choice, or missing uniform time/ell estimates.
% No solved collapse or actual Hartree existence is inferred.
%
% SOURCES / INDEPENDENCE: the exact native SC1 endpoint identity,
% explicit square-root form multipliers, AP-grid point-node counting,
% min--max and a written smooth functional-calculus decomposition;
% already paired pure-cut Schatten entropy bound. This substantive
% author route is saved before reading the newly saved RED SL1--SL4.
% No primary-source theorem, code or numerical execution is imported.

% SC6. ACTUAL SL1--SL6 COMPARISON AND MASSIVE REFINEMENT / SOURCES [R].
% Read the full actual independent RED SL1--SL4 only after SC1--SC5
% were substantively saved; then read its actual SC comparisons SL5--SL6.
% The following sharper massive construction originates in that saved
% independent RED route, and was checked here directly before use.
% At fixed mu>0, on a continuous path from h0 with uniform actual
% spectral gap (-delta,delta), bounded n,w, n>=n_*>0 and first/second
% spatial derivatives O(epsilon),O(epsilon^2), put
% B_ij(k)=-i h0(k)q_j(k)partial_i P0(k)/(2E(k)^2),
% a_ij(x)=partial_i w_j(x)/n(x),
% P1=(1/2)sum_ij{a_ij(X),B_ij(D)}.
% Since {h0,partial_i P0}=0, B_ij is Hermitian and
% [h0,B_ij]=-i q_j partial_i P0. The exact endpoint expansion gives
% [h,P0]=i sum_ij partial_i w_j q_j partial_i P0+O(epsilon^2);
% its first lapse-gradient term vanishes. Thus
% [h,P0+P1]=O(epsilon^2) and
% (P0+P1)^2-(P0+P1)=O(epsilon^2).
% All estimates hold in an exponentially weighted kernel/Schur norm:
% at mu>0 the native B_ij and P0 multipliers are analytic in a fixed
% strip, while Taylor remainders carry at most quadratic hopping offsets.
% The near-projector Riesz construction gives Q=P0+P1+O(epsilon^2).
% Block-diagonalizing h relative to Q changes it by O(epsilon^2).
% Its negative projector is Q by path continuity and the actual global
% gap. A resolvent identity on a common separated contour then proves
% P=P0+P1+R2, |(R2)_xy|<=Cepsilon^2 exp(-c|x-y|).
% The weighted resolvents are justified by conjugating with truncated
% exponential distance weights and a Neumann inverse at small fixed
% weight; the constants retain mu and the actual gap. A locally frozen
% symbol gap is not substituted for this pathwise operator gap.
%
% For cuts with #{x in A:x+z notin A}<=C A_lat(1+|z|), these kernels
% give trace-norm cross-cut R2=O(A_lat epsilon^2) and HS-square
% cross-cut (P-P0)=O(A_lat epsilon^2). Onsite Dirac time reversal fixes
% P0 and the cut, but sends B_ij to -B_ij, since it conjugates i and
% reverses both q_j and the momentum derivative. It sends P1 to -P1.
% Therefore the first trace variation of cut entropy at P0 vanishes
% for every such cut, without an additional centered-cut symmetry.
% This does not set the old literal-transport entropy rate to zero.
%
% For completeness, let Z be the Hermitian double of the cross block
% and F(x)=h2[(1-sqrt(1-4x^2))/2]. Exponential crossing counts give
% Tr|Z|^p<=C p^(-4)A_lat, 0<p<=1. Smooth F only at |x|<=d, with
% |F-F_d|<=C d^(2-p)|x|^p and ||F_d''||<=C(1+|log d|).
% The trace Hessian is bounded by ||F_d''|| ||delta Z||_HS^2,
% since its coefficients are divided differences of F_d'. The P1
% linear term vanishes by time reversal and the R2 term has the stated
% trace bound. Set d=epsilon,p=1/|log epsilon| to obtain
% |S_A(P)-S_A(P0)|/A_lat<=Cepsilon^2(1+|log epsilon|)^4.
% No entropy-eigenvalue gap or uniform mu->0 limit is concealed.
% This is stronger than SC2 on its more restrictive massive domain;
% SC4--SC5 supplies a different ground-sheet massless sufficient case.
%
% The actual SL6 check agrees with SC4--SC5, including its finite-rank
% interpolation: ||X-Y||_(5/3)^(5/3)<=C n_A d^(5/3) follows from
% rank <=4n_A and the HS bound, not from an operator-norm bound d.
% At mu=0 the homogeneous threshold value one still strictly orders
% each nonzero-energy momentum; it supplies no uniform gap at the nodes.
% The exact directional massive thresholds and actual l=0/l=1 adjoints
% agree. Neither route promotes a literal-transport control, unprojected
% spherical obstruction, refill rule or new dynamics to the binding model.
% Critical-first addendum read before this compact/comparison update:
% these saved proofs serve duty2d's printed local-area sentence only.
% No additional refinement, code/run, decimal or external X is asserted.
